\chapter{Mathematical Structures}
\section{Preliminary Remarks} \label{sec:preliminary-remarks-for-mathematical-structures}
In English, the term ``structure'' refers to the systematic organization of the components that constitute a complex entity.
Whenever one thinks of structures, what comes to mind is an ``order'' of some kind, or entities that have recognizable forms.
This intuitive notion is somewhat ubiquitous and can be found in almost every field of rational study:
\begin{itemize}
    \item In chemistry, the chemical structure of a water molecule is characterized by two hydrogen atoms attached to a single
          oxygen atom at a \( 104.45\degree \) angle.
    \item In computer science, a linked list is a data structure consisting of a collection of sequential nodes with each node
          containing data and a reference (link) to the next node in the sequence.
    \item In business and management, one speaks of an institute's leadership structure, which typically consists of a roster of
          executives along with a command chain that designates the level of authority each executive holds.
\end{itemize}
These examples all offer a glimpse of what a ``structure'' generally entails:
\begin{enumerate*}[label=(\alph*)]
    \item a collection of objects that are related to each other in a systematic and identifiable way;
    \item a set of rules/information describing the interrelationship between the objects in the collection.
\end{enumerate*}
Thus, at a conceptual level, a ``structure'' can be represented as follows:
\begin{equation*}
    \text{structure} = \text{set} + \text{descriptive information}
\end{equation*}
Simply put, a ``structure'' is a set endowed with additional features and properties that describe
the nature of the objects it contains. A collection of chess pieces is just a set; throw
a chessboard and a rulebook into the mix, and we have a structure. Instead of just knowing
what the pieces are, we also know how the pieces move and how they relate to each other.

\vspace{0.3cm}
\noindent
In mathematics, we seek to formalize the notion of a structure. That is, to define in abstract mathematical terms
what a structure is. With set theory serving as a foundation of mathematics, we can narrow the goal
even further: to define mathematical structures in terms of sets. As we shall soon see,
there are generally two classes of descriptive information that defines a mathematical structure:
\keyword{relations} and \keyword{operations}. These, too, have to be precisely defined as sets.
But before that can be done, some groundwork has to be laid.

\section{The Set of Natural Numbers}
Intuitively, we all know what the natural numbers are: a system of symbols representing the number of
objects in a finite collection of discrete elements. These symbols are abstract in nature, stripped of
any reference to the concrete qualities of the objects being counted. Be it two dogs, two oranges or two books, the
same symbol `2' is used to represent their quantities.
The power of these symbols lies not only in their expressiveness but also in the fact that they facilitate higher
computation, especially with the advent of arithmetic, which transformed the act of counting from a physical process (counting with fingers, etc.) to that of pure
symbolic manipulation.
Nevertheless, despite the highly intuitive and self-evident nature of the natural numbers,
for them to be truly rooted in the foundation of mathematics, they have to be rigorously defined in terms of
simpler and more fundamental mathematical concepts---namely, sets. Thus, the goal of this section is twofold:
\begin{enumerate*} [label=(\roman*)]
    \item to present a systematic method that assigns a unique set to each symbol \( 0,1,2,3,\dots \) of the natural numbers,
    \item to prove that there exists a set containing exactly those sets.
\end{enumerate*}

\subsection{Von Neumann Ordinals} \label{ssec:von-neumann-ordinals}
Assigning a unique set to each symbol \( 0,1,2,3,\dots \) is a fairly straightforward task. Since each natural number represents
a finite, discrete quantity, we simply have to assign to each natural number \( n \) a single representative of all sets containing
exactly \( n \) elements. Thus, the symbol 0 can be defined as a set with no elements, the symbol 1 a set with
exactly one element, the symbol 2 a set with exactly two elements, etc.

\vspace{0.3cm}
\noindent
For the symbol 0, the process of finding a representative set is trivial courtesy of 
\cref{thm:uniqueness-empty-set}, which asserts the uniqueness of the empty set.
\begin{equation*}
    0 \coloneq \emptyset
\end{equation*}
But what about the rest of the natural numbers? Surely there are more than one way to form sets containing a certain number of elements?
For example, the sets \( \mleft\{ 0 \mright\} \), \( \mleft\{ \mleft\{ 0 \mright\} \mright\} \), \( \mleft\{ \mleft\{ \mleft\{ 0 \mright\} \mright\} \mright\} \),
and \( \mleft\{ \orderedpair{0}{0} \mright\} \) all have exactly one element.
So which one of them should we pick to represent the symbol 1? There is no right answer, of course.
In fact, the exact definition doesn't even matter as long as it yields all the familiar properties of natural numbers.
Here, we will use the most straightforward and intuitively obvious approach:
\begin{align*}
    0 &\coloneq \emptyset = \mleft\{ \mright\} \\
    1 &\coloneq 0 \union \mleft\{ 0 \mright\} = \{0\} \\
    2 &\coloneq 1 \union \mleft\{ 1 \mright\} = \{0\} \union \{1\} = \{0,1\}\\
    3 &\coloneq 2 \union \mleft\{ 2 \mright\} = \{0,1\} \union \{2\} = \{0,1,2\}\\
    \vdots
\end{align*}
The idea is to define any natural number \( n \) as a set containing all the natural numbers that are smaller
than it, thereby ensuring that the set assigned to \( n \) always contains exactly \( n \) elements.
To construct the representative set for the ``next'' natural number,
we simply adjoin one more element (\( n \) itself) to \( n \).
Any set constructed by means of this recursive process is called a \keyword{von Neumann ordinal}.

\subsection{Set-Theoretic Definition of \naturalset}
In the previous section, we presented a scheme for assigning a unique set to the familiar symbols \( 0,1,2,3,\dots \) denoting
the natural numbers. The next task is to prove (on the basis of the axioms of set theory) that there exists a set containing
exactly the von Neumann ordinals.

\begin{dfn}[Successor]
    \label{dfn:successor}
    For any set \( x \), the \keyword{successor} of \( x \) is the set \( x \union \mleft\{ x \mright\} \).
\end{dfn}

\begin{dfn}[Inductive Set]
    \label{dfn:inductive-set}
    A set \( I \) is \keyword{inductive} if and only if the following conditions hold:
    \begin{enumerate}[topsep=2pt,noitemsep,label=\texttt{\textup{(\alph*)}}]
        \item \( \emptyset \in I \)
        \item \( \forall x \, \big[ x \in I \implies x \union \mleft\{ x \mright\} \in I \big] \)
    \end{enumerate}
\end{dfn}

\begin{rem}
    Essentially, an inductive set is any set that contains all the von Neumann ordinals.
\end{rem}

\begin{cor}
    \label{thm:any-von-neumann-ordinal-belongs-to-an-inductive-set}
    If \( I \) is inductive and \( x \) is a von Neumann ordinal, then \( x \in I \).
\end{cor}

\begin{dfn}
    \label{dfn:the-set-of-natural-numbers}
    \( \naturalset \) is the set given by,
    \begin{equation*}
        \naturalset = \setbuilderbigg{x}{\forall I \, \Big[ \emptyset \in I \land \forall z \, \big( z \in I \implies z \union \mleft\{ z \mright\} \in I \big) \implies x \in I \Big]}
    \end{equation*}
\end{dfn}

\begin{thm}
    \label{thm:existence-of-the-set-of-all-natural-numbers}
    \( \naturalset \) exists.
\end{thm}

\begin{proof}
    Let the predicate \( \varphi \) be defined as,
    \begin{equation}
        \label{eq:phi-is-the-predicate-for-the-set-of-N}
        \varphi(x) \coloneq \forall I \, \bigg[ \emptyset \in I \land
        \forall z \, \Big( z \in I \implies z \union \mleft\{ z \mright\} \in I \Big)
        \implies x \in I \bigg]
    \end{equation}
    Let \( I_{0} \) be an inductive set (note that \( I_{0} \) exists by the \nameref{axm:axiom-of-infinity}).
    By \cref{dfn:inductive-set}, we have,
    \begin{equation}
        \label{eq:I-is-inductive}
        \big( \emptyset \in I_{0} \big) \land \forall z \, \Big( z \in I_{0} \implies z \union \mleft\{ z \mright\} \Big)
    \end{equation}
    Let \( x_{0} \) be an arbitrary object such that \( \varphi(x_{0}) \) is satisfied. By \eqref{eq:phi-is-the-predicate-for-the-set-of-N},
    it follows that,
    \begin{equation}
        \label{eq:x0-satisfies-the-predicate-for-N}
        \forall I \, \bigg[ \emptyset \in I \land \forall z \, \Big( z \in I \implies z \union \mleft\{ z \mright\} \in I \Big) \implies x_{0} \in I \bigg]
    \end{equation}
    Applying \aref{rule-of-inference:universal-instantiation}{Universal Instantiation} on \eqref{eq:x0-satisfies-the-predicate-for-N} yields,
    \begin{equation*}
        \emptyset \in I_{0} \land \forall z \, \Big( z \in I_{0} \implies z \union \mleft\{ z \mright\} \in I_{0} \Big) \implies  x_{0} \in I_{0}
    \end{equation*}
    Using \eqref{eq:I-is-inductive}, we conclude by \aref{rule-of-inference:modus-ponens}{Modus Ponens} that,
    \begin{equation}
        \label{eq:x0-in-I0}
        x_{0} \in  I_{0}
    \end{equation}
    Since \eqref{eq:x0-in-I0} was derived on the assumption that \( \varphi(x_{0}) \) is satisfied, we conclude by the \aref{rule-of-inference:deduction-theorem}{Deduction Theorem} that,
    \begin{equation*}
        \varphi(x_{0}) \implies x_{0} \in I_{0}
    \end{equation*}
    Since \( x_{0} \) is arbitrary, we have by \aref{rule-of-inference:universal-generalization}{Universal Generalization},
    \begin{equation*}
        \forall x \, \big( \varphi(x) \implies x \in I_{0} \big)
    \end{equation*}
    In other words,
    \begin{equation*}
        \exists A \, \forall x \, \Big[ \varphi(x) \implies x \in A \Big]
    \end{equation*}
    Thus, the unique set \( \naturalset = \setbuilderbig{x}{\varphi(x)} \) exists by \cref{thm:sufficient-condition-for-unrestricted-comprehension}.
\end{proof}

\begin{thm}
    \label{thm:the-set-of-all-natural-numbers-is-inductive}
    \( \naturalset \) is inductive.
\end{thm}

\begin{proof}
    Let \( I_{0} \) be an arbitrary set. It is obvious from the laws of propositional logic that the statement,
    \begin{equation*}
        \emptyset \in I_{0} \land \forall z \, \Big( z \in I_{0} \implies z \union \mleft\{ z \mright\} \in I_{0} \Big) \implies \emptyset \in I_{0}
    \end{equation*}
    is trivially true. Since \( I_{0} \) is arbitrary, we conclude by \aref{rule-of-inference:universal-generalization}{Universal Generalization} that,
    \begin{equation*}
        \forall I \, \bigg[ \emptyset \in I \land \forall z \, \Big( z \in I \implies z \union \mleft\{ z \mright\} \in I \Big)
        \implies \emptyset \in I \bigg]
    \end{equation*}
    Thus, \( \emptyset \in \naturalset \) by \cref{dfn:the-set-of-natural-numbers}.

    \vspace{0.3cm}
    \noindent
    Now let \( x_{0} \in \naturalset \) be arbitrary. Thus, by \cref{dfn:the-set-of-natural-numbers},
    \begin{equation}
        \label{eq:x0-belongs-to-the-set-of-N}
        \forall I \, \bigg[ \emptyset \in I \land \forall z \, \Big( z \in I \implies z \union \mleft\{ z \mright\} \in I \Big) \implies x_{0} \in I \bigg]
    \end{equation}
    Let \( I_{0} \) be any inductive set. By \eqref{eq:x0-belongs-to-the-set-of-N}, it follows that \( x_{0} \in I_{0} \). Since \( I_{0} \) is inductive and
    \( x_{0} \in I_{0} \), necessarily \( x_{0} \union \mleft\{ x_{0} \mright\} \in I_{0} \) by \cref{dfn:inductive-set}. Since \( I_{0} \) is an arbitrary inductive
    set, we conclude by \aref{rule-of-inference:universal-generalization}{Universal Generalization} that,
    \begin{equation*}
        \forall I \, \bigg[ \emptyset \in I \land
        \forall z \, \Big( z \in I \implies z \union \mleft\{ z \mright\} \in I \Big) \implies
        x_{0} \union \mleft\{ x_{0} \mright\} \in  I \bigg]\\
    \end{equation*}
    And thus, by \cref{dfn:the-set-of-natural-numbers},
    \begin{equation*}
        x_{0} \union \mleft\{ x_{0} \mright\} \in \naturalset
    \end{equation*}
    In other words,
    \begin{equation*}
        \begin{aligned}
            x_{0} \in \naturalset \implies x_{0} \union \mleft\{ x_{0} \mright\} \in \naturalset&
        \end{aligned}
        \begin{aligned}
            &\text{for any \( x_{0} \)}
        \end{aligned}
    \end{equation*}
    Thus,
    \begin{equation*}
        \emptyset \in \naturalset \land \forall x \, \Big[ x \in \naturalset \implies x \union \mleft\{ x \mright\} \in \naturalset \Big]
    \end{equation*}
    By \cref{dfn:inductive-set}, \( \naturalset \) is necessarily inductive.
\end{proof}

\begin{thm}
    \label{thm:the-set-of-natural-numbers-is-a-subset-of-any-inductive-set}
    If \( I \) is inductive, then \( \naturalset \subseteq I \).
\end{thm}

\begin{proof}
    Let \( x_{0} \) be an arbitrary element of \naturalset.
    By \cref{dfn:the-set-of-natural-numbers}, it follows that,
    \begin{equation}
        \label{eq:x0-belongs-to-N}
        \forall I \, \bigg[ \emptyset \in I \land \forall z \, \Big( z \in I \implies z \union \mleft\{ z \mright\} \in I \Big) \implies x_{0} \in I \bigg]
    \end{equation}
    Let \( I_{0} \) be any inductive set. Thus, by \cref{dfn:inductive-set},
    \begin{equation}
        \label{eq:I0-is-inductive}
        \emptyset \in I_{0} \land \forall z \, \Big( z \in I_{0} \implies z \union \mleft\{ z \mright\} \in I_{0} \Big)
    \end{equation}
    Applying \aref{rule-of-inference:universal-instantiation}{Universal Instantiation} to \eqref{eq:x0-belongs-to-N}, we have,
    \begin{equation}
        \label{eq:instantiated-statement-for-I0-is-inductive}
        \emptyset \in I_{0} \land \forall z \, \Big( z \in I_{0} \in z \union \mleft\{ z \mright\} \in I_{0} \Big) \implies x_{0} \in I_{0}
    \end{equation}
    By \eqref{eq:I0-is-inductive} and \eqref{eq:instantiated-statement-for-I0-is-inductive}, we conclude by \aref{rule-of-inference:modus-ponens}{Modus Ponens} that \( x_{0} \in \naturalset \).
    Since \( x_{0} \) is arbitrary, we can generalize,
    \begin{equation*}
        \forall x \, \Big( x \in \naturalset \implies x \in I_{0} \Big)
    \end{equation*}
    By \cref{dfn:subset}, it follows that \( \naturalset \subseteq I_{0} \) for any inductive set \( I_{0} \).
\end{proof}

\begin{thm}
    \label{thm:the-set-of-all-natural-numbers-contains-exactly-the-von-Neumann-ordinals}
    For any set \( x \), \( x \) is a von Neumann ordinal if and only if \( x \in \naturalset \).
\end{thm}

\begin{proof}
    Let \( x_{0} \) be a von Neumann ordinal. Since \( \naturalset \) is inductive by \cref{thm:the-set-of-all-natural-numbers-is-inductive},
    it follows from \cref{thm:any-von-neumann-ordinal-belongs-to-an-inductive-set} that \( x_{0} \in \naturalset \).

    \vspace{0.3cm}
    \noindent
    To prove the converse, let \( I_{0} \) be any inductive set. From \cref{thm:any-von-neumann-ordinal-belongs-to-an-inductive-set},
    we know that \( I_{0} \) contains all von Neumann ordinals. Therefore, by the \nameref{axm:axiom-schema-of-specification}, there must exist a non-empty set
    \( S \subseteq I_{0} \) such that every element of \( S \) is a von Neumann ordinal.
    \begin{equation}
        \label{eq:S-is-a-subset-of-I0-that-contains-all-von-Neumann-ordinals}
        S = \setbuilder{x \in I_{0}}{\text{\( x \) is a von Neumann ordinal}}
    \end{equation}
    The next task is to prove that \( S \) is inductive.

    \vspace{0.3cm}
    \noindent
    Since \( I_{0} \) is an inductive set, necessarily \( \emptyset \in I_{0} \) by \cref{dfn:inductive-set}. Since \( \emptyset  \)
    is a von Neumann ordinal and \( \emptyset \in I_{0} \), necessarily \( \emptyset \in S \)
    by \eqref{eq:S-is-a-subset-of-I0-that-contains-all-von-Neumann-ordinals}.
    Now let \( x_{0} \in S \) be arbitrary. By \eqref{eq:S-is-a-subset-of-I0-that-contains-all-von-Neumann-ordinals},
    \( x_{0} \) must be a von Neumann ordinal. Since \( S \subseteq I_{0} \) and
    \( x_{0} \in S \), necessarily \( x_{0} \in I_{0} \). But \( I_{0} \) is inductive, which means that \( x_{0} \union \mleft\{ x_{0} \mright\} \) must
    also be an element of \( I_{0} \). And since \( x_{0} \) is a von Neumann ordinal, \( x_{0} \union \mleft\{ x_{0} \mright\} \) must
    be a von Neumann ordinal as well (see \cref{ssec:von-neumann-ordinals}). Thus, \( x_{0} \union \mleft\{ x_{0} \mright\} \in I_{0} \) and
    \( x_{0} \union \mleft\{ x_{0} \mright\} \) is a von Neuman ordinal. By \eqref{eq:S-is-a-subset-of-I0-that-contains-all-von-Neumann-ordinals},
    it follows that \( x_{0} \union \mleft\{ x_{0} \mright\} \in S \). Thus,
    \begin{equation*}
        \begin{aligned}
            x_{0} \in S \implies x_{0} \union \mleft\{ x_{0} \mright\} \in S&
        \end{aligned}
        \begin{aligned}
            &\text{for any \( x_{0} \)}
        \end{aligned}
    \end{equation*}
    By \aref{rule-of-inference:universal-generalization}{Universal Generalization}, we have,
    \begin{equation*}
        \emptyset \in S \land \forall x \, \Big[ x \in S \implies x \union \mleft\{ x \mright\} \in S \Big]
    \end{equation*}
    Thus, \( S \) must be an inductive set by \cref{dfn:inductive-set}.
    And by \cref{thm:the-set-of-natural-numbers-is-a-subset-of-any-inductive-set}, \( \naturalset \subseteq S \). Therefore,
    every element of \( \naturalset \) is necessarily a von Neumann ordinal.
\end{proof}

\begin{rem}
    It should be evident from \cref{thm:the-set-of-all-natural-numbers-contains-exactly-the-von-Neumann-ordinals} and
    the \nameref{axm:axiom-of-extensionality} that \( \naturalset \) contains all the von Neumann ordinals and
    nothing else. Since each symbol for the natural numbers is assigned a unique von Neumann ordinal,
    it is justified to think of \( \naturalset \) as the set of all natural numbers.
\end{rem}

\section{Ordered N-Tuples}
Having rigorously defined the natural numbers as sets and the set \( \naturalset \) which contains all the natural numbers,
we now turn to an immediate application of the natural numbers.

\begin{dfn}
    \label{dfn:n-tuples}
    For any \( n \in \naturalset \), an \keyword{n-tuple} is any function with \( n \) as its domain.
    An n-tuple is also called a \keyword{tuple} or \keyword{finite sequence} of length \( n \).
\end{dfn}

\begin{dfn}[Sequence Notation]
    \label{dfn:n-tuples-sequence-notation}
    For any \( n \in \naturalset \) where \( n \neq 0 \),
    \( t = \tuple{x_{0}, \dots, x_{n-1}} \) if and only if the following conditions hold:
    \begin{enumerate}[topsep=2pt,noitemsep,label=\texttt{\textup{(\alph*)}}]
        \item \( \surjection{t}{n}{\mleft\{ x_{0},\dots,x_{n-1} \mright\}} \)
        \item \( \forall i\,{\in}\,n \, \big( t(i) = x_{i} \big) \)
    \end{enumerate}
\end{dfn}

\begin{rem}
    An \( n \)-tuple is really just an indexed system of sets with a natural number \( n \) acting as
    the index set \( I \).
\end{rem}

\begin{cor}[0-Tuple]
    \label{thm:zero-tuple}
    The \keyword{0-tuple} (also known as the \keyword{empty sequence}) is simply the empty set \( \emptyset \)
    since it is an indexed system of sets \( t \)
    with \( \dom t = \ran t = 0 = \emptyset \).
\end{cor}

\begin{notation}
    We denote the \( 0 \)-tuple or empty sequence as \( \tuple{} \).
\end{notation}

\begin{dfn}
    \phantomsection \label{dfn:set-of-all-finite-sequences-of-elements-from-a-set}
    For any set \( A \), \( \finiteseq{A} \coloneq \Union_{n \,\in\, \naturalset} A^{n} \).
\end{dfn}

\begin{rem}
    \( \finiteseq{A} \) is the set of all finite sequences of elements of \( A \).
\end{rem}

\begin{thm}
    \phantomsection \label{thm:seq-empty-set-is-the-singleton-containing-the-empty-set}
    \( \finiteseq{\emptyset} = \mleft\{ \emptyset \mright\} \).
\end{thm}

\begin{proof}
    By \cref{dfn:set-of-all-finite-sequences-of-elements-from-a-set}, we have \( \finiteseq{\emptyset} = \Union_{n \,\in\, \naturalset} \emptyset^{n} \).
    By \cref{dfn:union-of-indexed-system-of-sets}, it follows that,
    \begin{equation*}
        \finiteseq{\emptyset} = \Union_{n \,\in\, \naturalset} \emptyset^{n} = \setbuilderBig{x}{\exists n\,{\in}\,\naturalset \, \big[ x \in \emptyset^{n} \big]}
    \end{equation*}
    We know that \( \emptyset \in \finiteseq{\emptyset} \)
    because \( \emptyset^{0} = \emptyset^{\emptyset} = \mleft\{ \emptyset \mright\} \) by \cref{thm:the-exponentiation-of-the-empty-set}
    and thus \( \emptyset \in \emptyset^{0} \).
    Now let \( x_{0} \in \finiteseq{\emptyset} \) be arbitrary.
    By the definition of \( \finiteseq{\emptyset} \),
    this means that \( x_{0} \in \emptyset^{n} \) for some \( n \in \naturalset \). But by \cref{thm:the-exponentiation-of-the-empty-set},
    \( \emptyset^{n} \neq \emptyset \) if and only if \( n = 0 \). Thus, \( x_{0} \in \emptyset^{\emptyset} = \mleft\{ \emptyset \mright\} \),
    so necessarily \( x_{0} = \emptyset \). In other words, \( x_{0} \in \finiteseq{\emptyset} {\iff} x_{0} = \emptyset \) for any arbitrary \( x_{0} \).
    By the \nameref{axm:axiom-of-pairing}, it follows that \( \finiteseq{\emptyset} \) is the \aref{dfn:unordered-pair}{singleton} \( \mleft\{ \emptyset \mright\} \).
\end{proof}

\begin{thm}
    \label{thm:every-non-zero-tuple-is-expressible-in-sequence-notation}
    If \( t \) is a tuple with length \( n \neq 0 \), then,
    \begin{equation*}
        \exists x_{0} \, \dots \exists x_{n-1} \, \Big[ t = \tuple{x_{0}, \dots, x_{n-1}} \Big]
    \end{equation*}
\end{thm}

\begin{proof}
    Since \( t \) is a function with \( \dom t = n \) where \( n \in \naturalset \) and \( n \neq \emptyset \),
    it follows from \cref{dfn:domain-of-a-set-of-ordered-pairs} that there exists some \( x_{i} = t(i) \) for
    every \( i \in n \). Let the set \( R \) be defined as,
    \begin{equation}
        \label{eq:list-out-all-elements-of-ran-t}
        R = \mleft\{ x_{0},\dots,x_{n-1} \mright\} = \setbuilderbigg{x}{\Lor_{i \,\in\, n} x = x_{i}}
    \end{equation}
    We can easily prove that \( R = \ran t \).

    \vspace{0.2cm}
    \noindent
    Let \( \alpha \in R \) be arbitrary. Clearly, \( \alpha = x_{i} = t(i) \) for some \( i \in n \); in other words,
    \( \orderedpair{i}{\alpha} \in t \) for some \( i \in n \). By \cref{dfn:range-of-a-set-of-ordered-pairs}, it
    follows that \( \alpha \in \ran t \). Conversely, let \( \alpha \in \ran t \) be arbitrary.
    Again, by \cref{dfn:range-of-a-set-of-ordered-pairs}, there must exist some \( i \in n \) such that
    \( \orderedpair{i}{\alpha} \in t \). Thus, \( \alpha = t(i) \) by \cref{dfn:functional-notation}.
    Since \( t(i) = x_{i} \) for all \( i \in n \), it follows that \( \alpha = x_{i} \).
    By \eqref{eq:list-out-all-elements-of-ran-t}, clearly \( \alpha \in R \).
    Thus, \( \ran t = R = \mleft\{ x_{0},\dots,x_{n-1} \mright\} \) so that,
    \begin{equation*}
        \surjection{t}{n}{\ran t} \iff \surjection{t}{n}{\mleft\{ x_{0},\dots,x_{n-1} \mright\}}
    \end{equation*}
    where \( t(i) = x_{i} \) for all \( i \in n \). Thus, \( t = \tuple{x_{0}, \dots, x_{n-1}} \) by \cref{dfn:n-tuples-sequence-notation}.
\end{proof}

\begin{thm}
    \label{thm:if-every-coordinate-of-an-n-tuple-belongs-to-A-then-the-n-tuple-belongs-to-the-power-set-of-the-cartesian-product-between-n-and-A}
    For any \( n \in \naturalset \) where \( n \neq 0 \),
    if \( x_{0},\dots,x_{n-1} \in A \) for some set \( A \), then \( \tuple{x_{0},\dots,x_{n-1}} \in \powerset{n \times A} \).
    \begin{equation*}
        \forall n\,{\in}\,\naturalset \,
        \bigg[ n \neq 0 \implies
        \forall A \, \forall x_{0} \, \dots \forall x_{n-1} \, \Big( \Land_{i \,\in\, n} x_{i} \in A \implies \tuple{x_{0}, \dots, x_{n-1}} \in \powerset{n \times A} \Big)
        \bigg]
    \end{equation*}
\end{thm}

\begin{proof}
    Let \( n \) be any non-zero natural number and \( A \) be any set.
    Let the objects \( x_{0},\dots,x_{n-1}\) be \( n \) arbitrary elements of \( A \).
    Then by \cref{dfn:n-tuples-sequence-notation},
    \begin{equation*}
        \surjection{ \tuple{x_{0},\dots,x_{n-1}} }{n}{\mleft\{ x_{0},\dots,x_{n-1} \mright\}}
    \end{equation*}
    Clearly, \( \ran(\tuple{x_{0},\dots,x_{n-1}}) = \mleft\{ x_{0},\dots,x_{n-1} \mright\} \subseteq A  \)
    and \( \dom \mleft( \tuple{x_{0}, \dots, x_{n-1}} \mright) = n \). Thus, by \cref{dfn:arrow-notation},
    \begin{equation*}
        \function{\tuple{x_{0},\dots,x_{n-1}}}{n}{A}
    \end{equation*}
    Since \( \tuple{x_{0},\dots,x_{n-1}} \) is a function on \( n \) into \( A \),
    it follows from \cref{thm:a-function-on-A-into-B-is-a-subset-of-the-cartesian-product-between-A-and-B}
    that,
    \begin{equation*}
        \tuple{x_{0},\dots,x_{n-1}} \subseteq n \times A
    \end{equation*}
    By \cref{dfn:power-set}, it follows that \( \tuple{x_{0},\dots,x_{n-1}} \in \powerset{n \times A} \).
    By the \aref{rule-of-inference:deduction-theorem}{Deduction Theorem}, we have,
    \begin{equation*}
        \Land_{i \,\in\, n} x_{i} \in A \implies \tuple{x_{0}, \dots, x_{n-1}} \in \powerset{n \times A}
    \end{equation*}
    Since \( x_{0},\dots,x_{n-1} \) are arbitrary, \aref{rule-of-inference:universal-generalization}{Universal Generalization} yields,
    \begin{equation*}
        \forall x_{0} \, \dots \forall x_{n-1} \, \bigg( \Land_{i \,\in\, n} x_{i} \in A \implies \tuple{x_{0}, \dots, x_{n-1}} \in \powerset{n \times A} \bigg)
    \end{equation*}
    Since \( A \) is also arbitrary, we can generalize once again,
    \begin{equation*}
        \forall A \, \forall x_{0} \, \dots \forall x_{n-1} \, \bigg( \Land_{i \,\in\, n} x_{i} \in A \implies \tuple{x_{0}, \dots, x_{n-1}} \in \powerset{n \times A} \bigg)
        \qedhere
    \end{equation*}
\end{proof}

\begin{thm}[Equality]
    \label{thm:tuple-equality}
    For any \( n \in \naturalset \) where \( n \neq 0 \),
    \begin{equation*}
        \tuple{x_{0},\dots,x_{n-1}} = \tuple{y_{0},\dots,y_{n-1}} \iff \Land_{i \,\in\, n} x_{i} = y_{i}
    \end{equation*}
\end{thm}

\begin{proof}
    Let \( t_{x} \) and \( t_{y} \) be two arbitrary n-tuples with \( n \neq 0 \).
    In other words,
    \begin{gather*}
        t_{x} = \tuple{x_{0},\dots,x_{n-1}}\\
        t_{y} = \tuple{y_{0},\dots,y_{n-1}}
    \end{gather*}
    for arbitrary objects \( x_{0},\dots,x_{n-1} \) and \( y_{0},\dots,y_{n-1} \).
    Thus, by \cref{dfn:n-tuples-sequence-notation},
    \begin{gather*}
        \surjection{t_{x}}{n}{\mleft\{ x_{0},\dots,x_{n-1} \mright\}} \\
        \surjection{t_{y}}{n}{\mleft\{ y_{0},\dots,y_{n-1}\mright\}}
    \end{gather*}
    The problem is thus reduced to that of the equality of functions. And by \cref{thm:equality-of-functions},
    \begin{equation*}
        t_{x} = t_{y} \iff \dom t_{x} = \dom t_{y} \land \forall i\,{\in}\,n \, \big[ t_{x}(i) = t_{y}(i) \big]
    \end{equation*}
    Since \( \dom t_{x} = \dom t_{y} = n \), the formula reduces to,
    \begin{equation}
        \label{eq:condition-for-tuple-equality-from-first-principle}
        t_{x} = t_{y} \iff \forall i\,{\in}\,n \, \big[ t_{x}(i) = t_{y}(i) \big]
    \end{equation}
    But by \cref{dfn:n-tuples},
    \begin{equation*}
        \forall i\,{\in}\,n \, \big[ x_{i} = t_{x}(i) \big] \land
        \forall i\,{\in}\,n \, \big[ y_{i} = t_{y}(i) \big]
    \end{equation*}
    Thus, \eqref{eq:condition-for-tuple-equality-from-first-principle} is equivalent to,
    \begin{equation*}
        t_{x} = t_{y} \iff \forall i\,{\in}\,n \, \big[ x_{i} = y_{i} \big]
        \iff \Land_{i \,{\in}\, n} x_{i} = y_{i} \qedhere
    \end{equation*}
\end{proof}

\begin{rem}
    Obviously, two tuples can only be equal if they are of the same length.
    After all, two functions cannot possibly be equal unless they have
    equal domains.
\end{rem}

\begin{thm}
    \label{thm:sufficient-condition-for-unrestricted-specification-for-tuples}
    If \( \varphi \) is an n-ary predicate with \( n \in \naturalset \) and \( n \neq 0 \), then,
    \begin{multline*}
        \exists A_{0} \,\dots \exists A_{n-1} \, \forall x_{0} \,\dots \forall x_{n-1} \, \bigg[ \varphi(x_{0},\dots,x_{n-1})
        \implies \Land_{i \,{\in}\, n} x_{i} \in A_{i} \bigg] \implies \\
        \exists! S \, \forall x \, \bigg[ x \in S \iff \exists a_{0} \, \dots \exists a_{n-1} \, \Big( x = \tuple{a_{0},\dots,a_{n-1}} \land \varphi(a_{0},\dots,a_{n-1}) \Big) \bigg]
    \end{multline*}
\end{thm}

\begin{proof}
    Let \( \varphi \) be an n-ary predicate with \( n \neq 0 \) such that,
    \begin{equation}
        \label{eq:n-ary-phi-implies-each-coordinate-belongs-to-an-existing-set}
        \forall x_{0} \, \dots \forall x_{n-1} \, \bigg[ \varphi(x_{0},\dots,x_{n-1}) \implies \Land_{i \,\in\, n} x_{i} \in A_{i} \bigg]
    \end{equation}
    for existing sets \( A_{0},\dots,A_{n-1} \).
    Let the unary predicate \( \Gamma \) be defined as follows:
    \begin{equation}
        \label{eq:express-n-ary-phi-as-unary-gamma}
        \Gamma(x) \coloneq \exists a_{0} \, \dots \exists a_{n-1} \, \Big[ x = \tuple{a_{0},\dots,a_{n-1}} \land \varphi(a_{0},\dots,a_{n-1}) \Big]
    \end{equation}
    Let \( \alpha \) be an arbitrary object such that \( \Gamma(\alpha) \) is satisfied.
    Therefore, it follows from \eqref{eq:express-n-ary-phi-as-unary-gamma} that,
    \begin{equation*}
        \begin{aligned}
            \alpha = \tuple{a_{0}, \dots, a_{n-1}} \land \varphi(a_{0}, \dots, a_{n-1})&
        \end{aligned}
        \begin{aligned}
            &\text{for some \( a_{0}, \dots, a_{n-1} \)}
        \end{aligned}
    \end{equation*}
    Since \( \varphi(a_{0}, \dots, a_{n-1}) \) is satisfied, we infer from
    \eqref{eq:n-ary-phi-implies-each-coordinate-belongs-to-an-existing-set} that,
    \begin{equation}
        \label{eq:ai-belongs-to-Ai}
        \Land_{i \,\in\, n} a_{i} \in A_{i}
    \end{equation}
    An immediate consequence of \eqref{eq:ai-belongs-to-Ai} is that,
    \begin{equation*}
        \Land_{i \,\in\, n} \bigg[a_{i} \in \Union_{i \,\in\, n} A_{i}\bigg]
    \end{equation*}
    Thus, by \cref{thm:if-every-coordinate-of-an-n-tuple-belongs-to-A-then-the-n-tuple-belongs-to-the-power-set-of-the-cartesian-product-between-n-and-A},
    \begin{equation*}
        \alpha = \tuple{a_{0}, \dots, a_{n-1}} \in \powerset{n \times \Union_{i \,\in\, n} A_{i}}
    \end{equation*}
    By the \aref{rule-of-inference:deduction-theorem}{Deduction Theorem}, \( \Gamma(\alpha) \implies \alpha \in \powerset{n \times \Union_{i \,\in\, n} A_{i}} \). Since \( \alpha \) is arbitrary,
    we can generalize the formula to obtain,
    \begin{equation*}
        \forall x \, \Bigg[ \Gamma(x) \implies x \in \powerset{n \times \Union_{i \,\in\, n} A_{i}} \Bigg]
    \end{equation*}
    Thus, the unique set \( \setbuilder{x}{\Gamma(x)} \) exists by \cref{thm:sufficient-condition-for-unrestricted-comprehension}.
\end{proof}

\begin{dfn}[Set-Builder Notation]
    \label{dfn:set-builder-notation-for-tuples}
    If \( \varphi \) is an n-ary predicate where \( n \in \naturalset \) and \( n \neq 0 \),
    then \( \setbuilder{\tuple{x_{0}, \dots, x_{n-1}}}{\varphi(x_{0},\dots,x_{n-1})} \)
    denotes the set of all n-tuples whose coordinates satisfy \( \varphi \).
    \begin{multline*}
        \forall x \, \bigg[ x \in \setbuilderbig{\tuple{x_{0}, \dots, x_{n-1}}}{\varphi(x_{0},\dots,x_{n-1})} \iff \\
        \exists a_{0} \, \dots \exists a_{n-1} \, \Big( x = \tuple{a_{0}, \dots, a_{n-1}} \land \varphi(a_{0},\dots,a_{n-1}) \Big)
        \bigg]
    \end{multline*}
\end{dfn}

\begin{rem}
    The existence of \( \setbuilderbig{\tuple{x_{0}, \dots, x_{n-1}}}{\varphi(x_{0},\dots,x_{n-1})} \)
    can be established with \cref{thm:sufficient-condition-for-unrestricted-specification-for-tuples}.
\end{rem}

\begin{thm}
    \label{thm:the-exponentiation-of-a-set-A-with-a-natural-number-as-exponent-in-set-builder-notation}
    If \( A \) is a set and \( n \in \naturalset \), then,
    \begin{equation*}
        A^{n} =
        \begin{cases}
            \mleft\{ \emptyset \mright\} & \text{\normalfont if} \ n = 0 \\
            \setbuilderbig{\tuple{x_{0}, \dots, x_{n-1}}}{\Land_{i \,\in\, n} x_{i} \in A} & \text{\normalfont if} \ n \neq 0
        \end{cases}
    \end{equation*}
\end{thm}

\begin{proof}
    By \cref{dfn:set-exponentiation},
    \begin{equation}
        \label{eq:An-in-basic-setbuilder-form}
        A^{n} = \setbuilder{f}{\function{f}{n}{A}}
    \end{equation}
    If \( n = 0 \), then clearly,
    \begin{equation*}
        A^{0} = A^{\emptyset} = \setbuilder{f}{\function{f}{\emptyset}{A}} = \mleft\{ \tuple{} \mright\} = \mleft\{ \emptyset \mright\}
    \end{equation*}
    If \( n \neq 0 \), we can apply \cref{dfn:n-tuples-sequence-notation} to \eqref{eq:An-in-basic-setbuilder-form}.
    First, define the predicate \( \varphi \) as follows:
    \begin{equation}
        \label{eq:predicate-for-tuples-with-coordinates-belonging-to-A}
        \varphi(f) \coloneq \exists x_{0} \, \dots \exists x_{n-1} \, \bigg[ f = \tuple{x_{0}, \dots, x_{n-1}} \land \Land_{i \,\in\, n} x_{i} \in A \bigg]
    \end{equation}
    Let \( f  \) be an arbitrary object. Suppose that \( f \in A^{n} \). Thus, \( \function{f}{n}{A} \) by \eqref{eq:An-in-basic-setbuilder-form}.
    Since \( n \neq \emptyset \) and \( \dom f = n \), it follows from \cref{dfn:domain-of-a-set-of-ordered-pairs} that
    \( \ran f \neq \emptyset \). In other words, there exist objects \( x_{0},\dots,x_{n-1} \) such that \( \ran f = \mleft\{ x_{0},\dots,x_{n-1} \mright\} \).
    Since \( n \in \naturalset \), it follows from \cref{dfn:n-tuples} that
    \( f = \tuple{x_{0}, \dots, x_{n-1}} \). Also, since \( x_{0},\dots,x_{n-1} \in \ran f \) and
    \( \ran f \subseteq A \), necessarily, \( \Land_{i \,\in\, n} x_{i} \in A \). In other words, \( \varphi(f) \) is satisfied.
    By the \aref{rule-of-inference:deduction-theorem}{Deduction Theorem},
    \begin{equation}
        \label{eq:f-in-An-implies-phi-f}
        \begin{aligned}
            f \in A^{n} \implies \varphi(f)&
        \end{aligned}
        \begin{aligned}
            &\text{for any \( f \)}
        \end{aligned}
    \end{equation}
    Conversely, suppose that \( \varphi(f) \) is satisfied. By \eqref{eq:predicate-for-tuples-with-coordinates-belonging-to-A},
    \( f \) is an n-tuple \( \tuple{x_{0}, \dots, x_{n-1}} \) where \( \Land_{i \,\in\, n} x_{i} \in A \).
    By \cref{dfn:n-tuples}, it follows that \( \forall i\,{\in}\,n \, \big( x_{i} = f(i) \big) \).
    Given that \( \Land_{i \,\in\, n} x_{i} \in A \), necessarily \( \forall i\,{\in}\,n \, \big( f(i) \in A \big) \).
    In other words, \( f \) is a function with \( \dom f = n \) and \( \ran f \subseteq A \).
    By \cref{dfn:arrow-notation}, it follows that \( \function{f}{n}{A} \). And finally by \eqref{eq:An-in-basic-setbuilder-form},
    we have \( f \in A^{n} \).
    \begin{equation}
        \label{eq:phi-f-implies-f-in-An}
        \begin{aligned}
            \varphi(f) \implies f \in A^{n}&
        \end{aligned}
        \begin{aligned}
            &\text{for any \( f \)}
        \end{aligned}
    \end{equation}
    Having derived \eqref{eq:f-in-An-implies-phi-f} and \eqref{eq:phi-f-implies-f-in-An}, we further derive 
    \( \varphi(f) \iff f \in A^{n} \). Since \( f \) is arbitrary, we conclude by
    \aref{rule-of-inference:universal-generalization}{Universal Generalization} and the \nameref{axm:axiom-of-extensionality} that,
    \begin{equation*}
        A^{n}
        = \setbuilderbigg{f}{\varphi(f)}
        = \setbuilderbigg{f}{\exists x_{0} \, \dots \exists x_{n-1} \, \Big[ f = \tuple{x_{0}, \dots, x_{n-1}} \land \Land_{i \,\in\, n} x_{i} \in A \Big]}
    \end{equation*}
    And finally, \cref{dfn:set-builder-notation-for-tuples} allows us to express the result using set-builder notation:
    \begin{equation*}
        A^{n} = \setbuilderbigg{\tuple{x_{0}, \dots, x_{n-1}}}{\Land_{i \,\in\, n} x_{i} \in A} \qedhere
    \end{equation*}
\end{proof}

\begin{dfn}[N-ary Product]
    \label{dfn:n-ary-product}
    Let \( A_{0},\dots,A_{n-1} \) be \( n \) sets where \( n \) is a non-zero natural number.
    The \keyword{n-ary product} (or simply the \keyword{product}) of \( A_{0},\dots,A_{n-1} \)
    is the set \( \prod \tuple{A_{0}, \dots, A_{n-1}} \).
\end{dfn}

\begin{rem}
    Since \( \tuple{A_{0}, \dots, A_{n-1}} \) is an indexed system of sets with \( I = n \),
    then by \cref{dfn:product-of-indexed-system-of-sets}, the n-ary product of \( A_{0},\dots,A_{n-1} \)
    is given by,
    \begin{equation*}
        \prod \tuple{A_{0}, \dots, A_{n-1}} =
        \setbuilderbigg{\function{f}{n}{\Union_{i \,\in\, n} A_{i}}}{\forall i \, \big( i \in n \implies f(i) \in A_{i} \big)}
    \end{equation*}
    Let \( f \in \prod \tuple{A_{0}, \dots, A_{n-1}} \) be arbitrary. Thus, \( \function{f}{n}{\Union_{i \,\in\, n} A_{i}} \)
    and \( \forall i \, \big( i \in n \implies f(i) \in A_{i} \big) \) by definition. Since \( f \) is a function on \( n \)
    where \( n \in \naturalset \), \( f \) must be an n-tuple by \cref{dfn:n-tuples}. Since \( n \neq 0 \), \( f \)
    can be expressed in sequence notation by \cref{thm:every-non-zero-tuple-is-expressible-in-sequence-notation}.
    Also, \( f(i) = x_{i} \in A_{i} \) for any \( i \in n \).
    Thus,
    \begin{equation*}
        \exists x_{0} \, \dots \exists x_{n-1} \, \bigg[ f = \tuple{x_{0}, \dots, x_{n-1}} \land \Land_{i \,\in\, n} x_{i} \in A_{i} \bigg]
    \end{equation*}
    By \cref{dfn:set-builder-notation-for-tuples},
    \begin{equation*}
        f \in \prod \tuple{A_{0}, \dots, A_{n-1}} \implies f \in \setbuilderbigg{\tuple{x_{0}, \dots, x_{n-1}}}{\Land_{i \,\in\, n} x_{i} \in A_{i}}
    \end{equation*}
    Since the converse is also true (obvious), we can conclude that,
    \begin{equation*}
        \prod \tuple{A_{0}, \dots, A_{n-1}} = \setbuilderbigg{\tuple{x_{0}, \dots, x_{n-1}}}{\Land_{i \,\in\, n} x_{i} \in A_{i}}
    \end{equation*}
    In other words, the n-ary product of \( A_{0},\dots,A_{n-1} \) is the set of all n-tuples \( \tuple{x_{0}, \dots, x_{n-1}} \)
    such that \( \Land_{i \,\in\, n} x_{i} \in A_{i} \).
    Also, for the special case where \( \Land_{i \,\in\, n} A_{i} = A \) for some set \( A \), then,
    \( \surjection{\tuple{A_{0}, \dots, A_{n-1}}}{n}{\mleft\{ A \mright\}} \),
    in which case \( \prod \tuple{A_{0}, \dots, A_{n-1}} = A^{n} \) by \cref{thm:product-of-an-indexed-singleton}.
\end{rem}

\begin{thm}
    \label{thm:n-ary-product-is-empty-if-and-only-if-one-of-the-sets-is-empty}
    If \( A_{0},\dots,A_{n-1} \) are \( n \) sets where \( n \neq 0 \), then \( \prod \tuple{A_{0}, \dots, A_{n-1}} = \emptyset \)
    if and only if \( \exists i\,{\in}\,n \, \big( A_{i} = \emptyset \big) \).
    \begin{equation*}
        \forall n\,{\in}\,\naturalset \,
        \Big[ n \neq 0 \implies \forall A_{0} \, \dots \forall A_{n-1} \, \Big( \prod \tuple{A_{0}, \dots, A_{n-1}} = \emptyset \iff
        \exists i\,{\in}\,n \, \mleft( A_{i} = \emptyset \mright)\Big) \Big]
    \end{equation*}
\end{thm}

\begin{proof}
    Let \( A_{0},\dots,A_{n-1} \) be \( n \) sets for some \( n \neq 0 \). If \( \prod \tuple{A_{0}, \dots, A_{n-1}} = \emptyset \),
    then we can derive the following biconditionals using \cref{dfn:n-ary-product},
    \begin{align}
        &\forall f \, \bigg[ f \notin \prod \tuple{A_{0}, \dots, A_{n-1}} \bigg] \notag \\
        &\iff
        \forall f \, \bigg[ \lnot \, \Big( \function{f}{n}{\Union_{i \,\in\, n} A_{i}} \land \forall i \, \big( i \in n \implies f(i) \in A_{i} \big) \Big) \bigg]
        \notag \\
        &\iff
        \forall f \, \bigg[ \function{f}{n}{\Union_{i \,\in\, n} A_{i}} \implies \lnot \, \forall i \, \big( i \in n \implies f(i) \in A_{i} \big) \bigg] \notag \\
        &\iff
        \forall f \, \bigg[ \function{f}{n}{\Union_{i \,\in\, n} A_{i}} \implies \exists i \, \big( i \in n \land f(i) \notin A_{i} \big) \bigg]
        \label{eq:main-biconditional-for-empty-product}
    \end{align}
    Suppose that \( \exists i\,{\in}\,n \, \big( A_{i} = \emptyset \big) \).
    Thus, \( A_{k} = \emptyset \) for some \( k \in n \).
    Next, let \( f \) be an any set such that \( \function{f}{n}{\Union_{i \,\in\, n} A_{i}} \).
    Since \( k \in n \) and \( f \) is a function on \( n \), there must exist a \( y \in \Union_{i \,\in\, n} A_{i}\) such that \( y = f(k) \).
    Since \( A_{k} = \emptyset \), it follows that \( f(k) \notin A_{k} \). In other words, \( \exists i \, \big( i \in n \land f(i) \notin A_{i} \big) \).
    By the \aref{rule-of-inference:deduction-theorem}{Deduction Theorem},
    \begin{equation*}
        \function{f}{n}{\Union_{i \,\in\, n} A_{i}} \implies \exists i \, \big( i \in n \land f(i) \notin A_{i} \big)
    \end{equation*}
    Since \( f \) is arbitrary, this can be generalized,
    \begin{equation}
        \label{eq:quantified-consequent-1}
        \forall f \, \Big[ \function{f}{n}{\Union_{i \,\in\, n} A_{i}} \implies \exists i \, \big( i \in n \land f(i) \notin A_{i} \big) \Big]
    \end{equation}
    Since \eqref{eq:quantified-consequent-1} is derived from the assumption that \( \exists i\,{\in}\,n \, \big( A_{i} = \emptyset \big) \),
    we conclude by the \aref{rule-of-inference:deduction-theorem}{Deduction Theorem} that,
    \begin{equation}
        \label{eq:at-least-one-empty-Ai-implies-empty-product}
        \exists i\,{\in}\,n \, \big( A_{i} = \emptyset \big) \implies
        \forall f \, \Big[ \function{f}{n}{\Union_{i \,\in\, n} A_{i}} \implies \exists i \, \big( i \in n \land f(i) \notin A_{i} \big) \Big]
    \end{equation}
    To prove the converse, we suppose \( \lnot \, \exists i\,{\in}\,n \, \big( A_{i} = \emptyset \big) \), which is logically equivalent to,
    \begin{equation*}
        \forall i\,{\in}\,n \, \big( A_{i} \neq \emptyset \big)
    \end{equation*}
    Since each set \( A_{0},\dots,A_{n-1} \) contains at least one element,
    it is possible to construct an n-tuple \( f = \tuple{a_{0}, \dots, a_{n-1}} \) where \( \Land_{i \,\in\, n} a_{i} \in A_{i} \).
    From the definition of tuples, it is clear that,
    \begin{equation*}
        \function{f}{n}{\Union_{i \,\in\, n} A_{i}} \land \forall i \, \big( i \in n \implies f(i) \in A_{i} \big)
    \end{equation*}
    which is logically equivalent to,
    \begin{align*}
        &\function{f}{n}{\Union_{i \,\in\, n} A_{i}} \land \lnot \, \exists i \, \big( i \in n \land f(i) \notin A_{i} \big) \\
        &\iff
        \lnot \, \Big[ \function{f}{n}{\Union_{i \,\in\, n} A_{i}} \implies \exists i \, \big( i \in n \land f(i) \notin A_{i} \big) \Big]
    \end{align*}
    By \aref{rule-of-inference:existential-generalization}{Existential Generalization}, this becomes,
    \begin{align}
        &\exists f \, \lnot \, \Big[ \function{f}{n}{\Union_{i \,\in\, n} A_{i}} \implies \exists i \, \big( i \in n \land f(i) \notin A_{i} \big) \Big] \notag \\
        &\iff
        \lnot \, \forall f \, \Big[ \function{f}{n}{\Union_{i \,\in\, n} A_{i}} \implies \exists i \, \big( i \in n \land f(i) \notin A_{i} \big) \Big]
        \label{eq:quantified-consequent-2}
    \end{align}
    Since \eqref{eq:quantified-consequent-2} is derived from the assumption that \( \lnot \, \exists i\,{\in}\,n \, \big( A_{i} = \emptyset \big) \), we conclude
    by the \aref{rule-of-inference:deduction-theorem}{Deduction Theorem} that,
    \begin{equation*}
        \lnot \, \exists i\,{\in}\,n \, \big( A_{i} = \emptyset \big)
        \implies \lnot \, \forall f \, \Big[ \function{f}{n}{\Union_{i \,\in\, n} A_{i}} \implies \exists i \, \big( i \in n \land f(i) \notin A_{i} \big) \Big]
    \end{equation*}
    By \aref{rule-of-inference:contraposition}{Contraposition}, this is logically equivalent to,
    \begin{equation}
        \label{eq:empty-product-implies-at-least-one-empty-Ai}
        \forall f \, \Big[ \function{f}{n}{\Union_{i \,\in\, n} A_{i}} \implies \exists i \, \big( i \in n \land f(i) \notin A_{i} \big) \Big]
        \implies \exists i\,{\in}\,n \, \big( A_{i} = \emptyset \big)
    \end{equation}
    Having established both \eqref{eq:at-least-one-empty-Ai-implies-empty-product} and \eqref{eq:empty-product-implies-at-least-one-empty-Ai},
    we can conclude the biconditional,
    \begin{equation}
        \label{eq:at-least-one-empty-Ai-and-empty-product-biconditional}
        \exists i\,{\in}\,n \, \big( A_{i} = \emptyset \big)
        \iff
        \forall f \, \Big[ \function{f}{n}{\Union_{i \,\in\, n} A_{i}} \implies \exists i \, \big( i \in n \land f(i) \notin A_{i} \big) \Big]
    \end{equation}
    Combining \eqref{eq:main-biconditional-for-empty-product} and \eqref{eq:at-least-one-empty-Ai-and-empty-product-biconditional}, we have,
    \begin{equation*}
        \prod \tuple{A_{0}, \dots, A_{n-1}} = \emptyset
        \iff
        \exists i\,{\in}\,n \, \Big[ A_{i} = \emptyset \Big]
    \end{equation*}
    Since \( n \) and \( A_{0},\dots,A_{n-1} \) are arbitrary, the theorem follows immediately from \aref{rule-of-inference:universal-generalization}{Universal Generalization}.
\end{proof}

\begin{thm}
    \label{thm:a-bijection-exists-between-A-and-A1}
    For any set \( A \), a bijection exists between \( A \) and \( A^{1} \).
    \begin{equation*}
        \forall A \, \exists h \, \Big[ \bijection{h}{A}{A^{1}} \Big]
    \end{equation*}
\end{thm}

\begin{proof}
    Let \( A \) be any set. Let the set \( h \) be defined as,
    \begin{equation}
        \label{eq:define-a-function-from-A-to-A1}
        h = \setbuilderBig{\orderedpair{a}{f}}{a \in A \land f = \mleft\{ \orderedpair{\emptyset}{a} \mright\}}
    \end{equation}
    We can prove that \( \bijection{h}{A}{A^{1}} \).

    \vspace{0.2cm}
    \noindent
    Let \( \alpha \), \( f_{1} \), and \( f_{2} \) be objects such that \( \orderedpair{\alpha}{f_{1}} \in h \)
    and \( \orderedpair{\alpha}{f_{2}} \in h \). By \eqref{eq:define-a-function-from-A-to-A1}, we have,
    \begin{gather*}
        \alpha \in A \land f_{1} = \mleft\{ \orderedpair{\emptyset}{\alpha} \mright\}\\
        \alpha \in A \land f_{2} = \mleft\{ \orderedpair{\emptyset}{\alpha} \mright\}
    \end{gather*}
    Obviously, \( f_{1} = f_{2} \), so by \cref{dfn:function-basic}, \( h \) is indeed a function. Next, we
    prove that \( \dom h = A \).
    Let \( \alpha \in A \) be arbitrary. Let \( f = \mleft\{ \orderedpair{\emptyset}{\alpha} \mright\} \).
    Then, by \eqref{eq:define-a-function-from-A-to-A1}, clearly \( \orderedpair{\alpha}{f} \in h \). So by
    \cref{dfn:domain-of-a-set-of-ordered-pairs}, it follows that \( \alpha \in \dom h \). Conversely,
    suppose that \( \alpha \in \dom h \) for some \( \alpha \). Again, \( a \in A \) by \eqref{eq:define-a-function-from-A-to-A1}.
    By \aref{rule-of-inference:universal-generalization}{Universal Generalization} and the \nameref{axm:axiom-of-extensionality}, we conclude that \( \dom h = A \).

    \vspace{0.2cm}
    \noindent
    Next, we prove that \( \ran h = A^{1} \). Let \( f \in \ran h \) be arbitrary. By \cref{dfn:range-of-a-set-of-ordered-pairs},
    there must exists an \( \alpha \) such that \( \orderedpair{\alpha}{f} \in h \). And by \eqref{eq:define-a-function-from-A-to-A1},
    this can only be true if and only if \( \alpha \in A \land f = \mleft\{ \orderedpair{\emptyset}{\alpha} \mright\} \).
    Since \( 1 = \mleft\{ \emptyset \mright\} \) and \( \alpha \in A \), it is clear that \( \function{f}{1}{A} \).
    By \cref{dfn:set-exponentiation}, this implies \( f \in A^{1} \).
    Conversely, let \( f \in A^{1} \) be arbitrary. Thus, \( \function{f}{1}{A} \) by \cref{dfn:set-exponentiation}.
    This means that \( f \) is a function with \( \dom f = 1 \) and \( \ran f \subseteq A \).
    Since \( \dom f = 1 = \mleft\{ \emptyset \mright\} \neq \emptyset \), it follows that \( \ran f \neq \emptyset \).
    Thus, there exists an \( \alpha \in A \) such that \( \orderedpair{\emptyset}{\alpha} \in f \).
    Since \( \emptyset \) is the only element of \( \dom f \), \( \orderedpair{\emptyset}{\alpha} \) must also be
    the only element of \( f \). In other words, \( f = \mleft\{ \orderedpair{\emptyset}{\alpha} \mright\} \) for some
    \( \alpha \in A \). And by \eqref{eq:define-a-function-from-A-to-A1}, clearly \( \orderedpair{\alpha}{f} \in h \),
    which means that \( f \in \ran h \). By \aref{rule-of-inference:universal-generalization}{Universal Generalization} and the \nameref{axm:axiom-of-extensionality},
    it follows that \( \ran h = A^{1} \).

    \vspace{0.2cm}
    \noindent
    So far, it has been proven that \( h \) is a function with \( \dom h = A \) and \( \ran h = A^{1} \). In other words,
    \begin{equation*}
        \surjection{h}{A}{A^{1}}
    \end{equation*}
    To prove that \( h \) is a bijection, it remains to show that \( h \) is injective. Let \( \alpha \), \( \beta \) and \( f \)
    be arbitrary objects such that \( \orderedpair{\alpha}{f} \in h \) and \( \orderedpair{\beta}{f} \in h \).
    By \eqref{eq:define-a-function-from-A-to-A1}, this means that,
    \begin{equation*}
        f = \mleft\{ \orderedpair{\emptyset}{\alpha} \mright\} = \mleft\{ \orderedpair{\emptyset}{\beta} \mright\}
    \end{equation*}
    Thus, \( \orderedpair{\emptyset}{\alpha} = \orderedpair{\emptyset}{\beta} \). By \cref{thm:equality-of-ordered-pairs},
    it follows that \( \alpha = \beta \), so \( h \) is necessarily injective.
    \begin{equation*}
        \therefore \, \surjection{h}{A}{A^{1}} \land \injection{h}{A}{A^{1}} \iff \bijection{h}{A}{A^{1}} \qedhere
    \end{equation*}
\end{proof}

\begin{thm}
    \label{thm:a-bijection-exists-between-A-times-A-and-A2}
    For any set \( A \), a bijection exists between \( A \times A \) and \( A^{2} \).
    \begin{equation*}
        \forall A \, \exists h \, \Big[ \bijection{h}{A \times A}{A^{2}} \Big]
    \end{equation*}
\end{thm}

\begin{proof}
    Let \( A \) be an arbitrary set. Let the set \( h \) be defined as,
    \begin{equation}
        \label{eq:define-a-function-from-A-times-A-to-A2}
        h = \setbuilderBig{\orderedpair{x}{f}}{\exists a \, \exists b \,
        \big[ a \in A \land b \in A \land x = \orderedpair{a}{b}
        \land
        f = \mleft\{ \orderedpair{0}{a}, \orderedpair{1}{b} \mright\} \big]}
    \end{equation}
    We can prove that \( \bijection{h}{A \times A}{A^{2}} \).

    \vspace{0.3cm}
    \noindent
    Let \( x \), \( f_{1} \), and \( f_{2} \) be arbitrary objects such that \( \orderedpair{x}{f_{1}} \in h \)
    and \( \orderedpair{x}{f_{2}} \in h \). By \eqref{eq:define-a-function-from-A-times-A-to-A2}, we have for some
    \( \alpha,\beta \in A \):
    \begin{align*}
        x &= \orderedpair{\alpha}{\beta} \\
        f_{1} &= \big\{ \orderedpair{0}{\alpha}, \orderedpair{1}{\beta} \big\} \\
        f_{2} &= \big\{ \orderedpair{0}{\alpha}, \orderedpair{1}{\beta} \big\}
    \end{align*}
    Clearly, \( f_{1} = f_{2}\). So \( h \) must be a function by \cref{dfn:function-basic}.

    \vspace{0.3cm}
    \noindent
    We proceed to prove that \( \dom h = A \times A \).

    \vspace{0.3cm}
    \noindent
    Let \( x_{0} \in A \times A \) be arbitrary. Thus, \( x_{0} = \orderedpair{\alpha}{\beta} \) for some \( \alpha,\beta \in A \).
    Let \( f \) be the set \( \big\{ \orderedpair{0}{\alpha}, \orderedpair{1}{\beta} \big\} \).
    By \eqref{eq:define-a-function-from-A-times-A-to-A2}, it follows that \( \orderedpair{x_{0}}{f} \in h \) so that \( x_{0} \in \dom h \).
    Thus, \( A \times A \subseteq \dom h \).
    Conversely, let \( x_{0} \in \dom h \). By \eqref{eq:define-a-function-from-A-times-A-to-A2}, this means that \( \orderedpair{x_{0}}{f} \in h \)
    for some \( f \) where \( x_{0} = \orderedpair{\alpha}{\beta} \) and \( f = \big\{ \orderedpair{0}{\alpha}, \orderedpair{1}{\beta} \big\} \) for
    some \( \alpha,\beta \in A \).
    Since \( x_{0} = \orderedpair{\alpha}{\beta} \) and \( \alpha,\beta \in A \), it follows that \( x_{0} \in A \times A \).
    Therefore, \( \dom h \subseteq A \times A \). Since \( \dom h \subseteq A \times A \) and \( A \times A \subseteq \dom h \) are both true,
    it follows from \cref{thm:set-equality-in-terms-of-subsets} that \( \dom h = A \times A \).

    \vspace{0.3cm}
    \noindent
    To prove that \( \ran h = A^{2} \), we let \( f \in \ran h \) be arbitrary.
    By \cref{dfn:range-of-a-set-of-ordered-pairs}, it it follows that \( \orderedpair{x_{0}}{f} \in h \)
    for some \( x_{0} \) where \( x_{0} = \orderedpair{\alpha}{\beta} \) and \( f = \big\{ \orderedpair{0}{\alpha},\orderedpair{1}{\beta} \big\} \)
    for some \( \alpha,\beta \in A \).
    \( f \) is clearly a function with \( \dom f = \mleft\{ 0,1 \mright\} \) and \( \ran f = \mleft\{ \alpha, \beta \mright\} \).
    And since \( 2 = \mleft\{ 0,1 \mright\} \) and \( \alpha,\beta \in A \), necessarily \( \dom f = 2 \) and \( \ran f \subseteq A \).
    Thus, \( \function{f}{2}{A} \) by \cref{dfn:arrow-notation} and \( f \in A^{2} \) by \cref{dfn:set-exponentiation}. This means that for any
    arbitrary \( f \),
    \begin{equation*}
        f \in \ran h \implies f \in A^{2}
    \end{equation*}
    Conversely, if \( f \) is any arbitrary element of \( A^{2} \), then \( \function{f}{2}{A} \) by \cref{dfn:set-exponentiation}.
    Since \( \dom f = 2 \) and \( 2 = \mleft\{ 0,1 \mright\} \), it follows that \( f = \big\{ \orderedpair{0}{\alpha},\orderedpair{1}{\beta} \big\} \)
    for some \( \alpha,\beta \in A \) because \( \ran f \subseteq A \). If we now let \( x_{0} = \orderedpair{\alpha}{\beta} \), then
    \( \orderedpair{x_{0}}{f} \) must belong to \( h \) by \eqref{eq:define-a-function-from-A-times-A-to-A2}.
    This means that there exists an \( x_{0} \) such that \( \orderedpair{x_{0}}{f} \in h \). By \cref{dfn:range-of-a-set-of-ordered-pairs},
    \( f \) must them belong to \( \ran h \). In other words, for any arbitrary \( f \),
    \begin{equation*}
        f \in A^{2} \implies f \in \ran h
    \end{equation*}
    Having established the inclusion in both directions, we can now conclude that \( \ran h = A^{2} \). Thus, \( h \) is a function with
    \( \dom h = A \times A \) and \( \ran h = A^{2} \). In other words,
    \begin{equation*}
        \surjection{h}{A \times A}{A^{2}}
    \end{equation*}
    To prove that \( h \) is a bijection, it remains to show that \( h \) is injective. Let \( x_{1} \), \( x_{2} \) and \( f \) be
    arbitrary objects such that \( \orderedpair{x_{1}}{f} \in h \) and \( \orderedpair{x_{2}}{f} \in h \).
    From \eqref{eq:define-a-function-from-A-times-A-to-A2}, we infer that,
    \begin{equation*}
        \begin{aligned}
            x_{1} = \orderedpair{a_{1}}{b_{1}} \land f = \big\{ \orderedpair{0}{a_{1}},\orderedpair{1}{b_{1}} \big\}&\\
            x_{2} = \orderedpair{a_{2}}{b_{2}} \land f = \big\{ \orderedpair{0}{a_{2}},\orderedpair{1}{b_{2}} \big\}&
        \end{aligned}
        \text{}
        \begin{aligned}
            &\quad\text{for some \( a_{1},b_{1} \in A \)}\\
            &\quad\text{for some \( a_{2},b_{2} \in A \)}
        \end{aligned}
    \end{equation*}
    Clearly, \( f = \big\{ \orderedpair{0}{a_{1}},\orderedpair{1}{b_{1}} \big\} = \big\{ \orderedpair{0}{a_{2}},\orderedpair{1}{b_{2}} \big\}\).
    The \nameref{axm:axiom-of-extensionality} guarantees that the follows statement is true:
    \begin{equation*}
        \orderedpair{0}{a_{1}} = \orderedpair{0}{a_{2}} \lor
        \orderedpair{0}{a_{1}} = \orderedpair{1}{b_{2}}
    \end{equation*}
    Since \( 0 \neq 1 \), the second disjunct must be false.
    Thus, by \aref{rule-of-inference:disjunctive-syllogism}{Disjunctive Syllogism}, we conclude that \( \orderedpair{0}{a_{1}} = \orderedpair{0}{a_{2}}\).
    By \cref{thm:equality-of-ordered-pairs}, this implies that \( a_{1} = a_{2} \). Similarly, we have
    \( \orderedpair{1}{b_{1}} = \orderedpair{1}{b_{2}} \), which in turns implies \( b_{1} = b_{2} \).
    And with \( a_{1} = a_{2} \) and \( b_{1} = b_{2} \) both being true,
    it follows that \( x_{1} = \orderedpair{a_{1}}{b_{1}} = \orderedpair{a_{2}}{b_{2}} = x_{2} \).
    Thus, \( h \) is indeed injective by \cref{dfn:injection}.
\end{proof}

\begin{rem}
    The significance of \cref{thm:a-bijection-exists-between-A-and-A1,thm:a-bijection-exists-between-A-times-A-and-A2}
    will only become clear in the context of isomorphisms.
\end{rem}

\section{Relations}
Select any two pieces from a chess set. Now try to describe each piece \emph{in relation} to the other.
What immediately comes to mind?
We might note, for instance, that one of the selected pieces is physically taller than the other.
Their colors may also be different, with one having a brighter hue than the other. Physical characteristics aside,
the pieces may also differ in material worth as it pertains to the game of chess---e.g., a Queen is worth more than a Pawn, a Rook
more than a Bishop, etc. In all of these examples, the properties of the selected pieces were compared to determine whether
they satisfy a given condition. If the condition is met, then we make an assertion along the lines of ``chess piece \( A \) has an
\( X \) relationship with chess piece \( B \)'', \( X \) being a placeholder for ``is taller than,'' ``has a brighter
color than,'' and ``is more valuable than'' respectively.
The symbol `\( X \)' denotes a \emph{relation} between two objects,
and the proposition ``\( A \) has an \( X \) relationship with \( B \)'' is usually abbreviated as \( X(A,B) \) or \( AXB \).

\keyconcept{
    A \emph{relation} is a symbolic representation of a condition that an object(s) can satisfy. It is a \emph{predicate}.
}

\noindent
A relation can involve any number of objects, of course.
In general, the arity of the predicate corresponds to the number of objects involved in that relation.
For example, the ``is taller than'' relationship involves exactly two objects,
so a binary predicate can be used to represent it.

\vspace{0.3cm}
\noindent
The order in which the objects are specified to the predicate also matters.
As in, if \( \alpha \) and \( \beta \) are two distinct
objects in the universe of discourse and \( \varphi \) is a binary predicate,
then \( \varphi(\alpha, \beta) \) does not necessarily imply \( \varphi(\beta,\alpha) \).
For example, if \( \varphi \) represents the ``is taller than'' relationship and \( \varphi(\alpha,\beta) \)
is true, then \( \varphi(\beta,\alpha) \) must be false by virtue of the semantics of the ``is taller than''
relationship.

\vspace{0.3cm}
\noindent
So far, so good.
Conceptually, we know that a relation is synonymous to a predicate.
But how can we express this in terms of sets? To answer
this question, let us work through another example.
Consider a set \( S \) of all students in a classroom.
Suppose we want a set \( R \) that describes the
relationship ``\( x \) is taller than \( y \).''\,What information should
\( R \) contain so that it can serve as an adequate representation of the relation
``\( x \) is taller than \( y \)?'' The answer should be obvious. The set
\( R \) must be a list of all the student pairs \( \tuple{x,y} \)
in which student \( x \) is physically taller than student \( y \)!

\vspace{0.3cm}
\noindent
In general, for a set \( R \) to represent a relation between \( n \) objects from a set,
each element of \( R \) must be an ordered list of \( n \) objects that satisfy the given relation.
In other words, each element of \( R \) must be a mathematical construct that implements
the notion of an ordered list of \( n \) objects. But such a mathematical construct has already been defined---the ordered n-tuples.
Therefore, if \( R \) represents a relation between \( n \) objects from a set \( A \), then \( R \) is
a set of n-tuples \( \tuple{a_{0}, \dots, a_{n-1}} \) formed from the elements of \( A \), which
leads us to the following definition.

\begin{dfn}[N-Ary Relation]
    \label{dfn:n-ary-relation-on-a-set}
    For any set \( A \) and any \( n \in \naturalset \), \( R \) is an \keyword{n-ary relation} on \( A \) if and only if \( R \subseteq A^{n} \).
    The natural number \( n \) is called the \keyword{arity} of \( R \).
\end{dfn}

\begin{rem}
    An n-ary relation is always a subset of \( A^{n} \) because by
    \cref{thm:the-exponentiation-of-a-set-A-with-a-natural-number-as-exponent-in-set-builder-notation},
    \( A^{n} \) is the set of all possible n-tuples that can be formed using the elements of \( A \).
\end{rem}

\begin{notation}[Arity]
    \phantomsection \label{notation:arity-of-n-ary-relations}
    If \( R \) is a relation on \( A \) we write \( \arity{R} \) to denote the arity of \( R \).
\end{notation}

\subsection{Binary Relations}
Having precisely defined n-ary relations, we now turn to the most common type of relation in mathematics---binary relations.
\begin{dfn}
    \label{dfn:binary-relation-on-a-set}
    For any set \( A \), \( R \) is a \keyword{binary relation} on \( A \) if and only if
    \( R \subseteq A^{2} \).
\end{dfn}

\begin{dfn}[Identity Relation]
    \label{dfn:identity-relation}
    For any set \( A \), the \keyword{identity relation} on \( A \) is a binary relation defined as,
    \begin{equation*}
        \identityrel{A} = \setbuilderbig{\tuple{x,y}}{x \in A \land y \in A \land x = y}
    \end{equation*}
\end{dfn}
\noindent
Here are some important properties that are often used to classify binary relations.

\begin{dfn}[Reflexivity]
    \label{dfn:reflexivity-of-binary-relations}
    If \( R \) is a binary relation on \( A \), then \( R \) is \keyword{reflexive}
    if and only if,
    \begin{equation*}
        \forall x\,{\in}\,A \, \Big[ \tuple{x,x} \in R \Big]
    \end{equation*}
\end{dfn}

\begin{dfn}[Symmetry]
    \label{dfn:symmetry-of-binary-relations}
    If \( R \) is a binary relation on \( A \), then \( R \) is \keyword{symmetric}
    if and only if,
    \begin{equation*}
        \forall x\,{\in}\,A \, \forall y\,{\in}\,A \, \Big[ \tuple{x,y} \in R \implies \tuple{y,x} \in R \Big]
    \end{equation*}
\end{dfn}

\begin{dfn}[Anti-Symmetry]
    \label{dfn:anti-symmetry-of-binary-relations}
    If \( R \) is a binary relation on \( A \), then \( R \) is \keyword{anti-symmetric} if and only if,
    \begin{equation*}
        \forall x\,{\in}\,A \, \forall y\,{\in}\,A \, \Big[ \tuple{x,y} \in R \land \tuple{y,x} \in R \implies x = y \Big]
    \end{equation*}
\end{dfn}

\begin{dfn}[Asymmetry]
    \label{dfn:asymmetry-of-binary-relations}
    If \( R \) is a binary relation on \( A \), then \( R \) is \keyword{asymmetric} if and only if,
    \begin{equation*}
        \forall x\,{\in}\,A \, \forall y\,{\in}\,A \, \Big[ \tuple{x,y} \in R \implies \tuple{y,x} \notin R \Big]
    \end{equation*}
\end{dfn}

\begin{rem}
    The difference between \keyword{anti-symmetry} and \keyword{asymmetry} lies in the fact that \( \tuple{x,x} \in R \)
    is possible in the former but impossible in the latter. Thus, any binary relation that is asymmetric must also be
    \keyword{irreflexive} (not reflexive).
\end{rem}

\begin{dfn}[Transitivity]
    \label{dfn:transitivity-of-binary-relations}
    If \( R \) is a binary relation on \( A \), then \( R \) is \keyword{transitive} if and only if,
    \begin{equation*}
        \forall x\,{\in}\,A \, \forall y\,{\in}\,A \, \forall z\,{\in}\,A \,
        \Big[ \orderedpair{x}{y} \in R \land \orderedpair{y}{z} \in R \implies \orderedpair{x}{z} \in R \Big]
    \end{equation*}
\end{dfn}

\subsubsection{Equivalence Relations}
There are numerous situations in which it is convenient and even advantageous to
ignore the differences between two objects and treat them as though they are one
and the same even though they are, in fact, distinct objects.
For example, imagine trying to explain to a novice how the Knights move in the game of chess.
Is there a point in providing a different explanation for how the white Knight moves
versus how the black Knight moves?
Obviously not.
After all, the color of a chess piece has no bearing whatsoever on its designated movement.
Therefore, for the purpose of explaining the rules of chess to a novice player, the differences
between a white Knight and a black Knight can be disregarded because them are irrelevant.
There are many other examples to this effect:
\begin{itemize}[noitemsep]
    \item \( \frac{1}{2} \) and \( \frac{3}{6} \) can be treated as the same fraction even though they are denoted by different symbols.
    \item In a card game, the same rules apply regardless of the pack of cards being used; the King of Spades from one pack is no different
          from the King of Spades from another pack.
\end{itemize}
In general, when we say that ``\( A \) is equivalent to \( B \),'' what we're really saying is that
\( A \) and \( B \) are \emph{effectively equal} to one another. More precisely, it means that \( A \)
and \( B \) can be treated as though they are the same object even though they may, in fact, be different objects.
We also note the following:
\vspace{0cm}
\begin{enumerate}[topsep=2pt,noitemsep,label=\texttt{\textup{(\alph*)}}]
    \item \( A \) is equivalent to \( A \)
    \item If \( A \) is equivalent to \( B \), then \( B \) must also be equivalent to \( A \)
    \item If \( A \) is equivalent to \( B \) and \( B \) to \( C \), then \( A \) must be equivalent to \( C \)
\end{enumerate}
These properties arise from the meaning ascribed to the
``is equivalent to'' relationship.
Recall that ``\( A \) is equivalent to \( B \)'' means we are treating \( A \)
and \( B \) as the same object while ignoring all potential differences between them.
Another notable feature of the ``is equivalent to'' relationship is that it partitions
a set into non-empty, pairwise disjoint sets.

\vspace{0.3cm}
\noindent
The following definitions formalize the above discussion.

\begin{dfn}[Equivalence Relation]
    \label{dfn:equivalence-relation}
    For any set \( A \), the set \( E \subseteq A^{2} \) is an \keyword{equivalence relation} on \( A \)
    if and only if it is reflexive, symmetric, and transitive. Moreover, we say that \( x \) is equivalent to
    \( y \) modulo \( E \) if and only if \( \tuple{x,y} \in E \).
\end{dfn}

\begin{dfn}[Equivalence Class]
    \label{dfn:equivalence-class}
    If \( E \subseteq A^{2} \) is an equivalence relation on \( A \) and \( \alpha \) is an element of \( A \),
    then the \keyword{equivalence class} of \( \alpha \) modulo \( E \) is the set:
    \begin{equation*}
        \equivclass{\alpha}_{E} = \setbuilderbig{x \in A}{\tuple{x,\alpha} \in E}
    \end{equation*}
\end{dfn}

\begin{rem}
    When the context makes it clear, we can use the notation \( a \equivrel b \) for \( \tuple{a,b} \in E \) and
    \( \equivclass{\alpha} \) for \( \equivclass{\alpha}_{E} \).
\end{rem}

\begin{cor}
    \label{thm:existence-of-equivalence-classes}
    The equivalence class of \( \alpha \) exists for any \( \alpha \in A \) because \( E \subseteq A^{2} \) by definition.
    As in,
    \begin{equation*}
        \forall x \, \Big[ \tuple{x,\alpha} \in E \implies x \in A \Big]
    \end{equation*}
\end{cor}

\begin{cor}
    \label{thm:equivalence-class-belongs-to-the-power-set}
    For any \( \alpha \in A \), \( \equivclass{\alpha}_{E} \in \powerset{A} \).
\end{cor}

\vspace{0.2cm}
\noindent
Now that we have given a precise definition for equivalence relations, we can define other mathematical constructs
in terms of it.

\begin{dfn}[Quotient Set]
    \label{dfn:quotient-set}
    For any set \( A \) and any equivalence relation \( E \) on \( A \), the \keyword{quotient set} of \( A \) with respect to
    \( E \) is defined as,
    \begin{equation*}
        A / E = \setbuilderbig{\equivclass{x}_{E}}{x \in A}
    \end{equation*}
\end{dfn}

\begin{rem}
    Since \( A / E \) is a division of \( A \) into partitions, the term \keyword{quotient} is appropriate.
\end{rem}

\begin{dfn}[Partition]
    \label{dfn:partition-of-a-set}
    A system of sets \( P \) is a \keyword{partition} of another set \( A \) if and only if the following
    conditions hold:
    \begin{enumerate}[topsep=2pt,label=\texttt{\textup{(\alph*)}}]
        \item \( \forall x\,{\in}\,P \, \big( x \neq \emptyset \big) \)
        \item \( \forall x\,{\in}\,P \, \forall y\,{\in}\,P \, \big( x \neq y \implies x \intersect y = \emptyset \big) \)
        \item \( \Union P = A \)
    \end{enumerate}
\end{dfn}

\begin{dfn}
    \label{dfn:set-of-all-partitions-of-a-set-and-the-set-of-all-equivalence-relations-on-a-set}
    For any set \( A \), we denote the set of all partitions of \( A \) as \partitions{A} and
    the set of all equivalent relations on \( A \) as \( \equivrels{A} \).
\end{dfn}

\begin{dfn}[Class Invariance]
    \label{dfn:class-invariance}
    Let \( A \) be a set and \( E \) be an equivalence relation on \( A \).
    If \( f \) is a function on \( A \), then \( f \) is \keyword{class invariant}
    under \( E \) if and only if,
    \begin{equation*}
        \forall x\,{\in}\,A \, \forall y\,{\in}\,A \, \Big[ \tuple{x,y} \in E \implies f(x) = f(y) \Big]
    \end{equation*}
\end{dfn}

\begin{rem}
    The term \keyword{class invariant} is apt because \( f \) is invariant for elements within the same
    equivalence class; i.e., picking a different element from the same equivalence class will not alter the
    function's value.
\end{rem}

\begin{dfn}[Equivalence Kernel]
    \label{dfn:equivalence-kernel-of-a-function}
    If \( f \) is a function, then the \keyword{equivalence kernel} of \( f \) is the equivalence relation
    \( K \subseteq \mleft( \dom f \mright)^{2} \) such that,
    \begin{equation*}
        \forall x\,{\in}\,\dom f \, \forall y\,{\in}\,\dom f \, \Big[ \tuple{x,y} \in K \iff f(x) = f(y) \Big]
    \end{equation*}
\end{dfn}
\noindent
Here are some important theorems involving equivalence relations.

\begin{thm}
    \label{thm:existence-of-the-quotient-set}
    For any set \( A \) and any equivalence relation \( E \) on \( A \), the quotient set \( A/E \) exists.
\end{thm}

\begin{proof}
    Let the predicate \( \varphi \) be defined as,
    \begin{equation*}
        \varphi(x) \coloneq \exists a\,{\in}\,A \, \big( x = \equivclass{a}_{E} \big)
    \end{equation*}
    Let \( x_{0} \) be an arbitrary object such that \( \varphi(x_{0}) \) is satisfied. That is,
    \( x_{0} = \equivclass{a}_{E} \) for some \( a \in A \).
    By \cref{thm:equivalence-class-belongs-to-the-power-set}, it follows that \( x_{0} \in \powerset{A} \).
    \begin{equation*}
        \therefore \, \forall x \, \big( \varphi(x) \implies x \in \powerset{A} \big)
    \end{equation*}
    By \cref{thm:sufficient-condition-for-unrestricted-comprehension}, the set
    \( A / E = \setbuilderbig{\equivclass{x}_{E}}{x \in A} \) exists.
\end{proof}

\begin{thm}
    \label{thm:existence-of-the-set-of-all-partitions}
    For any set \( A \), \( \partitions{A} \) exists.
\end{thm}

\begin{proof}
    Let the predicate \( \varphi \) be defined as,
    \begin{equation*}
        \varphi(x) \coloneq \text{\( x \) is a partition of \( A \)}
    \end{equation*}
    Let \( X \) be an arbitrary object such that \( \varphi(X) \) is satisfied.
    By \cref{dfn:partition-of-a-set}, it follows that \( \Union X = A \). Thus,
    by \cref{dfn:union-of-sets}, we have,
    \begin{equation}
        \label{eq:Union-X-equals-A}
        \forall x \, \Big[ \exists z \, \big( z \in X \land x \in z \big) \iff x \in A  \Big]
    \end{equation}
    Let \( z_{0} \in X \) be arbitrary. If \( x_{0} \in z_{0} \), then necessarily \( x_{0} \in A \)
    by \eqref{eq:Union-X-equals-A}. In other words, for any \( z_{0} \in X \),
    \( \forall x \, \big[ x \in z_{0} \implies x \in A \big] \) and thus \( z_{0} \subseteq A\).
    \begin{equation*}
        \therefore \, \forall z \, \Big[ z \in X \implies z \in \powerset{A} \Big]
    \end{equation*}
    In other words, \( X \subseteq \powerset{A} \) and thus \( X \in \powerset{\powerset{A}} \).
    By the \aref{rule-of-inference:deduction-theorem}{Deduction Theorem}, we conclude,
    \begin{equation*}
        \varphi(X) \implies X \in \powerset{\powerset{A}}
    \end{equation*}
    Since \( X \) is arbitrary, \aref{rule-of-inference:universal-generalization}{Universal Generalization} yields,
    \begin{equation*}
        \forall x \, \Big[ \varphi(x) \implies x \in \powerset{\powerset{A}} \Big]
    \end{equation*}
    Thus, the set \( \partitions{A} = \setbuilderbig{x}{\varphi(x)} \) exists
    by \cref{thm:sufficient-condition-for-unrestricted-comprehension}.
\end{proof}

\begin{thm}
    \label{thm:the-only-partition-for-the-empty-set-is-the-empty-set}
    \( \partitions{\emptyset} = \mleft\{ \emptyset \mright\} \)
\end{thm}

\begin{proof}
    Let \( P \) be any element of \( \partitions{\emptyset} \).
    By \cref{dfn:partition-of-a-set}, it follows that \( \Union P = \emptyset \).
    Thus, we have by \cref{dfn:empty-set,dfn:union-of-sets},
    \begin{equation*}
        \forall x \, \big[ \lnot \, \exists z\,{\in}\,P \, \big( x \in z \big) \big]
    \end{equation*}
    Which is equivalent to,
    \begin{equation}
        \label{eq:the-union-of-P-is-empty}
        \forall x \, \forall z \, \big[ z \in P \implies x \notin z \big]
    \end{equation}
    Suppose that \( P \neq \emptyset \). Thus, there exists a \( z_{0} \) such that \( z_{0} \in P \).
    Applying \aref{rule-of-inference:universal-instantiation}{Universal Instantiation} on \eqref{eq:the-union-of-P-is-empty} yields,
    \begin{equation}
        \label{eq:the-union-of-P-is-empty-instantiated}
        \forall x \, \big[ z_{0} \in P \implies x \notin z_{0} \big]
    \end{equation}
    Let \( x_{0} \) be an arbitrary object. By \aref{rule-of-inference:universal-instantiation}{Universal Instantiation} on \eqref{eq:the-union-of-P-is-empty-instantiated},
    we have,
    \begin{equation*}
        z_{0} \in P \implies x_{0} \notin z_{0}
    \end{equation*}
    Since \( z_{0} \in P \), we conclude by \aref{rule-of-inference:modus-ponens}{Modus Ponens} that \( x_{0} \notin z_{0} \). Since \( x_{0} \) is arbitrary,
    it follows that \( \forall x \, \big( x \notin z_{0} \big) \) so that \( z_{0} = \emptyset \), which is a contradiction
    since \( P \) is a partition and no element of a partition can be empty by \cref{dfn:partition-of-a-set}. Clearly, the
    supposition \( P \neq \emptyset \) leads to a contradiction, so it must be false by \aref{rule-of-inference:reductio-ad-absurdum}{Reductio ad Absurdum}.
    \begin{equation*}
        P \in \partitions{\emptyset} \implies P = \emptyset
    \end{equation*}
    To prove the converse, we must prove that \( \emptyset \in \partitions{\emptyset} \).
    Obviously, if \( P = \emptyset \), then the statements \( \forall x \, \big( x \in P \implies x \neq \emptyset \big) \)
    and \( \forall x\,{\in}\,P \, \forall y\,{\in}\,P \, \big( x \neq y \implies x \intersect y = \emptyset \big) \)
    are both vacuously true. Also, since \( P = \emptyset \), the statement \( \exists z\,{\in}\,P \, \big( x \in z \big) \)
    is necessarily false, which means that \( \Union P = \emptyset \). In other words, \( \emptyset \) satisfies all the conditions
    stated in \cref{dfn:partition-of-a-set}. Thus, \( \emptyset \in \partitions{\emptyset} \) as required.
    \begin{equation*}
        P = \emptyset \implies P \in \partitions{\emptyset}
    \end{equation*}
    This establishes the biconditional \( P \in \partitions{\emptyset} \iff P = \emptyset \). And since \( P \) is arbitrary,
    we can generalize,
    \begin{equation*}
        \forall x \, \big[ x \in \partitions{\emptyset} \iff x = \emptyset \big]
    \end{equation*}
    By \cref{dfn:unordered-pair}, it follows that \( \partitions{\emptyset} = \mleft\{ \emptyset \mright\} \).
\end{proof}

\begin{rem}
    From \cref{thm:existence-of-the-set-of-all-partitions}, we know that \( \partitions{A} \subseteq \powerset{\powerset{A}} \).
    Therefore, \( \partitions{\emptyset} \subseteq \powerset{\powerset{\emptyset}} \). Since \( \powerset{\emptyset} = \mleft\{ \emptyset \mright\} \),
    it follows that \( \powerset{\powerset{\emptyset}} = \mleft\{ \emptyset, \mleft\{ \emptyset \mright\} \mright\} \). This is in agreeement with
    the fact that \( \partitions{\emptyset} = \mleft\{ \emptyset \mright\} \).
\end{rem}

\begin{thm}
    \label{thm:existence-of-the-set-of-all-equivalence-relations}
    For any set \( A \), \( \equivrels{A} \) exists.
\end{thm}

\begin{proof}
    By \cref{dfn:set-of-all-partitions-of-a-set-and-the-set-of-all-equivalence-relations-on-a-set},
    it follows that \( \forall x \, \big[ x \in \equivrels{A} \implies x \subseteq A^{2} \big] \).
    Therefore, \( \forall x \, \big[ x \in \equivrels{A} \implies x \in \powerset{A^{2}} \big] \).
    \cref{thm:sufficient-condition-for-unrestricted-comprehension} then guarantees the existence
    of \( \equivrels{A} \).
\end{proof}

\begin{thm}
    \label{thm:the-only-equivalence-relation-on-the-empty-set-is-the-empty-set}
    \( \equivrels{\emptyset} = \mleft\{ \emptyset \mright\} \)
\end{thm}

\begin{proof}
    Let \( E \in \equivrels{\emptyset} \) be arbitrary.
    By \cref{dfn:set-of-all-partitions-of-a-set-and-the-set-of-all-equivalence-relations-on-a-set,dfn:equivalence-relation},
    it follows that \( E \subseteq \emptyset^{2} \).
    Since \( 2 \neq \emptyset \), we have by \cref{thm:the-exponentiation-of-the-empty-set} that \( \emptyset^{2} = \emptyset \).
    Since \( E \subseteq \emptyset^{2} \) and \( \emptyset^{2} = \emptyset \),
    it follows that \( E \subseteq \emptyset \) and thus \( E = \emptyset \) by
    \cref{thm:only-the-empty-set-can-be-a-subset-of-the-empty-set}.
    \begin{equation*}
        E \in \equivrels{\emptyset} \implies E = \emptyset
    \end{equation*}
    Conversely, suppose that \( E = \emptyset \). Clearly, \( E \subseteq \emptyset^{2} \) because the empty set is a subset
    of any set by \cref{thm:empty-set-universal-subset}. Moreover, the following statements are all vacuously true:
    \begin{gather*}
        \forall x\,{\in}\,\emptyset \, \big[ \tuple{x,x} \in E \big]\\
        \forall x\,{\in}\,\emptyset \, \forall y\,{\in}\,\emptyset \, \big[ \tuple{x,y} \in E \implies \tuple{y,x} \in E \big]\\
        \forall x\,{\in}\,\emptyset \, \forall y\,{\in}\,\emptyset \, \forall z\,{\in}\,\emptyset \,
        \big[ \tuple{x,y} \in E \land \tuple{y,z} \in E \implies \tuple{x,z} \in E \big]
    \end{gather*}
    Thus, \( E \) must be an equivalence relation on \( \emptyset \) by \cref{dfn:equivalence-relation}, so \( E \in \equivrels{\emptyset} \).
    \begin{equation*}
        E = \emptyset \implies E \in \equivrels{\emptyset}
    \end{equation*}
    This establishes the biconditional \( E \in \equivrels{\emptyset} \iff E = \emptyset \) for any \( E \).
    Thus, \( \equivrels{\emptyset} = \mleft\{ \emptyset \mright\} \).
\end{proof}

\begin{cor}
    \label{thm:the-quotient-set-of-empty-set-with-respect-to-empty-set-is-the-empty-set}
    \( \emptyset / \emptyset = \emptyset \).
\end{cor}

\begin{proof}
    \( \emptyset \) is the only equivalence relation on \( \emptyset \)
    by \cref{thm:the-only-equivalence-relation-on-the-empty-set-is-the-empty-set}.
    By \cref{dfn:quotient-set},
    \( \emptyset / \emptyset =
    \setbuilderbig{x}{\exists a \, ( a \in \emptyset \land x = \equivclass{a}_{\emptyset} )}
    = \big\{ \, \big\} = \emptyset \).
\end{proof}

\begin{thm}
    \label{thm:equivalent-classes-are-pairwise-disjoint}
    Let \( A \) be a set and \( E \) be an equivalence relation on \( A \).
    Let \( a \) and \( b \) be arbitrary elements of \( A \). Then the following
    statements hold:
    \begin{enumerate}[topsep=2pt,noitemsep,label=\texttt{\textup{(\alph*)}}]
        \item \( \tuple{a,b} \in E \iff \equivclass{a}_{E} = \equivclass{b}_{E} \) \label{enum:condition-for-ab-in-E}
        \item \( \tuple{a,b} \notin E \iff \equivclass{a}_{E} \intersect \equivclass{b}_{E} = \emptyset \) \label{enum:condition-for-ab-not-in-E}
    \end{enumerate}
\end{thm}

\begin{proof}
    Since \( a,b \in A\), the equivalence classes \( \equivclass{a}_{E} \) and \( \equivclass{b}_{E} \) exist by
    \cref{thm:existence-of-equivalence-classes}. Suppose that \( \tuple{a,b} \in E \). By \cref{dfn:equivalence-class},
    we have,
    \begin{gather}
        \forall x \, \big( x \in \equivclass{a}_{E} \iff \tuple{x,a} \in E \big) \label{eq:predicate-for-equivalent-class-of-a} \\
        \forall x \, \big( x \in \equivclass{b}_{E} \iff \tuple{x,b} \in E \big) \label{eq:predicate-for-equivalent-class-of-b}
    \end{gather}
    Let \( x_{0} \in A \) be arbitrary. If \( \tuple{x_{0}, a} \in E \), then \( \tuple{x_{0}, b} \in E \) follows from the fact that
    \( \tuple{a,b} \in E \) and the fact that \( E \) is transitive.
    \begin{equation*}
        \tuple{x_{0}, a} \in E \implies \tuple{x_{0}, b} \in E
    \end{equation*}
    Conversely, if \( \tuple{x_{0},b} \in E \), then \( \tuple{x_{0}, a} \in E \) must also be true because \( \tuple{b,a} \in E \) and
    \( E \) is transitive.
    \begin{equation*}
        \tuple{x_{0},b} \in E \implies \tuple{x_{0},a} \in E
    \end{equation*}
    By \aref{rule-of-inference:universal-generalization}{Universal Generalization},
    \begin{equation}
        \label{eq:xa-in-E-iff-xb-in-E}
        \forall x \, \big[ \tuple{x,a} \in E \iff \tuple{x,b} \in E \big]
    \end{equation}
    Applying \eqref{eq:predicate-for-equivalent-class-of-a} and \eqref{eq:predicate-for-equivalent-class-of-b}
    to \eqref{eq:xa-in-E-iff-xb-in-E} yields,
    \begin{equation*}
        \forall x \, \big( x \in \equivclass{a}_{E} \iff x \in \equivclass{b}_{E} \big)
    \end{equation*}
    By the \nameref{axm:axiom-of-extensionality}, it follows that \( \equivclass{a}_{E} = \equivclass{b}_{E} \).
    And by the \aref{rule-of-inference:deduction-theorem}{Deduction Theorem},
    \begin{equation}
        \label{eq:ab-in-E-implies-equal-equivalent-classes}
        \tuple{a,b} \in E \implies \equivclass{a}_{E} = \equivclass{b}_{E}
    \end{equation}
    To prove the converse, suppose that \( \equivclass{a}_{E} = \equivclass{b}_{E} \) for some arbitrary \( a,b \in A \).
    Since \( E \) is reflexive, we know that \( \tuple{a,a} \in E \) so that \( a \in \equivclass{a}_{E} \).
    But \( \equivclass{a}_{E} = \equivclass{b}_{E} \), which means that \( a \in \equivclass{b}_{E} \) must also be true.
    Since \( a \in \equivclass{b}_{E} \), it follows from \cref{dfn:equivalence-class} that \( \tuple{a,b} \in E \). Thus,
    \begin{equation}
        \label{eq:equal-equivalent-classes-implies-ab-in-E}
        \equivclass{a}_{E} = \equivclass{b}_{E} \implies \tuple{a,b} \in E
    \end{equation}
    Together, \eqref{eq:ab-in-E-implies-equal-equivalent-classes} and \eqref{eq:equal-equivalent-classes-implies-ab-in-E}
    establish the biconditional \( \tuple{a,b} \in E \iff \equivclass{a}_{E} = \equivclass{b}_{E} \). With both \( a \) and
    \( b \) being arbitrary, \ref{enum:condition-for-ab-in-E} follows from \aref{rule-of-inference:universal-generalization}{Universal Generalization}.

    \vspace{0.3cm}
    \noindent
    To prove \ref{enum:condition-for-ab-not-in-E}, suppose that \( \tuple{a,b} \in E \) for some
    arbitrary \( a,b \in A \). By \ref{enum:condition-for-ab-in-E}, it follows that \( \equivclass{a}_{E} = \equivclass{b}_{E} \),
    and thus \( \equivclass{a}_{E} \intersect \equivclass{b}_{E} = \equivclass{a}_{E} \intersect \equivclass{a}_{E} = \equivclass{a}_{E} \)
    by \cref{set-identities:idempotence-of-intersections}. Since \( E \) is reflexive, we know that \( \tuple{a,a} \in E \) and thus
    \( a \in \equivclass{a}_{E} \). In other words, \( \equivclass{a}_{E} \neq \emptyset \).
    Since \( \equivclass{a}_{E} = \equivclass{b}_{E} = \equivclass{a}_{E} \),
    it follows that \( \equivclass{a}_{E} \intersect \equivclass{b}_{E} \neq \emptyset \).
    Thus, by the \aref{rule-of-inference:deduction-theorem}{Deduction Theorem}, we conclude,
    \begin{equation*}
        \tuple{a,b} \in E \implies \equivclass{a}_{E} \intersect \equivclass{b}_{E} \neq \emptyset
    \end{equation*}
    By \aref{rule-of-inference:contraposition}{Contraposition}, this is equivalent to,
    \begin{equation}
        \label{eq:disjoint-equivalent-classes-implies-ab-notin-E}
        \equivclass{a}_{E} \intersect \equivclass{b}_{E} = \emptyset \implies \tuple{a,b} \notin E
    \end{equation}
    Conversely, suppose that \( \equivclass{a}_{E} \intersect \equivclass{b}_{E} \neq \emptyset \) for some \( a,b \in A \).
    Thus, there exists some \( \alpha \in A \) such that \( \alpha \in \equivclass{a}_{E} \) and \( \alpha \in \equivclass{b}_{E} \).
    In other words, \( \tuple{\alpha,a} \in E \) and \( \tuple{\alpha,b} \in E \). By the reflexivity of \( E \), it follows that
    \( \tuple{a, \alpha} \in E \). Thus, \( \tuple{a,b} \in E \) must also be true by the transitivity of \( E \).
    By the \aref{rule-of-inference:deduction-theorem}{Deduction Theorem}, we conclude that,
    \begin{equation*}
        \equivclass{a}_{E} \intersect \equivclass{b}_{E} \neq \emptyset \implies \tuple{a,b} \in E
    \end{equation*}
    Again, by \aref{rule-of-inference:contraposition}{Contraposition}, this is equivalent to,
    \begin{equation}
        \label{eq:ab-not-in-E-implies-disjoint-equivalent-classes}
        \tuple{a,b} \notin E \implies \equivclass{a}_{E} \intersect \equivclass{b}_{E} = \emptyset
    \end{equation}
    Together, \eqref{eq:disjoint-equivalent-classes-implies-ab-notin-E} and \eqref{eq:ab-not-in-E-implies-disjoint-equivalent-classes}
    establish \ref{enum:condition-for-ab-not-in-E}.
\end{proof}

\begin{lem}
    \label{lem:the-quotient-set-of-a-set-is-a-partition-of-the-set}
    For any set \( A \) and any equivalence relation \( E \) on \( A \),
    \( A / E \) is a partition of \( A \).
    \begin{equation*}
        \forall A \, \forall E \, \Big[ E \in \equivrels{A} \implies A / E \in \partitions{A} \Big]
    \end{equation*}
\end{lem}

\begin{proof}
    Let \( A \) be any set. By the \aref{rule-of-inference:law-of-excluded-middle}{Law of Excluded Middle},
    \begin{equation*}
        A = \emptyset \lor A \neq \emptyset
    \end{equation*}

    \noindent
    {\scshape Case 1}: \( A = \emptyset \). If \( A = \emptyset \), then \( E \in \equivrels{A} \iff E = \emptyset \)
    by \cref{thm:the-only-equivalence-relation-on-the-empty-set-is-the-empty-set}. If \( A = \emptyset \) and \( E = \emptyset \),
    then \( A / E = \emptyset \) by \cref{thm:the-quotient-set-of-empty-set-with-respect-to-empty-set-is-the-empty-set}. Since
    \( \partitions{\emptyset} = \mleft\{ \emptyset \mright\} \) by \cref{thm:the-only-partition-for-the-empty-set-is-the-empty-set},
    necessarily \( A / E \in \partitions{A} \).
    \begin{equation*}
        \therefore \, \forall E \, \Big[ E \in \equivrels{A} \implies A / E \in \partitions{A} \Big]
    \end{equation*}

    \noindent
    {\scshape Case 2}: \( A \neq \emptyset \). Let \( E \in \equivrels{A} \) be arbitrary. Thus, \( E \subseteq A^{2}\) and \( E \)
    is reflexive, symmetric, and transitive. Since \( E \) is reflexive and \( A \neq \emptyset \), necessarily \( E \neq \emptyset \).
    Clearly, \( A / E\) exists by \cref{thm:existence-of-the-quotient-set}.
    And to prove that \( A / E \in \partitions{A} \), it is necessary to show that \( A / E \) satisfies the conditions stated in
    \cref{dfn:partition-of-a-set}.

    \vspace{0.3cm}
    \noindent
    The first condition is easy to prove because it follows immediately from the reflexivity of \( E \); i.e., since \( \tuple{a,a} \in E \)
    for any \( a \in A \) and \( A \neq \emptyset \), it follows that every equivalence class has at least one element.

    \vspace{0.3cm}
    \noindent
    To prove the second condition, suppose that \( x \) and \( y \) are two distinct elements of \( A / E \). That is, \( x = \equivclass{a}_{E} \)
    and \( y = \equivclass{b}_{E} \) for some \( a \in A \) and \( b \in A \). Since \( x \neq y \), \( x \intersect y = \emptyset \) follows from
    \cref{thm:equivalent-classes-are-pairwise-disjoint}, establishing the second condition.

    \vspace{0.3cm}
    \noindent
    For the third condition, we apply the definition for the union of sets. Thus, for any arbitrary \( x_{0} \),
    \begin{equation*}
        x_{0} \in \Union A / E
        \iff
        \exists z \, \Big[ z \in A / E \land x_{0} \in z \Big]
    \end{equation*}
    Using \cref{dfn:quotient-set}, the statement can then be expressed equivalently as,
    \begin{equation*}
        x_{0} \in \Union A / E
        \iff
        \exists z \, \Big[ \exists a \, \big( a \in A \land z = \equivclass{a}_{E} \big) \land x_{0} \in z \Big]
    \end{equation*}
    With \cref{tautology:grouping-conjuncts-into-existentially-quantified-formula,tautology:using-an-existentially-quantified-formula-to-express-the-fact-that-an-object-satisfies-a-formula},
    we can further simplify the formula like so,
    \begin{align}
        x_{0} \in \Union A / E &\iff \exists z \, \exists a \, \Big[ a \in A \land z = \equivclass{a}_{E} \land x_{0} \in z \Big] \notag \\
        &\iff \exists a \, \Big[ a \in A \land \exists z \, \big( z = \equivclass{a}_{E} \land x_{0} \in z \big) \Big] \notag \\
        &\iff \exists a \, \Big[ a \in A \land x_{0} \in \equivclass{a}_{E} \Big] \label{eq:x0-belongs-to-some-equivalence-class}
    \end{align}
    If \( x_{0} \in A \), then obviously \( x_{0} \in \equivclass{x_{0}}_{E} \land x_{0} \in A \), so necessarily \( x_{0} \in \Union A/E \) by
    \eqref{eq:x0-belongs-to-some-equivalence-class}. Conversely, if \( x_{0} \in \Union A / E \), then \( x_{0} \in \equivclass{a}_{E} \) for
    some \( a \in A \) by \eqref{eq:x0-belongs-to-some-equivalence-class}. Since \( E \subseteq A^{2} \) and \( \tuple{x_{0}, a} \in E \),
    necessarily \( x_{0} \in A \). In other words, \( x_{0} \in \Union A / E \iff x_{0} \in A \) for any \( x_{0} \). Thus, \( \Union A / E = A \)
    as required for the third condition to hold.

    \vspace{0.3cm}
    \noindent
    Having proved that \( A / E \) satisfies all three conditions in \cref{dfn:partition-of-a-set}, we conclude that \( A / E \in \partitions{A} \).
    \begin{equation*}
        \therefore \, \forall E \, \Big[ E \in \equivrels{A} \implies A / E \in \partitions{A} \Big]
    \end{equation*}
    Since it holds for all possible cases, the statement \( \forall E \, \big[ E \in \equivrels{A} \implies A / E \in \partitions{A} \big] \) follows from Disjuntion Elimination.
    With \( A \) being arbitrary, the lemma follows from \aref{rule-of-inference:universal-generalization}{Universal Generalization}.
\end{proof}

\begin{lem}
    \label{lem:any-partition-of-a-set-defines-an-equivalece-relation-on-the-set}
    If \( P \) is a partition of \( A \), then the set,
    \begin{equation*}
        E = \setbuilderbig{\tuple{x,y}}{\exists z\,{\in}\,P \, \mleft( x \in z \land y \in z \mright)}
    \end{equation*}
    is an equivalence relation on \( A \). Moreover, \( P = A / E \).
\end{lem}

\begin{proof}
    We start by proving that \( E \subseteq A^{2} \). Let \( x_{0} \in E \) be arbitrary. Thus,
    \( x_{0} = \tuple{a,b} \) where \( a \in z_{0} \) and \( b \in z_{0} \) for some \( z_{0} \in P \).
    By \cref{dfn:union-of-sets}, it is clear that \( a \in \Union P \) and \( b \in \Union P \). But
    \( P \) is a partition of \( A \), so by \cref{dfn:partition-of-a-set}, necessarily \( \Union P = A \).
    Therefore, \( a \in A \) and \( b \in A \) and \( x_{0} = \tuple{a,b} \in A^{2} \). In other words,
    \( x_{0} \in E \implies x_{0} \in A^{2} \) for any \( x_{0} \), which implies that \( E \subseteq A^{2} \).
    The next step is to establish that \( E \) is reflexive, symmetric, and transitive.

    \vspace{0.3cm}
    \noindent
    Reflexivity: Let \( a \in A \) be arbitrary. Since \( \Union P = A \), necessarily \( a \in \Union P \).
    Thus, there must exist a \( z \in P \) such that \( a \in z \). Since \( x \in z \land x \in z \) is
    trivially true, it follows that \( \exists z\,{\in}\,P \, \big[ x \in z \land x \in z \big] \) so that
    \( \tuple{a,a} \in E \) by definition. Thus, \( E \) is reflexive.

    \vspace{0.3cm}
    \noindent
    Symmetric: Let \( a \) and \( b \) be arbitrary elements of \( A \) such that \( \tuple{a,b} \in E \).
    By definition, there must exist some \( z \in P \) such that \( a \in z \) and \( b \in z \).
    But if \( a \in z \land b \in z \) is true, then \( b \in z \land a \in z \) is also necessarily true.
    In other words, \( \exists z\,{\in}\,P \, \big[ b \in z \land a \in z \big] \), and \( \tuple{b,a} \in E \)
    follows by definition. Thus, \( E \) is symmetric.

    \vspace{0.3cm}
    \noindent
    Transitive: Let \( a \), \( b \), and \( c \) be arbitrary members of \( A \) such that \( \tuple{a,b} \in E
    \land \tuple{b,c} \in E \). By the definition of \( E \), it follows that \( a \in z_{1} \land b \in z_{1}\)
    and \( b \in z_{2} \land c \in z_{2} \) for some \( z_{1} \in P \) and \( z_{2} \in P \). Since \( b \in z_{1}\)
    and \( b \in z_{2} \) are both true, it follows that \( b \in z_{1} \intersect z_{2} \), and thus
    \( z_{1} \intersect z_{2} \neq \emptyset \). But \( P \) is a partition of \( A \), so
    \( z_{1} \intersect z_{2} \neq \emptyset \) implies \( z_{1} = z_{2} \) by \cref{dfn:partition-of-a-set}.
    Since \( z_{1} = z_{2} \), we have \( a \in z_{1} \land c \in z_{1} \) for some \( z_{1} \in P \) so that
    \( \tuple{a,c} \in E \) by the definition of \( E \). Thus, \( E \) is transitive.

    \vspace{0.3cm}
    \noindent
    We have established that \( E \subseteq A^{2} \) and that it is reflexive, symmetric, and transitive.
    By \cref{dfn:equivalence-relation}, \( E \) is an equivalence relation on \( A \). And since \( E \)
    is an equivalence relation, it follows from \cref{thm:existence-of-the-quotient-set} that \( A / E \)
    exists. The final step, then, is to show that \( P = A / E \).

    \vspace{0.3cm}
    \noindent
    Let \( x_{0} \) be an arbitrary member of \( A / E \). In other words, \( x_{0} = \equivclass{a}_{E} \)
    for some \( a \in A \). We know that \( E \) is reflexive, so clearly \( a \in x_{0} \) and \( \tuple{a,a} \in E \).
    By the definition of \( E \), there must exists some \( z_{0} \in P \) such that \( a \in z_{0} \). We will prove that
    \( x_{0} \) and \( z_{0} \) are in fact the same set.

    \vspace{0.3cm}
    \noindent
    Let \( y_{0} \in x_{0} \) be arbitrary. Again by \cref{dfn:equivalence-class}, it follows
    that \( \tuple{y_{0},a} \in E \) and thus \( y_{0} \in z_{1} \) and \( a \in z_{1} \) for some \( z_{1} \in P \).
    But \( P \) is a partition, so \( a \in z_{0} \land a \in z_{1} \) implies that \( z_{0} = z_{1} \).
    Thus, \( y_{0} \) must also belong to \( z_{0} \). In other words, \( \forall y \, \big( y \in x_{0} \implies y \in z_{0} \big) \)
    and thus \( x_{0} \subseteq z_{0} \). Conversely, if \( y_{0} \) is an arbitrary element of \( z_{0} \),
    then \( \exists z\,{\in}\,P \, \mleft( y_{0} \in z \land a \in z \mright) \) because \( a \in z_{0} \).
    This means that \( \tuple{y_{0},a} \in E \), so \( y_{0} \in x_{0} = \equivclass{a}_{E} \) by definition.
    In other words, \( \forall y \, \big( y \in z_{0} \implies y \in x_{0} \big) \), so \( z_{0} \subseteq x_{0} \).

    \vspace{0.3cm}
    \noindent
    Having established both \( x_{0} \subseteq z_{0} \) and \( z_{0} \subseteq x_{0} \), we conclude that \( x_{0} = z_{0} \).
    Since \( z_{0} \in P \), it follows that \( x_{0} \in P \). By the \aref{rule-of-inference:deduction-theorem}{Deduction Theorem} and \aref{rule-of-inference:universal-generalization}{Universal Generalization}, we have,
    \begin{equation}
        \label{eq:quotient-set-is-a-subset-of-the-partition}
        \forall x \, \big[ x \in A / E \implies x \in P \big]
    \end{equation}
    To prove the converse, we let \( x_{0} \) be an arbitrary element of \( P \). Since \( P \) is a partition of \( A \)
    and \( x_{0} \in P \), it follows from \cref{dfn:partition-of-a-set} that \( x_{0} \neq \emptyset \). Thus, there must
    exists some \( a \) such that \( a \in x_{0} \). Since \( \Union P = A \), it follows \( a \in A \). We now proceed
    to show that \( x_{0} = \equivclass{a}_{E} \). Suppose that \( y_{0} \) is some arbitrary element of \( x_{0} \). Given
    that \( a \in x_{0} \) and \( y_{0} \in x_{0} \) with \( x_{0} \in P \),
    it follows that \( \exists z\,{\in}\,P \, \big( y_{0} \in z \land a \in z \big) \) and thus \( \tuple{y_{0},a} \in E \).
    By \cref{dfn:equivalence-class}, this means that \( y_{0} \in \equivclass{a}_{E} \).
    Therefore, \( \forall y \, \big( y \in x_{0} {\implies} y \in \equivclass{a}_{E} \big) \) and \( x_{0} \subseteq \equivclass{a}_{E} \).

    \vspace{0.3cm}
    \noindent
    Conversely, if \( y_{0} \in \equivclass{a}_{E} \), then necessarily \( \tuple{y_{0},a} \in E \). By the definition of \( E \), there must exist
    some \( z_{0} \in P \) such that \( y_{0} \in z_{0} \), \( a \in z_{0} \). But \( P \) is a partition on \( A \), so \( a \in x_{0} \land a \in z_{0} \)
    implies \( x_{0} = z_{0} \). Given that \( y_{0} \in z_{0} \), \( y_{0} \in x_{0} \) must also be true.
    Therefore, \( \forall y \, \big[ y \in \equivclass{a}_{E} {\implies} y \in x_{0} \big] \) and \( \equivclass{a}_{E} \subseteq x_{0} \).

    \vspace{0.3cm}
    \noindent
    Having established both \( x_{0} \subseteq \equivclass{a}_{E} \) and \( \equivclass{a}_{E} \subseteq x_{0} \), we can then conclude that
    \( x_{0} = \equivclass{a}_{E} \). Clearly, \( \equivclass{a}_{E} \in A / E \), so \( x_{0} \) must also be an element of \( A / E \). By the
    \aref{rule-of-inference:deduction-theorem}{Deduction Theorem} and \aref{rule-of-inference:universal-generalization}{Universal Generalization}, we have,
    \begin{equation}
        \label{eq:partition-is-a-subset-of-the-quotient-set}
        \forall x \, \Big[ x \in P \implies x \in A / E \Big]
    \end{equation}
    \( P = A / E \) follows from \eqref{eq:quotient-set-is-a-subset-of-the-partition}, \eqref{eq:partition-is-a-subset-of-the-quotient-set},
    and \nameref{axm:axiom-of-extensionality}.
\end{proof}

\begin{thm}[Fundamental Theorem of Equivalence Relations]
    \label{thm:fundamental-theorem-of-equivalence-relations}
    For any set \( A \), a bijection exists between \( \equivrels{A} \) and \( \partitions{A} \).
    \begin{equation*}
        \forall A \, \exists f \, \big[ \bijection{f}{\equivrels{A}}{\partitions{A}} \big]
    \end{equation*}
\end{thm}

\begin{proof}
    Let \( A \) be any set. Define the set \( f \) as follows:
    \begin{equation}
        \label{eq:define-bijection-f-between-eqv-A-and-par-A}
        f = \setbuilderBig{\orderedpair{E}{P}}{E \in \equivrels{A} \land P = A / E }
    \end{equation}
    We have to prove the following three claims:
    \begin{enumerate*}[label=(\alph*)]
        \item \( f \) is a function on \( \equivrels{A} \),
        \item \( \ran f = \partitions{A} \),
        \item \( f \) is injective.
    \end{enumerate*}

    \vspace{0.5cm}
    \noindent
    {\scshape Claim} (a): \( f \) is a function and \( \dom f = \equivrels{A} \).

    \vspace{0.2cm}
    \noindent
    Let \( E \), \( P_{1} \), and \( P_{2} \) be arbitrary objects such that
    \( \orderedpair{E}{P_{1}} \in f \) and \( \orderedpair{E}{P_{2}} \in f \). By \eqref{eq:define-bijection-f-between-eqv-A-and-par-A},
    it follows that \( E \in \equivrels{A} \), \( P_{1} = A / E \), and \( P_{2} = A / E \). Clearly, \( P_{1} = P_{2} \),
    so \( f \) is necessarily a function. The fact that \( \dom f \subseteq \equivrels{A} \) is already implied by
    \eqref{eq:define-bijection-f-between-eqv-A-and-par-A}. Thus, it remains to prove the inclusion in the opposite direction.
    Let \( E \in \equivrels{A} \) be arbitrary. Clearly, \( A / E \) exists by \cref{thm:existence-of-the-quotient-set}, and
    it follows from \eqref{eq:define-bijection-f-between-eqv-A-and-par-A} that \( \orderedpair{E}{A / E} \in f \).
    By \cref{dfn:domain-of-a-set-of-ordered-pairs}, we have \( E \in \dom f \). In other words, \( \equivrels{A} \subseteq \dom f \).
    Having proven the inclusion in the opposite direction, we conclude by \cref{thm:set-equality-in-terms-of-subsets} that \( \dom f = \equivrels{A} \).

    \vspace{0.5cm}
    \noindent
    {\scshape Claim} (b): \( \ran f = \partitions{A} \).

    \vspace{0.2cm}
    \noindent
    Let \( P \in \partitions{A} \) be arbitrary. Since \( P \) is a partition of \( A \),
    it follows from \cref{lem:any-partition-of-a-set-defines-an-equivalece-relation-on-the-set}
    that the set,
    \begin{equation*}
        E = \setbuilderBig{\tuple{x,y}}{\exists z\,{\in}\,P \, \big( x \in z \land y \in z \big)}
    \end{equation*}
    is an equivalence relation on \( A \) and that \( A / E = P \). Applying these facts to \eqref{eq:define-bijection-f-between-eqv-A-and-par-A},
    we deduce that \( \orderedpair{E}{P} \in f \). Therefore, \( P \in \ran f \). In other words, \( \partitions{A} \subseteq \ran f \).
    To prove the inclusion in the other direction, we let \( P \in \ran f \) be arbitrary. By \eqref{eq:define-bijection-f-between-eqv-A-and-par-A}, it follows
    that \( P = A / E \) for some \( E \in \equivrels{A} \). Since \( E \) is an equivalence relation on \( A \), it follows from
    \cref{lem:the-quotient-set-of-a-set-is-a-partition-of-the-set} that \( P = A / E \) is a partition on \( A \) and thus \( P \in \partitions{A} \).
    In other words, \( \ran f \subseteq \partitions{A} \). With both \( \ran f \subseteq \partitions{A} \) and \( \partitions{A} \subseteq \ran f \)
    established, we can conclude that \( \ran f = \partitions{A} \).

    \vspace{0.5cm}
    \noindent
    {\scshape Claim} (c): \( f \) is injective.

    \vspace{0.2cm}
    \noindent
    Let \( E_{1} \), \( E_{2} \), and \( P \) be arbitrary objects such that \( \orderedpair{E_{1}}{P} \in f \)
    and \( \orderedpair{E_{2}}{P} \in f \). By \eqref{eq:define-bijection-f-between-eqv-A-and-par-A}, it follows
    that \( E_{1} \in \equivrels{A} \), \( E_{2} \in \equivrels{A} \), and \( A / E_{1} = A / E_{2} = P \).
    Since \( E_{1} \) and \( E_{2} \) are both equivalence relations on \( A \),
    \( P \) must be a partition on \( A \) by \cref{lem:the-quotient-set-of-a-set-is-a-partition-of-the-set}.
    By \cref{dfn:partition-of-a-set}, it follows that,
    \begin{equation}
        \label{eq:A-equals-to-the-union-of-P}
        \Union P = A
    \end{equation}
    Now let \( b \in A \) be arbitrary. By \eqref{eq:A-equals-to-the-union-of-P}, it follows that \( b \in \Union P \).
    Thus,
    \begin{equation*}
        \begin{aligned}
            b \in z_{0}&
        \end{aligned}
        \begin{aligned}
            &\text{for some \( z_{0} \in P \)}
        \end{aligned}
    \end{equation*}
    But \( P = A / E_{1} = A / E_{2} \), so necessarily \( z_{0} \in A / E_{1} \) and \( z_{0} \in A / E_{2} \). By \cref{dfn:quotient-set},
    it follows that \( z_{0} = \equivclass{\alpha}_{E_{1}} \) for some \( \alpha \in A \) and \( z_{0} = \equivclass{\beta}_{E_{2}} \) for
    some \( \beta \in A \).
    Since \( b \in z_{0} \), it follows that \( b \in \equivclass{\alpha}_{E_{1}} \) and
    \( b \in \equivclass{\beta}_{E_{2}} \). By \cref{thm:equivalent-classes-are-pairwise-disjoint}, necessarily
    \( \equivclass{b}_{E_{1}} = \equivclass{\alpha}_{E_{1}} = z_{0} = \equivclass{\beta}_{E_{2}} = \equivclass{b}_{E_{2}} \). Thus,
    \begin{equation*}
        \forall a\,{\in}\,A \, \Big[ \tuple{a,b} \in E_{1} \iff \tuple{a,b} \in E_{2} \Big]
    \end{equation*}
    And since \( b \) is arbitrary, we have by \aref{rule-of-inference:universal-generalization}{Universal Generalization},
    \begin{equation*}
        \forall a\,{\in}\,A \, \forall b\,{\in}\,A \, \Big[ \tuple{a,b} \in E_{1} \iff \tuple{a,b} \in E_{2} \Big]
    \end{equation*}
    Since \( E_{1} \) and \( E_{2} \) are both sets of 2-tuples, it is obvious that \( E_{1} = E_{2} \). With that,
    we have shown that \( \orderedpair{E_{1}}{P} \in f \land \orderedpair{E_{2}}{P} \in f \implies E_{1} = E_{2} \). Thus, \( f \)
    must be injective by \cref{dfn:injection}.

    \vspace{0.3cm}
    \noindent
    Since \( f \) is both injective and surjective with \( \dom f = \equivrels{A} \) and \( \ran f = \partitions{A} \), it
    follows that,
    \begin{equation*}
        \bijection{f}{\equivrels{A}}{\partitions{A}} \qedhere
    \end{equation*}
\end{proof}

\subsubsection{Orderings} \label{ssec:order-relations}
Just like how equivalence relations generalize the concept of equality, the notion of ordering
(dictionary ordering of words, the ordering of numbers by magnitude, etc.) can also be generalized and defined in
abstract terms without any reference to the concrete objects being ordered.

\begin{dfn}[Partial Ordering]
    \label{dfn:partial-ordering}
    A binary relation \( \leq \,\, \subseteq A^{2} \) is a \keyword{partial ordering} on \( A \)
    if and only if it is reflexive, anti-symmetric, and transitive. Partial orderings are also called
    \keyword{weak} or \keyword{non-strict partial orderings}.
\end{dfn}

\begin{dfn}[Strict Partial Ordering]
    \label{dfn:strict-partial-ordering}
    A binary relation \( < \,\, \subseteq A^{2} \) is a \keyword{strict partial ordering}
    on \( A \) if and only if it is asymmetric, and transitive. Strict partial
    orderings are also called \keyword{strong partial orderings}.
\end{dfn}

\begin{rem}
    It follows that any strict partial ordering must be irreflexive due to its asymmetry.
\end{rem}

\begin{lem}
    \label{lem:any-partial-ordering-corresponds-to-a-strict-partial-ordering}
    If \( \leq \) is a partial ordering on \( A \), then the set,
    \begin{equation*}
        S = \setbuilderbig{\tuple{x,y}}{x \leq y \land x \neq y}
    \end{equation*}
    is a strict ordering on \( A \).
\end{lem}

\begin{proof}
    Let \( \leq \) be a partial ordering on \( A \). Let \( S = \setbuilderbig{\tuple{x,y}}{x \leq y \land x \neq y} \).
    We have to prove that the binary relation \( S \) is asymmetric, and transitive.

    \vspace{0.3cm}
    \noindent
    Asymmetry: Let \( x_{0} \in A \) and \( y_{0} \in A \) be arbitrary. Suppose that \( \tuple{x_{0}, y_{0}} \in S \). Thus,
    \( x_{0} \leq y_{0} \land x_{0} \neq y_{0} \) by definition. If \( \tuple{y_{0},x_{0}} \in S \) were true, then necessarily
    \( y_{0} \leq x_{0} \land y_{0} \neq x_{0} \). This means that \( x_{0} \leq y_{0} \) and \( y_{0} \leq x_{0} \) are both
    true, which implies \( x_{0} = y_{0} \) due to the anti-symmetric property of \( \leq \).
    In other words, \( x_{0} = y_{0} \) and \( x_{0} \neq y_{0} \) are both true, which is an impossibility.
    By \aref{rule-of-inference:reductio-ad-absurdum}{Reductio ad Absurdum}, we conclude that \( \tuple{y_{0}, x_{0}} \notin S \). And by the \aref{rule-of-inference:deduction-theorem}{Deduction Theorem},
    we have, \( \tuple{x_{0},y_{0}} \in S {\implies} \tuple{y_{0}, x_{0}} \notin S \). Since \( x_{0} \) and \( y_{0} \)
    are arbitrary, the asymmetry of \( S \) follows from \aref{rule-of-inference:universal-generalization}{Universal Generalization}.

    \vspace{0.3cm}
    \noindent
    Transitivity: Let \( x_{0} \), \( y_{0} \), and \( z_{0} \) be arbitrary elements of \( A \). Suppose that \( \tuple{x_{0},y_{0}} \in S \)
    and \( \tuple{y_{0},z_{0}} \in S \). Then \( x_{0} \leq y_{0} \land x_{0} \neq y_{0} \) and \( y_{0} \leq z_{0} \land y_{0} \neq z_{0} \)
    by definition. Since \( \leq \) is transitive, \( x_{0} \leq y_{0} \land y_{0} \leq z_{0} \) implies \( x_{0} \leq z_{0} \). Now, suppose
    that \( x_{0} = z_{0} \). But this would mean that \( x_{0} \leq y_{0} \) and \( y_{0} \leq x_{0} \) are both true. And since \( \leq \)
    is anti-symmetric, it follows that \( x_{0} = y_{0} \), contradicting the fact that \( x_{0} \neq y_{0} \). Thus, by \aref{rule-of-inference:reductio-ad-absurdum}{Reductio ad Absurdum},
    the supposition that \( x_{0} = z_{0} \) must be false. In other words, \( x_{0} \neq z_{0} \). Since \( x_{0} \leq z_{0} \) and
    \( x_{0} \neq z_{0} \), \( \tuple{x_{0},z_{0}} \in S \) by definition. Thus, we have by the \aref{rule-of-inference:deduction-theorem}{Deduction Theorem},
    \( \tuple{x_{0},y_{0}} \in S \land \tuple{y_{0},z_{0}} \in S {\implies} \tuple{x_{0},z_{0}} \in S \) for any \( x_{0}, y_{0}, z_{0} \in A \).
    The transitivity of \( S \) follows immediately from \aref{rule-of-inference:universal-generalization}{Universal Generalization}.

    \vspace{0.3cm}
    \noindent
    Since \( S \) is asymmetric, and transitive, it follows that \( S \) is a strict partial ordering on \( A \).
\end{proof}

\begin{lem}
    \label{lem:any-strict-partial-ordering-corresponds-to-a-partial-ordering}
    If \( < \) is a strict partial ordering on \( A \), then the set,
    \begin{equation*}
        R = \setbuilderbig{\tuple{x,y}}{x < y \lor x = y}
    \end{equation*}
    is a partial ordering on \( A \).
\end{lem}

\begin{proof}
    Let \( < \) be a strict partial ordering on \( A \). Let \( R = \setbuilderbig{\tuple{x,y}}{x < y \lor x = y} \).
    We have to prove that \( R \) is reflexive, anti-symmetric, and transitive.

    \vspace{0.3cm}
    \noindent
    Reflexivity: Let \( x_{0} \in A \) be arbitrary. Since \( x_{0} = x_{0} \) is trivially true, \( x_{0} < x_{0} \lor x_{0} = x_{0} \)
    must also be true. Therefore, \( \tuple{x_{0},x_{0}} \in R \) by definition. With \( x_{0} \) being arbitrary, we can conclude that
    \( \forall x\,{\in}\,A \, \big[ \tuple{x,x} \in R \big] \). \( R \) is reflexive.

    \vspace{0.3cm}
    \noindent
    Anti-symmetry: Let \( x_{0} \in A \) and \( y_{0} \in A \) such that \( \tuple{x_{0},y_{0}} \in R \land \tuple{y_{0},x_{0}} \in R \).
    Thus, \( x_{0} < y_{0} \lor x_{0} = y_{0} \) and \( y_{0} < x_{0} \lor y_{0} = x_{0} \) by definition. Suppose that \( x_{0} \neq y_{0} \).
    Then \( x_{0} < y_{0} \) and \( y_{0} < x_{0} \) would both be true by \aref{rule-of-inference:disjunctive-syllogism}{Disjunctive Syllogism}. But this is a contradiction because \( < \)
    is asymmetric where \( x_{0} < y_{0} \) implies \( \lnot \, y_{0} < x_{0} \). In other words, \( x_{0} < y_{0} \) and \( y_{0} < x_{0} \)
    can never both be true at the same time. Clearly, the supposition \( x_{0} \neq y_{0} \) leads to a contradiction. So by Reductio ad
    Absurdum, it follows that \( x_{0} = y_{0} \). And by the \aref{rule-of-inference:deduction-theorem}{Deduction Theorem}, we have \( x_{0} < y_{0} \land y_{0} < x_{0} {\implies} x_{0} = y_{0} \)
    for any \( x_{0} \in A \) and \( y_{0} \in A \). Thus, \( R \) is anti-symmetric.

    \vspace{0.3cm}
    \noindent
    Transitivity: Let \( x_{0} \in A \), \( y_{0} \in A \), and \( z_{0} \in A \) such that \( \tuple{x_{0},y_{0}} \in R \land \tuple{y_{0},z_{0}} \in R \).
    Thus, \( x_{0} < y_{0} \lor x_{0} = y_{0} \) and \( y_{0} < z_{0} \lor y_{0} = z_{0} \) by definition.
    Since \( x_{0} < y_{0} \) and \( x_{0} = y_{0} \) cannot simultaneously be true (\( < \) is irreflexive) we only need to consider two separate cases.

    \vspace{0.3cm}
    \noindent
    {\scshape Case 1}: \( x_{0} < y_{0} \) and \( x_{0} \neq y_{0} \). Clearly, either \( y_{0} < z_{0} \land y_{0} \neq z_{0} \) or
    \( y_{0} = z_{0} \land \lnot \, \mleft( y_{0} < z_{0} \mright) \) because \( < \) is irreflexive. If \( y_{0} < z_{0} \), then
    obviously \( x_{0} < z_{0} \) because \( < \) is transitive. If \( y_{0} = z_{0} \), then \( x_{0} < z_{0} \) must also be true
    because \( x_{0} < y_{0} \) and \( y_{0} = z_{0} \). Since \( x_{0} < z_{0} \) is implied in both possibilities, it is necessarily
    true by \aref{rule-of-inference:disjunction-elimination}{Disjunction Elimination}. Since \( x_{0} < z_{0} \), it follows that \( x_{0} < z_{0} \lor x_{0} = z_{0} \) is true. Thus,
    \( \tuple{x_{0},z_{0}} \in R \) by definition.

    \vspace{0.3cm}
    \noindent
    {\scshape Case 2}: \( x_{0} = y_{0} \) and \( \lnot \, \mleft( x_{0} < y_{0} \mright) \). Again, either \( y_{0} < z_{0} \land y_{0} \neq z_{0} \)
    or \( y_{0} = z_{0} \land \lnot \, \mleft( y_{0} < z_{0} \mright) \) and both possibilities are mutually exclusive. If \( y_{0} < z_{0} \), then
    \( x_{0} < z_{0} \) must be true because \( x_{0} = y_{0} \) and \( y_{0} < z_{0} \). On the other hand, if \( y_{0} = z_{0} \), then \( x_{0} = y_{0} = z_{0} \).
    By \aref{rule-of-inference:disjunction-elimination}{Disjunction Elimination}, we derive \( x_{0} < z_{0} \lor x_{0} = z_{0} \), which implies \( \tuple{x_{0},z_{0}} \in R \) by definition.

    \vspace{0.3cm}
    \noindent
    Thus, \( \tuple{x_{0},z_{0}} \in R \) is implied in all cases. By \aref{rule-of-inference:disjunction-elimination}{Disjunction Elimination} and the \aref{rule-of-inference:deduction-theorem}{Deduction Theorem}, we conclude that
    \( \tuple{x_{0},y_{0}} \in R \land \tuple{y_{0},z_{0}} {\implies} \tuple{x_{0},z_{0}} \) for any \( x_{0},y_{0},z_{0} \in A \). Thus,
    \( R \) is transitive.

    \vspace{0.3cm}
    \noindent
    Since \( R \) is reflexive, anti-symmetric, and transitive, it must be a partial ordering.
\end{proof}

\begin{thm}
    \label{thm:a-bijection-exists-between-all-partial-orderings-and-strict-partial-orderings}
    Let \( A \) be any set. Let \( P \) be the set of all partial orderings on \( A \) and \( S \)
    be the set of all strict partial orderings on \( A \). Then a bijection exists between \( P \)
    and \( S \).
\end{thm}

\begin{proof}
    Let the set \( f \) be defined as,
    \begin{equation}
        \label{eq:define-the-bijection-f-between-P-and-S}
        f = \setbuilderBig{\orderedpair{\rho}{\sigma}}{ \rho \in P \land \sigma = \rho - \identityrel{A} }
    \end{equation}
    The goal is to prove that:
    \begin{enumerate*}[label=(\alph*)]
        \item \( f \) is a function,
        \item \( \dom f = P \),
        \item \( \ran f = S \),
        \item \( f \) is injective.
    \end{enumerate*}

    \vspace{0.5cm}
    \noindent
    {\scshape Claim} (a): \( f \) is a function.

    \vspace{0.2cm}
    \noindent
    Let \( \rho \), \( \sigma_{1} \), and \( \sigma_{2} \) be arbitrary objects such that \( \orderedpair{\rho}{\sigma_{1}} \in f \)
    and \( \orderedpair{\rho}{\sigma_{2}} \in f \). By \eqref{eq:define-the-bijection-f-between-P-and-S}, it follows that \( \sigma_{1} = \rho - \identityrel{A} \)
    and \( \sigma_{2} = \rho - \identityrel{A} \). Clearly, \( \sigma_{1} = \sigma_{2} \). So \( f \) is indeed a function by \cref{dfn:function-basic}.

    \vspace{0.5cm}
    \noindent
    {\scshape Claim} (b): \( \dom f = P \).

    \vspace{0.2cm}
    \noindent
    Let \( \rho \in \dom f \) be arbitrary. By \cref{dfn:domain-of-a-set-of-ordered-pairs}, it follows that \( \orderedpair{\rho}{\sigma} \in f \)
    for some \( \sigma \). Therefore, \( \rho \in P \) by \eqref{eq:define-the-bijection-f-between-P-and-S}. In other words, \( \dom f \subseteq P \).
    Conversely, let \( \rho \in P \) be arbitrary. Clearly, \( \orderedpair{\rho}{\rho - \identityrel{A}} \in f \) by \eqref{eq:define-the-bijection-f-between-P-and-S},
    so \( \rho \in \dom f \) and thus \( P \subseteq \dom f \). Since \( \dom f \subseteq P \) and \( P \subseteq \dom f \) are both true, necessarily \( \dom f = P \)
    by \cref{thm:set-equality-in-terms-of-subsets}.

    \vspace{0.5cm}
    \noindent
    {\scshape Claim} (c): \( \ran f = S \)

    \vspace{0.2cm}
    \noindent
    Let \( \sigma \in \ran f \) be arbitrary. Thus, there exists some \( \rho \in P \) such that \( \sigma = \rho - \identityrel{A} \). Since \( P \) is a set of all
    partial orderings on \( A \) and \( \rho \in P \), it follows that \( \rho \) is a partial ordering on \( A \).
    \begin{align*}
        \sigma &= \rho - \setbuilderbig{\tuple{x,y}}{x \in A \land y \in A \land x \neq y}\\
               &= \setbuilderbig{\tuple{x,y}}{\tuple{x,y} \in \rho \land x \neq y}
    \end{align*}
    Since \( \rho \) is a partial ordering on \( A \), \( \sigma \) must be a strict partial ordering on \( A \)
    by \cref{lem:any-partial-ordering-corresponds-to-a-strict-partial-ordering}.
    Thus, \( \sigma \in S \). In other words, \( \sigma \in \ran f {\implies} \sigma \in S \) for any \( \sigma \). Thus, \( \ran f \subseteq S \). To prove the inclusion
    in the opposite direction, we let \( \sigma \in S \) be arbitrary. Since \( S \) is the set of all strict partial orderings on \( A \), it follows that
    \( \sigma \) is a strict partial ordering on \( A \). And by \cref{lem:any-strict-partial-ordering-corresponds-to-a-partial-ordering}, the set,
    \begin{equation*}
        \rho = \setbuilderbig{\tuple{x,y}}{\tuple{x,y} \in \sigma \lor x = y}
    \end{equation*}
    is a partial ordering on \( A \). Let \( x_{0} \) be an arbitrary object. Thus,
    \begin{align*}
        x_{0} \in \rho
        &\iff
        \exists a \, \exists b \, \big[ x_{0} = \tuple{a,b} \land \mleft( \tuple{a,b} \in \sigma \lor a = b \mright) \big]\\
        &\iff
        \exists a \, \exists b \, \big[ x_{0} = \tuple{a,b} \land \tuple{a,b} \in \sigma \lor x_{0} = \tuple{a,b} \land a = b \big]\\
        &\iff
        \exists a \, \exists b \, \big[ x_{0} = \tuple{a,b} \land \tuple{a,b} \in \sigma \big] \lor
        \exists a \, \exists b \, \big[ x_{0} = \tuple{a,b} \land a = b \big]\\
        &\iff
        x_{0} \in \sigma \lor x_{0} \in \identityrel{A}\\
        &\iff
        x_{0} \in \big( \sigma \union \identityrel{A} \big)
    \end{align*}
    By \aref{rule-of-inference:universal-generalization}{Universal Generalization} and the \nameref{axm:axiom-of-extensionality}, we conclude that \( \rho = \sigma \union \identityrel{A} \).
    Clearly, \( \rho - \identityrel{A} = \big( \sigma \union \identityrel{A} \big) - \identityrel{A} \), so by
    \cref{set-identities:right-distributivity-of-difference,set-identities:commutativity-of-union,set-identities:de-morgan-laws},
    \begin{align*}
        \rho - \identityrel{A}
        &= \mleft( \sigma - \identityrel{A} \mright) \union \mleft( \identityrel{A} - \identityrel{A} \mright)
        = \mleft( \sigma - \identityrel{A} \mright) \union \emptyset\\
        &= \emptyset \union \mleft( \sigma - \identityrel{A} \mright)
        = \mleft( \sigma - \sigma \mright) \union \mleft( \sigma - \identityrel{A} \mright)\\
        &= \sigma - \mleft( \sigma \intersect \identityrel{A} \mright)
    \end{align*}
    But \( \sigma \) is a strict partial ordering on \( A \), which means that \( \sigma \) is irreflexive. In other words, \( \sigma \intersect \identityrel{A} = \emptyset \).
    \begin{equation*}
        \therefore \, \rho - \identityrel{A} = \sigma - \emptyset = \sigma
    \end{equation*}
    Since \( \rho \in P \) and \( \rho - \identityrel{A} = \sigma \), clearly \( \orderedpair{\rho}{\sigma} \in f \)
    by \eqref{eq:define-the-bijection-f-between-P-and-S} so that \( \sigma \in \ran f \).
    In other words, \( \sigma \in S {\implies} \sigma \in \ran f \) for any \( \sigma \). Thus, \( S \subseteq \ran f \).
    Since \( \ran f \subseteq S \) and \( S \subseteq \ran f \) are both true, it follows that \( \ran f = S \).

    \vspace{0.5cm}
    \noindent
    {\scshape Claim} (d): \( f \) is injective.

    \vspace{0.2cm}
    \noindent
    Let \( \sigma_{1} \), \( \sigma_{2} \), and \( \rho \) be arbitrary objects such that
    \( \orderedpair{\sigma_{1}}{\rho} \in f \) and \( \orderedpair{\sigma_{2}}{\rho} \in f \).
    By \eqref{eq:define-the-bijection-f-between-P-and-S}, it follows that,
    \begin{gather*}
        \rho \in P \land \sigma_{1} = \rho - \identityrel{A}\\
        \rho \in P \land \sigma_{2} = \rho - \identityrel{A}
    \end{gather*}
    Clearly, \( \sigma_{1} = \sigma_{2} \). Thus, \( f \) must be injective by \cref{dfn:injection}.

    \vspace{0.3cm}
    \noindent
    We have shown that \( \bijection{f}{P}{S} \). Thus, a bijection exists between \( P \) and \( S \).
\end{proof}

\vspace{0.3cm}
\noindent
The use of the term ``partial'' in \cref{dfn:partial-ordering,dfn:strict-partial-ordering} merits further discussion.
In short, the term ``partial'' is used to emphasized that not every pair of elements from \( A \) needs to be related
by the \( \leq \) relation. For example, consider the set \( A = \mleft\{ a,b,c \mright\} \) and the binary relation
\( \leq \,\, = \mleft\{ \tuple{a,a}, \tuple{b,b}, \tuple{c,c}, \tuple{a,b} \mright\} \) on \( A \). \( \leq \) is obviously an
ordering on \( A \) because it is reflexive, anti-symmetric, and transitive, but notice that \( a \) and \( c \) aren't related to
each other by the \( \leq \) relation at all because \( \tuple{a,c} \in \, \leq \) and \( \tuple{c,a} \in \, \leq \) are both false.
That is why \( \leq \) is called a \emph{partial} ordering because not every pair of objects from \( A \) is ordered (related by \( \leq \)).
This leads us to the following definition.

\begin{dfn}[Total Ordering]
    \label{dfn:total-non-strict-ordering}
    A binary relation \( \leq \) is a \keyword{total} ordering on \( T \) if and only if:
    \begin{enumerate}[topsep=2pt,label=\texttt{\textup{(\alph*)}}]
        \item \( \leq \) is reflexive, anti-symmetric, and transitive.
        \item \( \forall x\,{\in}\,T \, \forall y\,{\in}\,T \, \big( x \leq y \lor y \leq x \big) \)
    \end{enumerate}
\end{dfn}

\noindent
A total ordering is also called a \keyword{linear} ordering because all the elements of \( A \)
can be arranged linearly from `smallest' to `largest' with respect to the \( \leq \) order relation.
This is possible because every pair of elements from \( A \) is related by \( \leq \). Note also that
totality and partiality are not contradictory concepts. Totality does not imply non-partiality and
partiality does not imply non-totality; rather, total orderings are a \emph{special case} of a partial
orderings where every pair of elements in the set happened to be related by the order relation.

\vspace{0.3cm}
\noindent
Totality or linearity can also be equivalently formulated in terms of strict orderings.
Suppose, for instance, that \( \leq \) is a non-strict ordering on some set \( A \).
\cref{thm:a-bijection-exists-between-all-partial-orderings-and-strict-partial-orderings,lem:any-strict-partial-ordering-corresponds-to-a-partial-ordering}
guarantee the existence of a unique binary relation \( < \) such that \( < \) is a strict ordering on \( A \) and that,
\begin{equation*}
    \forall x\,{\in}\,A \, \forall y\,{\in}\,A \, \Big[ x \leq y \iff \big( x < y \lor x = y \big) \Big]
\end{equation*}
The condition for linearity given in \cref{dfn:total-non-strict-ordering} is that,
\begin{equation*}
    \forall x\,{\in}\,A \, \forall y\,{\in}\,A \, \Big[ x \leq y \lor y \leq x \Big]
\end{equation*}
Using elementary logic, we can re-express that condition as,
\begin{align*}
    &\forall x\,{\in}\,A \, \forall y\,{\in}\,A \, \Big[ \mleft( x < y \lor x = y \mright) \lor \mleft( y < x \lor y = x \mright) \Big]\\
    &\iff
    \forall x\,{\in}\,A \, \forall y\,{\in}\,A \, \Big[ x = y \lor x < y \lor y < x \Big]
\end{align*}
Thus, we have the following definition.

\begin{dfn}[Total Strict Ordering]
    \label{dfn:total-strict-ordering}
    A binary relation \( < \) on \( A \) is a \keyword{total} strict ordering on \( T \) if and only if:
    \begin{enumerate}[topsep=2pt,label=\texttt{\textup{(\alph*)}}]
        \item \( < \) is asymmetric and transitive.
        \item \( \forall x\,{\in}\,T \, \forall y\,{\in}\,T \, \big( x = y \lor x < y \lor y < x \big) \)
    \end{enumerate}
\end{dfn}

\begin{dfn}[Poset]
    \label{dfn:poset}
    If the binary relation \( \leq \) is a partial ordering on \( P \), then the tuple \( \tuple{P, \leq} \) is called a \keyword{partially ordered set}
    (\keyword{poset} for short) and \( \tuple{P, <} \) is called a \keyword{strict poset}.
\end{dfn}

\begin{dfn}[Toset]
    \label{dfn:toset}
    If the binary relation \( \leq \) is a total ordering on \( T \), then the tuple \( \tuple{T, \leq} \) is called a \keyword{totally ordered set}
    (\keyword{toset} for short) and \( \tuple{T, <} \) is called a \keyword{strict toset}.
\end{dfn}

\begin{rem}
    When the context makes it clear, the term poset (toset) can be used for both \( \tuple{P, \leq} \) and \( \tuple{P, <} \).
    The symbolic representation for the order relation usually suffices to convey the strictness of the ordering
    (e.g,. \( < \) for strict orderings and \( \leq \) for non-strict orderings).
\end{rem}

\vspace{0.3cm}
\noindent
Other concepts can be defined in terms of order relations such as the basic notion of extremum.

\begin{dfn}[Least Element]
    \label{dfn:least-element}
    \hfill
    \begin{enumerate}[topsep=2pt,label=\texttt{\textup{(\alph*)}},itemsep=5pt]
        \item If \( \tuple{P, \leq} \) is a poset and \( A \subseteq P \), then \( b \) is the \keyword{least element} of \( A \)
              if and only if \( b \in A \land \forall x\,{\in}\,A \, \big[ b \leq x \big] \).
        \item If \( \tuple{P, <} \) is a poset and \( A \subseteq P \), then \( b \) is the \keyword{least element} of \( A \)
              if and only if \( b \in A \land \forall x\,{\in}\,A \, \big[ b < x \lor x = b \big] \).
    \end{enumerate}
\end{dfn}

\begin{dfn}[Greatest Element]
    \label{dfn:greatest-element}
    \hfill
    \begin{enumerate}[topsep=2pt,label=\texttt{\textup{(\alph*)}},itemsep=5pt]
        \item If \( \tuple{P, \leq} \) is a poset and \( A \subseteq P \), then \( u \) is the \keyword{greatest element} of \( A \)
              if and only if \( u \in A \land \forall x\,{\in}\,A \, \big[ x \leq u \big] \).
        \item If \( \tuple{P, <} \) is a poset and \( A \subseteq P \), then \( u \) is the \keyword{greatest element} of \( A \)
              if and only if \( u \in A \land \forall x\,{\in}\,A \, \big[ x < u \lor x = u \big] \).
    \end{enumerate}
\end{dfn}

\begin{dfn}[Minimal Element]
    \label{dfn:minimal-element}
    \hfill
    \begin{enumerate}[topsep=2pt,label=\texttt{\textup{(\alph*)}},itemsep=5pt]
        \item If \( \tuple{P, \leq} \) is a poset and \( A \subseteq P \), then \( b \) is a \keyword{minimal element} of \( A \)
              if and only if \( b \in A \land \lnot \, \exists x\,{\in}\,A \, \big[ x \leq b \land x \neq b \big] \).
        \item If \( \tuple{P, <} \) is a poset and \( A \subseteq P \), then \( b \) is a \keyword{minimal element} of \( A \)
              if and only if \( b \in A \land \lnot \, \exists x\,{\in}\,A \, \big[ x < b \big] \).
    \end{enumerate}
\end{dfn}

\begin{dfn}[Maximal Element]
    \label{dfn:maximal-element}
    \hfill
    \begin{enumerate}[topsep=2pt,label=\texttt{\textup{(\alph*)}},itemsep=5pt]
        \item If \( \tuple{P, \leq} \) is a poset and \( A \subseteq P \), then \( u \) is a \keyword{maximal element} of \( A \)
              if and only if \( u \in A \land \lnot \, \exists x\,{\in}\,A \, \big[ u \leq x \land x \neq u \big]\).
        \item If \( \tuple{P, <} \) is a poset and \( A \subseteq P \), then \( u \) is a \keyword{maximal element} of \( A \)
              if and only if \( u \in A \land \lnot \, \exists x\,{\in}\,A \, \big[ u < x \big] \)
    \end{enumerate}
\end{dfn}

\begin{rem}
    It may be instructive to think of \keyword{minimal} and \keyword{maximal} elements as \emph{possible candidates} for the least and greatest element.
\end{rem}

\begin{thm}
    \label{thm:uniqueness-of-least-element}
    Let \( \tuple{P, <} \) be a poset and \( A \) be any subset of \( P \). The least element of \( A \) (if it exists) is unique.
\end{thm}

\begin{proof}
    Let \( \tuple{P, <} \) be a poset. Let \( A \subseteq P \) be arbitrary. Suppose that \( A \) has a least element \( b \).
    By \cref{dfn:least-element}, it follows that,
    \begin{equation}
        \label{eq:b-is-the-least-element-of-A}
        b \in A \land \forall x\,{\in}\,A \, \big[ b < x \lor x = b \big]
    \end{equation}
    Suppose that there exists another \( b^{\prime} \in A \) such that the conditions in \cref{dfn:least-element} hold. That is,
    \begin{equation}
        \label{eq:bprime-is-the-least-element-of-A}
        b^{\prime} \in A \land \forall x\,{\in}\,A \, \big[ b^{\prime} < x \lor x = b^{\prime} \big]
    \end{equation}
    Thus, the application of \aref{rule-of-inference:universal-instantiation}{Universal Instantiation} to \eqref{eq:b-is-the-least-element-of-A} and \eqref{eq:bprime-is-the-least-element-of-A}
    yields,
    \begin{equation*}
        \mleft( b < b^{\prime} \lor b^{\prime} = b \mright)
        \land
        \mleft( b^{\prime} < b \lor b = b^{\prime} \mright)
    \end{equation*}
    Which is equivalent to,
    \begin{equation*}
        \mleft( b < b^{\prime} \land b^{\prime} < b \mright) \lor \mleft( b = b^{\prime} \mright)
    \end{equation*}
    Since \( < \) is asymmetric, \( \mleft( b < b^{\prime} \land b^{\prime} < b \mright) \) cannot possibly be true. Thus, we conclude
    by \aref{rule-of-inference:disjunctive-syllogism}{Disjunctive Syllogism} that \( b = b^{\prime} \). Thus, the least element (if it exists) is unique.
\end{proof}

\begin{thm}
    \label{thm:uniqueness-of-greatest-element}
    Let \( \tuple{P, <} \) be a poset and \( A \) be any subset of \( P \). The greatest element of \( A \) (if it exists) is unique.
\end{thm}

\begin{proof}
    Let \( \tuple{P, <} \) be a poset. Let \( A \subseteq P \) be arbitrary. Suppose that \( u \) is the greatest element of \( A \).
    By \cref{dfn:greatest-element}, it follows that,
    \begin{equation}
        \label{eq:u-is-the-greatest-element-of-A}
        u \in A \land \forall x\,{\in}\,A \, \big[ x < u \lor x = u \big]
    \end{equation}
    Let \( u^{\prime} \) be another element of \( A \) that satisfies the conditions stated in \cref{dfn:greatest-element}. Thus,
    \begin{equation}
        \label{eq:uprime-is-the-greatest-element-of-A}
        u^{\prime} \in A \land \forall x\,{\in}\,A \, \big[ x < u^{\prime} \lor x = u^{\prime} \big]
    \end{equation}
    By applying \aref{rule-of-inference:universal-instantiation}{Universal Instantiation} on \eqref{eq:u-is-the-greatest-element-of-A} and \eqref{eq:uprime-is-the-greatest-element-of-A},
    we derive,
    \begin{equation*}
        \mleft( u^{\prime} < u \lor u^{\prime} = u \mright) \land \mleft( u < u^{\prime} \lor u = u^{\prime} \mright)
    \end{equation*}
    Which is equivalent to,
    \begin{equation*}
        \mleft( u^{\prime} < u \land u < u^{\prime} \mright) \lor \mleft( u = u^{\prime} \mright)
    \end{equation*}
    However, the asymmetric property of \( < \) means that \( u^{\prime} < u \land u < u^{\prime} \) must be false. Thus, we conclude
    by \aref{rule-of-inference:disjunctive-syllogism}{Disjunctive Syllogism} that \( u = u^{\prime} \).
\end{proof}

\vspace{0.2cm}
\noindent
Since the uniqueness of the least and greatest elements of a set is guaranteed by
\cref{thm:uniqueness-of-least-element,thm:uniqueness-of-greatest-element},
we can introduce special notations to denote the least and greatest elements of a set.

\begin{dfn}
    \label{dfn:least-element-notation}
    For any poset \( \tuple{P, <} \) and any \( A \subseteq P \), the least and greatest elements of \( A \) (if they exist)
    are denoted by \( \minimum A \) and \( \maximum A \) respectively.
\end{dfn}

\begin{thm}
    \label{thm:least-element-is-also-a-minimal-element}
    For any poset \( \tuple{P, <} \) and any \( A \subseteq P \), the least element of \( A \) (if it exists) is also a
    minimal element of \( A \).
\end{thm}

\begin{proof}
    Let \( b \) be the least element of \( A \). By \cref{dfn:least-element}, it follows that,
    \begin{equation}
        \label{eq:express-the-fact-that-b-is-the-least-element-of-A}
        b \in A \land \forall x\,{\in}\,A \, \big[ b < x \lor x = b \big]
    \end{equation}
    Let \( x_{0} \in A \) be arbitrary. By \aref{rule-of-inference:universal-instantiation}{Universal Instantiation} of \eqref{eq:express-the-fact-that-b-is-the-least-element-of-A},
    we have,
    \begin{equation}
        \label{eq:two-disjuncts-for-b-and-x0}
        b < x_{0} \lor x_{0} = b
    \end{equation}
    If \( b < x_{0} \), then clearly \( \lnot \, \mleft( x_{0} < b \mright) \) due to the asymmetric property of \( < \).
    If \( x_{0} = b \), then \( \lnot \, \mleft( x_{0} < b \mright) \) is also true because \( < \) is irreflexive. Since
    both disjuncts in \eqref{eq:two-disjuncts-for-b-and-x0} implies \( \lnot \, \mleft( x_{0} < b \mright) \), we can derive
    \( \lnot \, \mleft( x_{0} < b \mright) \) by \aref{rule-of-inference:disjunction-elimination}{Disjunction Elimination}.
    And by the \aref{rule-of-inference:deduction-theorem}{Deduction Theorem}, we have \( x_{0} \in A {\implies} \lnot \, \mleft( x_{0} < b \mright) \) for any \( x_{0} \).
    \begin{align*}
        &\forall x\,{\in}\,A \, \big[ \lnot \, \mleft( x < b \mright) \big]\\
        \therefore \, \lnot \, &\exists x\,{\in}\,A \, \big[ x < b \big]
    \end{align*}
    Since \( b \in A \) and \( \lnot \, \exists x\,{\in}\,A \, \big[ x < b \big] \),
    it follows from \cref{dfn:minimal-element} that \( b \) is a minimal element.
\end{proof}

\begin{thm}
    \label{thm:greatest-element-is-also-a-maximal-element}
    For any poset \( \tuple{P, <} \) and any \( A \subseteq P \), the greatest element of \( A \) (if it exists) is also a
    maximal element of \( A \).
\end{thm}

\begin{proof}
    Let \( u \) be the greatest element of \( A \). By \cref{dfn:greatest-element}, we have,
    \begin{equation}
        \label{eq:express-the-fact-that-u-is-the-greatest-element}
        u \in A \land \forall x\,{\in}\,A \, \big[ x < u \lor x = u \big]
    \end{equation}
    Let \( x_{0} \in A \) be arbitrary. By \aref{rule-of-inference:universal-instantiation}{Universal Instantiation} of \eqref{eq:express-the-fact-that-u-is-the-greatest-element}, we derive,
    \begin{equation*}
        x_{0} < u \lor x_{0} = u
    \end{equation*}
    If \( x_{0} < u \), then \( \lnot \, \mleft( u < x_{0} \mright) \) due to the asymmetric property of \( < \). If \( x_{0} = u \), then
    \( \lnot \, \mleft( u < x_{0} \mright) \) is also true because \( < \) is irreflexive. Since \( \lnot \, \mleft( u < x_{0} \mright) \)
    is implied in all possible cases, it follows from \aref{rule-of-inference:disjunction-elimination}{Disjunction Elimination} and the \aref{rule-of-inference:deduction-theorem}{Deduction Theorem} that
    \( x_{0} \in A {\implies} \lnot \, \mleft( u < x_{0} \mright) \) for any \( x_{0} \).
    \begin{align*}
        &\forall x\,{\in}\,A \, \big[ \lnot \, \mleft( u < x \mright) \big]\\
        \therefore \, \lnot \, &\exists x\,{\in}\,A \, \big[ u < x \big]
    \end{align*}
    Since \( u \in A \) and \( \lnot \, \exists x\,{\in}\,A \, \big[ u < x \big] \), it follows from \cref{dfn:maximal-element} that
    \( u \) is a maximal element.
\end{proof}

\begin{thm}
    \label{thm:minimal-element-is-the-least-element-for-tosets}
    For any toset \( \tuple{T, <} \) and any \( A \subseteq T \), the minimal element of \( A \) is also the least element of \( A \).
\end{thm}

\begin{proof}
    Let \( \tuple{T, <} \) be a toset and \( A \) be an arbitrary subset of \( T \). Suppose that \( b \) is a minimal element of \( A \).
    So by \cref{dfn:minimal-element}, it follows that \( b \in A \), and also \( \lnot \, \exists x\,{\in}\,A \, \mleft[ x < b \mright] \), which is
    logically equivalent to \( \forall x\,{\in}\,A \, \mleft[ \lnot \, \mleft( x < b \mright) \mright] \). Now let \( x_{0} \in A \) be arbitrary.
    Applying \aref{rule-of-inference:universal-instantiation}{Universal Instantiation} on \( \forall x\,{\in}\,A \, \mleft[ \lnot \, \mleft( x < b \mright) \mright] \) yields \( \lnot \, \mleft( x_{0} < b \mright) \).
    But \( \tuple{T, <} \) is a toset and \( A \subseteq T \), so it follows that,
    \begin{equation*}
        \mleft(x_{0} = b\mright) \lor \mleft(x_{0} < b\mright) \lor \mleft(b < x_{0}\mright)
    \end{equation*}
    Since \( \lnot \, \mleft( x_{0} < b \mright) \), we conclude by \aref{rule-of-inference:disjunctive-syllogism}{Disjunctive Syllogism} that,
    \begin{equation*}
        \mleft( x_{0} = b \mright) \lor \mleft( b < x_{0} \mright)
    \end{equation*}
    Since \( x_{0} \in A \) is arbitrary, \aref{rule-of-inference:universal-generalization}{Universal Generalization} yields,
    \begin{equation*}
        \forall x\,{\in}\,A \, \big[ \mleft( x_{0} = b \mright) \lor \mleft( b < x_{0} \mright) \big]
    \end{equation*}
    Since \( b \in A \), it follows from \cref{dfn:least-element} that \( b = \min A \).
\end{proof}

\begin{thm}
    \label{thm:maximal-element-is-the-greatest-element-for-tosets}
    For any toset \( \tuple{T,<} \) and any \( A \subseteq T \), the maximal element of \( A \) is also
    the greatest element of \( A \).
\end{thm}

\begin{proof}
    Let \( \tuple{T, <} \) be a toset and \( A \) be an arbitrary subset of \( T \). Suppose that \( u \) is a maximal element of \( A \).
    By \cref{dfn:maximal-element}, it follows that \( u \in A \) and also \( \lnot \, \exists x\,{\in}\,A \, \mleft[ u < x \mright] \). Note that
    the latter is logically equivalent to \( \forall x\,{\in}\,A \, \mleft[ \lnot \, \mleft( u < x \mright) \mright] \).
    Let \( x_{0} \in A \) be arbitrary. Instantiating \( \forall x\,{\in}\,A \, \mleft[ \lnot \, \mleft( u < x \mright) \mright] \) yields,
    \begin{equation*}
        \lnot \, \mleft( u < x_{0} \mright)
    \end{equation*}
    Moreover, since \( \tuple{T,<} \) is a toset, it follows that,
    \begin{equation*}
        \mleft( x_{0} = u \mright) \lor \mleft( x_{0} < u \mright) \lor \mleft( u < x_{0} \mright)
    \end{equation*}
    Since \( \lnot \, \mleft( u < x_{0} \mright) \), we conclude by \aref{rule-of-inference:disjunctive-syllogism}{Disjunctive Syllogism} that,
    \begin{equation*}
        \mleft( x_{0} = u \mright) \lor \mleft( x_{0} < u \mright)
    \end{equation*}
    Since \( x_{0} \in A \) is arbitrary, we can then generalize,
    \begin{equation*}
        \forall x\,{\in}\,A \, \mleft[ \mleft( x = u \mright) \lor \mleft( x < u \mright) \mright]
    \end{equation*}
    It follows immediately from \cref{dfn:greatest-element} that \( u \) is the greatest element of \( A \).
\end{proof}

\begin{thm}
    \label{thm:any-subset-of-a-poset-forms-a-poset}
    If \( \tuple{P, <} \) is a poset and \( \prec \,\,=\,\, < \intersect \, A^{2} \) for some \( A \subseteq P \),
    then \( \tuple{A, \prec} \) is also a poset.
\end{thm}

\begin{proof}
    Let \( \tuple{P, <} \) be a poset and \( A \subseteq P \) be arbitrary. Let the binary relation \( \prec \) be defined
    as,
    \begin{equation*}
        \prec \,\,=\,\, < \intersect \, A^{2}
    \end{equation*}
    To prove that \( \tuple{A, \prec} \) is a poset, we must prove several claims:
    \begin{enumerate*}[label=(\alph*)]
        \item \( \prec \,\, \subseteq A^{2} \),
        \item \( \prec \) is asymmetric,
        \item \( \prec \) is transitive.
    \end{enumerate*}

    \vspace{0.5cm}
    \noindent
    {\scshape Claim} (a): \( \prec \,\, \subseteq A^{2} \).

    \vspace{0.2cm}
    \noindent
    Since \( \prec \,\, = \,\, < \intersect \, A^{2} \), it is obvious that \( \prec \,\, \subseteq A^{2} \) because
    for any \( x \),
    \begin{equation*}
        \mleft( x \in \, \prec \mright) \implies \mleft( x \in \, < \mright) \land \mleft( x \in A^{2} \mright)
        \implies \mleft( x \in A^{2} \mright)
    \end{equation*}

    \noindent
    {\scshape Claim} (b): \( \prec \) is asymmetric.

    \vspace{0.2cm}
    \noindent
    Let \( x_{0} \in A \) and \( y_{0} \in A \) be arbitrary. If \( x_{0} \prec y_{0} \), then \( x_{0} < y_{0} \) by definition.
    But \( < \) is asymmetric, so it follows that \( \lnot \, \mleft( y_{0} < x_{0} \mright) \). Thus \( \lnot \, \mleft( y_{0} \prec x_{0} \mright) \)
    by definition. \( \prec \) is asymmetric.

    \vspace{0.5cm}
    \noindent
    {\scshape Claim} (c): \( \prec \) is transitive.

    \vspace{0.2cm}
    \noindent
    Let \( x_{0} \), \( y_{0} \), and \( z_{0} \) be arbitrary elements of \( A \). Suppose that \( x_{0} \prec y_{0} \) and \( y_{0} \prec z_{0} \).
    Thus, \( x_{0} < y_{0} \) and \( y_{0} < z_{0} \) by definition. But \( < \) is transitive, so we have \( x_{0} < z_{0} \). Since \( x_{0} < z_{0} \)
    and \( x_{0} \in A \land z_{0} \in A \), \( x_{0} \prec z_{0} \) follows by definition. Thus, \( \prec \) is indeed transitive.

    \vspace{0.3cm}
    \noindent
    We have shown that \( \prec \) is a binary relation on \( A \) that is asymmetric and transitive. Thus, \( \tuple{A, \prec} \) is a poset.
\end{proof}

\begin{dfn}[Upper Bound]
    \label{dfn:upper-bound}
    If \( \tuple{P, <} \) is a poset and \( A \subseteq P \), then \( b \in P \) is an \keyword{upper bound} or \keyword{majorant} of \( A \)
    if and only if \( \forall x\,{\in}\,A \, \mleft[ x < b \lor x = b \mright] \). The set of all upper bounds of \( A \)
    is denoted by \( \upperbounds{A} \).
    \begin{equation*}
        \upperbounds{A} = \setbuilderbig{x \in P}{\forall z\,{\in}\,A \, \mleft( z < x \lor x = z \mright)}
    \end{equation*}
\end{dfn}

\begin{dfn}[Lower Bound]
    \label{dfn:lower-bound}
    If \( \tuple{P, <} \) is a poset and \( A \subseteq P \), then \( b \in P \) is a \keyword{lower bound} or \keyword{minorant} of \( A \)
    if and only if \( \forall x\,{\in}\,A \, \mleft[ b < x \lor x = b \mright] \). The set of all lower bounds of \( A \) is denoted by
    \( \lowerbounds{A} \).
    \begin{equation*}
        \lowerbounds{A} = \setbuilderbig{x \in P}{\forall z\,{\in}\,A \, \mleft( x < z \lor x = z \mright)}
    \end{equation*}
\end{dfn}

\begin{dfn}[Infimum]
    \label{dfn:infimum}
    For any poset \( \tuple{P, <} \) and any \( A \subseteq P \), \( b \) is the \keyword{infimum} of \( A \) in \( \tuple{P, <} \)
    if and only if it is the greatest element of \( \lowerbounds{A} \). The infimum of \( A \) can be denoted as \( \infimum A \).
\end{dfn}

\begin{dfn}[Supremum]
    \label{dfn:supremum}
    For any poset \( \tuple{P, <} \) and any \( A \subseteq P \), \( b \) is the \keyword{supremum} of \( A \) in \( \tuple{P, <} \)
    if and only if it is the least element of \( \upperbounds{A} \). The supremum of \( A \) can be denoted as \( \supremum A \).
\end{dfn}

\begin{thm}
    \label{thm:least-element-is-the-infimum}
    For any poset \( \tuple{P, <} \) and any \( A \subseteq P \), if \( b = \minimum A \), then \( b = \infimum A \).
\end{thm}

\begin{proof}
    Let \( \tuple{P, <} \) be a poset and \( A \subseteq P \). Clearly, \( \lowerbounds{A} \) exists and is given by,
    \begin{equation}
        \label{eq:write-out-the-definition-for-the-set-of-all-lower-bounds-of-A}
        \lowerbounds{A} = \setbuilderBig{x \in P}{\forall z\,{\in}\,A \, \big( x < z \lor x = z  \big)}
    \end{equation}
    Suppose that \( b = \minimum A \). Then by \cref{dfn:least-element}, we have,
    \begin{equation}
        \label{eq:b-is-the-least-element-of-A-2}
        b \in A \land \forall x\,{\in}\,A \, \big[ b < x \lor x = b \big]
    \end{equation}
    Since \( b \in A \) and \( A \subseteq P \), it follows that \( b \in P \). By \eqref{eq:write-out-the-definition-for-the-set-of-all-lower-bounds-of-A}
    and \eqref{eq:b-is-the-least-element-of-A-2}, it is clear that \( b \in \lowerbounds{A} \).
    Now let \( x_{0} \in \lowerbounds{A} \) be arbitrary. Thus, we have by \eqref{eq:write-out-the-definition-for-the-set-of-all-lower-bounds-of-A},
    \begin{equation}
        \label{eq:applying-definition-of-set-of-all-lower-bounds-to-x0}
        x_{0} \in P \land \forall z\,{\in}\,A \, \big[ x_{0} < z \lor x_{0} = z \big]
    \end{equation}
    Instantiating \eqref{eq:applying-definition-of-set-of-all-lower-bounds-to-x0} with \( b \) yields,
    \begin{equation*}
        b \in A \implies \big( x_{0} < b \lor x_{0} = b \big)
    \end{equation*}
    Since \( b \in A \), we conclude by \aref{rule-of-inference:modus-ponens}{Modus Ponens} that \( \mleft( x_{0} < b \mright) \lor \mleft( x_{0} = b \mright) \). Since \( x_{0} \in \lowerbounds{A} \)
    is arbitrary, \aref{rule-of-inference:universal-generalization}{Universal Generalization} applies, and we have, \( \forall x\,{\in}\,\lowerbounds{A} \, \big[ x < b \lor x = b \big] \). Since \( b \in \lowerbounds{A} \)
    is also true, it follows from \cref{dfn:greatest-element} that \( b = \maximum \lowerbounds{A} \). Thus, \( b = \infimum A \) by \cref{dfn:infimum}.
\end{proof}

\begin{thm}
    \label{thm:greatest-element-is-the-supremum}
    For any poset \( \tuple{P, <} \) and any \( A \subseteq P \), if \( u = \maximum A \), then \( u = \supremum A \).
\end{thm}

\begin{proof}
    Let \( \tuple{P, <} \) be a poset and \( A \subseteq P \). Clearly, \( \upperbounds{A} \) exists and is given by,
    \begin{equation}
        \label{eq:state-the-definition-of-the-set-of-all-upper-bounds-of-A}
        \upperbounds{A} = \setbuilderBig{x \in P}{\forall z\,{\in}\,A \, \big( z < x \lor z = x \big)}
    \end{equation}
    If \( u = \maximum A \), then by \cref{dfn:greatest-element}, it follows that,
    \begin{equation}
        \label{eq:u-is-the-greatest-element-of-A-2}
        u \in A \land \forall x\,{\in}\,A \, \big[ x < u \lor x = u \big]
    \end{equation}
    Since \( A \subseteq P \) and \( u \in A \) by \eqref{eq:u-is-the-greatest-element-of-A-2}, we deduce,
    \begin{equation*}
        u \in P \land \forall z\,{\in}\,A \, \big[ z < u \lor z = u \big]
    \end{equation*}
    Thus, \( u \in \upperbounds{A} \) by \eqref{eq:state-the-definition-of-the-set-of-all-upper-bounds-of-A}.

    \vspace{0.5cm}
    \noindent
    Now let \( x_{0} \in \upperbounds{A} \) be arbitrary. By \eqref{eq:state-the-definition-of-the-set-of-all-upper-bounds-of-A},
    it follows that,
    \begin{equation}
        \label{eq:applying-the-definition-of-the-set-of-all-upper-bounds-of-A-to-x0}
        \mleft( x_{0} \in P \mright) \land \forall z\,{\in}\,A \, \big[ z < x_{0} \lor z = x_{0} \big]
    \end{equation}
    Instantiating \eqref{eq:applying-the-definition-of-the-set-of-all-upper-bounds-of-A-to-x0} with \( u \) yields,
    \begin{equation*}
        \mleft( u \in A \mright) \implies \big[ \mleft( u < x_{0} \mright) \lor \mleft( u = x_{0} \mright) \big]
    \end{equation*}
    We know from \eqref{eq:u-is-the-greatest-element-of-A-2} that \( u \in A \).
    Thus, we conclude by \aref{rule-of-inference:modus-ponens}{Modus Ponens} that,
    \( \mleft( u < x_{0} \mright) \lor \mleft( u = x_{0} \mright)\).
    Since \( x_{0} \in \upperbounds{A} \)
    is arbitrary, we can generalize:
    \begin{equation*}
        \forall x\,{\in}\,\upperbounds{A} \, \big[ u < x \lor u = x \big]
    \end{equation*}
    Since \( u \in \upperbounds{A} \), it follows from \cref{dfn:least-element} that \( u = \minimum \upperbounds{A} \). Thus, \( u = \supremum A \)
    by \cref{dfn:supremum}.
\end{proof}

\noindent
We conclude this section by introducing an important concept.
\begin{dfn}[Well-Ordering]
    \label{dfn:well-ordering}
    For any set \( W \), a binary relation \( < \) is a \keyword{well-ordering} on \( W \) if and only if:
    \begin{enumerate}[topsep=2pt,label=\texttt{\textup{(\alph*)}}]
        \item \( < \) is a total ordering,
        \item \( \forall S \,{\subseteq}\, W \, \big[ \mleft( S \neq \emptyset \mright) {\implies}
                \exists b\,{\in}\,S \, \mleft( b = \minimum S \mright) \big] \).
    \end{enumerate}
\end{dfn}

\begin{dfn}[Woset]
    \label{dfn:woset}
    If the binary relation \( < \) is a well-ordering on \( W \), then the tuple \( \tuple{W, <} \) is
    called a \keyword{well-ordered set} (or \keyword{woset} for short).
\end{dfn}

\begin{thm}
    \label{thm:every-element-of-a-woset-as-long-as-it-is-not-the-greatest-element-has-a-unique-successor}
    If \( \tuple{W, <} \) is a woset and \( x \) is an element of \( W \) that isn't the greatest element of \( W \),
    then there exists a unique \( y \in W \) such that there exists no element that lies in between \( x \) and \( y \).
    \begin{equation*}
        \forall x\,{\in}\,W \, \Big[ x \neq \maximum W
        \implies
        \exists! y\,{\in}\,W \, \Big( x < y \land \forall z\,{\in}\,W \, \mleft( x < z \implies \mleft( y < z \lor y = z \mright) \mright) \Big)  \Big]
    \end{equation*}
\end{thm}

\begin{proof}
    Let \( \tuple{W, <} \) be a woset. Let \( x_{0} \in W \) such that \( x_{0} \) is not the greatest element of \( W \). Now consider the set,
    \begin{equation*}
        G = \setbuilderBig{x \in W}{x_{0} < x}
    \end{equation*}
    Clearly, \( G \subseteq W \). And since \( x_{0} \) is not the greatest element of \( W \), it follows from \cref{dfn:greatest-element} that,
    \begin{equation*}
        x_{0} \notin W \lor \lnot \, \forall x\,{\in}\,W \, \big[ x < x_{0} \lor x = x_{0} \big]
    \end{equation*}
    Since \( x_{0} \in W \), we conclude by \aref{rule-of-inference:disjunctive-syllogism}{Disjunctive Syllogism} that \( \lnot \, \forall x\,{\in}\,W \, \mleft[ x < x_{0} \lor x = x_{0} \mright] \),
    which is logically equivalent to \( \exists x\,{\in}\,W \, \mleft[ \lnot \, \mleft( x < x_{0} \mright) \land x \neq x_{0} \mright] \). Therefore,
    there exists some \( w \in W \) such that \( \lnot \, (w < x_{0}) \land w \neq x_{0} \). Every woset is a toset, so \( \tuple{W, <} \) is clearly a
    toset. Thus, \( \lnot \, \mleft( w < x_{0} \mright) \) implies \( \mleft(  x_{0} < w \mright)  \lor \mleft(x_{0} = w \mright) \).
    But it has already been proven that \( x_{0} \neq w \), so necessarily \( x_{0} < w \) by \aref{rule-of-inference:disjunctive-syllogism}{Disjunctive Syllogism}.
    In other words, \( w \in G \). Thus,
    \( G \neq \emptyset \).

    \vspace{0.3cm}
    \noindent
    Since \( G \subseteq W \) and \( G \neq \emptyset \), it follows from \cref{dfn:woset} that \( G \) has a least element \( y_{0} \). Since
    \( y_{0} \in G \), necessarily \( x_{0} < y_{0} \). Let \( z_{0} \in W \) such that \( x_{0} < z_{0} \). Since \( y_{0} = \minimum G \), it
    follows from \cref{dfn:least-element} that \( y_{0} < z_{0} \lor y_{0} = z_{0} \). By \aref{rule-of-inference:universal-generalization}{Universal Generalization},
    \begin{equation*}
        \forall z\,{\in}\,W \, \big[ x_{0} < z \implies \mleft( y_{0} < z \lor y_{0} = z \mright) \big]
    \end{equation*}
    Since \( x_{0} < y_{0} \), we have,
    \begin{equation*}
        x_{0} < y_{0} \land \forall z\,{\in}\,W \, \big[ x_{0} < z \implies \mleft( y_{0} < z \lor y_{0} = z \mright) \big]
    \end{equation*}
    In other words,
    \begin{equation*}
        \exists y\,{\in}\,W \, \Big( x_{0} < y \land \forall z\,{\in}\,W \, \big[ x_{0} < z \implies \mleft( y < z \lor y = z \mright) \big] \Big)
    \end{equation*}
    Lastly, to prove uniqueness, suppose that \( y \) and \( y^{*} \) are two elements of \( W \) such that,
    \begin{gather*}
        x_{0} < y \land \forall z\,{\in}\,W \, \big[ x_{0} < z \implies \big( y < z \lor y = z \big) \big]\\
        x_{0} < y^{*} \land \forall z\,{\in}\,W \, \big[ x_{0} < z \implies \big( y^{*} < z \lor y^{*} = z \big) \big]
    \end{gather*}
    Using the definition for \( G \), we can rewrite the right conjuncts as,
    \begin{gather*}
        \forall z \, \big[ \mleft( z \in W \land z \in G \mright) \implies \mleft( y < z \lor y = z \mright) \big] \\
        \forall z \, \big[ \mleft( z \in W \land z \in G \mright) \implies \mleft( y^{*} < z \lor y^{*} = z \mright) \big]
    \end{gather*}
    Since \( G \subseteq W \), it follows that \( z \in G {\implies} z \in W \) for any \( z \) so that \( z \in W \land z \in G {\iff} z \in G \) for any \( z \).
    Thus, the formulas can be equivalently expressed as,
    \begin{gather*}
        \forall z \, \big[ z \in G \implies \mleft( y < z \lor y = z \mright) \big]\\
        \forall z \, \big[ z \in G \implies \mleft( y^{*} < z \lor y^{*} = z \mright) \big]
    \end{gather*}
    Clearly, \( y \) and \( y^{*} \) are the least elements of \( G \). But by \cref{thm:uniqueness-of-least-element}, the least element
    is unique, so necessarily \( y = y^{*} \).
\end{proof}

\newpage
\section{Operations}
An n-ary relation on a set \( A \) tells us whether \( n \) objects from \( A \) are related to one another; an n-ary operation, on the other hand, \emph{operates}
on a tuple of objects from \( A \) to yield another object from \( A \).

\begin{dfn}[N-Ary Operation]
    \label{dfn:n-ary-operation}
    For any non-empty set \( A \) and \( n \in \naturalset \), an \keyword{n-ary operation} on \( A \) is a function
    \( \function{f}{D}{A} \) where \( D \subseteq A^{n} \) and \( D \neq \emptyset \).
    The natural number \( n \) is called the \keyword{arity} of \( f \).
\end{dfn}

\begin{cor}
    \phantomsection \label{cor:0-ary-operations}
    If \( f \) is a \( 0 \)-ary operation on a non-empty set \( A \), then \( f = \mleft\{ \orderedpair{\emptyset}{a} \mright\} \)
    for some \( a \in A \).
\end{cor}

\begin{rem}
    The semantics of classical first-order logic requires n-ary operations to have the entirety of \( A^{n} \) as their domains; however, it can be shown
    that the semantics of first-order languages can be easily modified to accomodate partial functions on \( A^{n} \).
\end{rem}

\begin{notation}
    \phantomsection \label{notation:functional-notation-tuple-version}
    For any function \( f \), if \( \orderedpair{x}{y} \in f \) for some \( x \in \dom f \) where \( x = \tuple{a_{0}, \dots, a_{n-1}} \) for some
    \( a_{0},\dots,a_{n-1} \), then \( y \) is called the \keyword{value} of \( f \) at \( \tuple{a_{0}, \dots, a_{n-1}} \) and is denoted by \( f(a_{0},\dots,a_{n-1}) \).
    This notational convention is just an abbreviation for the notation introduced in \cref{dfn:functional-notation}. As in,
    \( f(a_{0},\dots, a_{n-1}) \coloneq  f(\tuple{a_{0}, \dots, a_{n-1}}) \).
\end{notation}

\begin{notation}[Arity]
    \phantomsection \label{notation:arity-of-n-ary-operations}
    If \( f \) is an operation on \( A \), we denote the arity of \( f \) as \( \arity{f} \).
\end{notation}

\begin{dfn}
    \label{dfn:closed-operation}
    Let \( f \) be an n-ary operation on a non-empty set \( A \). Let \( S \) be a non-empty subset of \( A \).
    We say that \( S \) is \keyword{closed under} \( f \) if and only if,
    \begin{equation*}
        \forall x_{0}\,{\in}\,S \, \dots \forall x_{n-1}\,{\in}\,S \, \Big[ \tuple{x_{0}, \dots, x_{n-1}} \in \dom f \implies f(x_{0},\dots,x_{n-1}) \in S \Big]
    \end{equation*}
\end{dfn}

\noindent
We now introduce several important concepts pertaining to binary operations.

\begin{notation}[Infix]
    \label{notation:infix-notation-for-binary-operations}
    Let \( A \) be a non-empty set and \( \ast \) be a binary operation on \( A \). If \( x \) and \( y \) are elements of \( A \),
    then the value of \( \ast \) at \( \tuple{x,y} \) can be denoted as \( x \ast y \) if and only if \( \tuple{x,y} \in \dom (\ast) \).
    
\end{notation}

\begin{dfn}[Idempotence]
    \phantomsection
    \label{dfn:idempotence-of-element-under-binary-operations}
    If \( A \) is a non-empty set and \( \ast \) is a binary operation on \( A \), then an element \( x \in A \) is \keyword{idempotent} under
    \( \ast \) if and only if \( x \ast x = x \).
\end{dfn}

\begin{dfn}[Associativity]
    \label{dfn:associativity-of-binary-operations}
    Let \( A \) be a non-empty set and \( \ast \) be a binary operation on \( A \). The operation \( \ast \)
    is \keyword{associative} if and only if \( x \ast \mleft( y \ast z \mright) = \mleft( x \ast y \mright) \ast z \)
    for any \( x, y, z \in A \) such that \( x \ast y \), \( y \ast z \),
    \( x \ast \mleft( y \ast z \mright) \) and \( \mleft( x \ast y \mright) \ast z \) are all defined.
\end{dfn}

\begin{dfn}[Commutativity]
    \label{dfn:commutativity-of-binary-operations}
    Let \( A \) be a non-empty set and \( \ast \) be a binary operation on \( A \). The operation \( \ast \)
    is \keyword{commutative} if and only if \( x \ast y = y \ast x \) for any \( x,y \in A \) such that \( x \ast y \)
    and \( y \ast x \) are both defined.
\end{dfn}

\begin{dfn}[Distributivity]
    \label{dfn:distributivity-of-binary-operations}
    Let \( A \) be a non-empty set and \( \ast \) and \( \diamond \) be binary operations on \( A \). The operation
    \( \ast \) is:
    \begin{enumerate}[topsep=2pt,label=\textup{\texttt{(\alph*)}}]
        \setlength{\itemsep}{0.3em}
        \item \keyword{left distributive} over \( \diamond \) if and only if
              \( x \ast \mleft( y \diamond z \mright) = \mleft( x \ast y \mright) \diamond \mleft( x \ast z \mright) \)
              for all \( x,y,z \in A \) when all of the operations involved are defined.
        \item \keyword{right distributive} over \( \diamond \) if and only if
              \( \mleft( x \diamond y \mright) \ast z = \mleft( x \ast z \mright) \diamond \mleft( y \ast z \mright) \)
              for all \( x,y,z \in A \) when all of the operations involved are defined.
        \item \keyword{distributive} over \( \diamond \) if and only if \( \ast \) is both left and right distributive
               over \( \diamond \).
    \end{enumerate}
\end{dfn}

\begin{dfn}[Identity]
    \phantomsection
    \label{dfn:identity-element-of-binary-operations}
    Let \( A \) be a non-empty set and \( \ast \) be a binary operation on \( A \). An element \( e \in A \) is:
    \begin{enumerate}[topsep=2pt,noitemsep,label=\textup{\texttt{(\alph*)}}]
        \item a \keyword{left identity} of \( A \) with respect to \( \ast \) if and only if,
            \begin{equation*}
                \forall x\,{\in}\,A \, \big[ \tuple{e, x} \in \dom (\ast) \implies \orderedpair{\tuple{e,x}}{x} \in \ast \big]
            \end{equation*}
        \item a \keyword{right identity} of \( A \) with respect to \( \ast \) if and only if,
            \begin{equation*}
                \forall x\,{\in}\,A \, \big[ \tuple{x, e} \in \dom (\ast) \implies \orderedpair{\tuple{x,e}}{x} \in \ast \big]
            \end{equation*}
        \item a \keyword{two-sided identity} or simply an \keyword{identity} of \( A \) with respect to \( \ast \) if and only if it is both a left and right identity.
    \end{enumerate}
\end{dfn}

\begin{notation}
    If \( \ast \) is a \emph{total} function on \( A^{2} \), then clearly
    \( \forall x\,{\in}\,A \, \big[ \tuple{e,x} \in \dom (\ast) \,\,{\land}\,\, \tuple{x,e} \in \dom (\ast) \big] \).
    And since \( e \ast x \) and \( x \ast e \) is defined for all \( x \in A \), we can replace
    \( \orderedpair{\tuple{e,x}}{x} \in \ast \) and \( \orderedpair{\tuple{x,e}}{x} \in \ast \) in \cref{dfn:identity-element-of-binary-operations}
    with \( e \ast x = x \) and \( x \ast e = x \).
\end{notation}

\begin{thm}
    \label{thm:uniqueness-of-identity-element}
    Let \( A \) be a non-empty set and \( \ast \) be a binary operation on \( A \). If \( e \in A \) and \( e^{\prime} \in A \) are
    identities of \( A \) with respect to \( \ast \) and either \( e \ast e^{\prime} \) or \( e^{\prime} \ast e \) is defined, then
    \( e^{\prime} = e \).
\end{thm}

\begin{proof}
    We have to prove that
    \( \big( \tuple{e,e^{\prime}} \in \dom (\ast) \,\,{\lor}\,\, \tuple{e^{\prime},e} \in \dom (\ast) \big) {\implies} \mleft( e^{\prime} = e \mright) \).
    Since \( e \) and \( e^{\prime} \) are identities of \( A \) with respect to \( \ast \) and \( e, e^{\prime} \in A \),
    we have by \cref{dfn:identity-element-of-binary-operations}:
    \begin{gather}
        \tuple{e,e^{\prime}} \in \dom (\ast)
        \implies
        \orderedpair{\tuple{e,e^{\prime}}}{e^{\prime}} \in \ast \label{eq:e-star-eprime-is-defined-implies-e-star-eprime-equals-eprime}\\
        \tuple{e^{\prime},e} \in \dom (\ast)
        \implies
        \orderedpair{\tuple{e^{\prime},e}}{e^{\prime}} \in \ast \label{eq:eprime-star-e-is-defined-implies-eprime-star-e-equals-eprime}\\
        \tuple{e^{\prime},e} \in \dom (\ast)
        \implies
        \orderedpair{\tuple{e^{\prime},e}}{e} \in \ast \label{eq:eprime-star-e-is-defined-implies-eprime-star-e-equals-e}\\
        \tuple{e,e^{\prime}} \in \dom (\ast)
        \implies
        \orderedpair{\tuple{e,e^{\prime}}}{e} \in \ast \label{eq:e-star-eprime-is-defined-implies-e-star-eprime-equals-e}
    \end{gather}

    \noindent
    Now suppose that \( \tuple{e,e^{\prime}} \in \dom (\ast) \). It follows from \eqref{eq:e-star-eprime-is-defined-implies-e-star-eprime-equals-eprime}
    and \eqref{eq:e-star-eprime-is-defined-implies-e-star-eprime-equals-e} that \( \orderedpair{\tuple{e,e^{\prime}}}{e^{\prime}} \in \ast \) and
    \( \orderedpair{\tuple{e,e^{\prime}}}{e} \in \ast \). Since \( \ast \) is a function, \( e = e^{\prime} \) follows necessarily from \cref{dfn:function-basic}.
    Since \( e = e^{\prime} \) can be derived from the assumption that \( \tuple{e,e^{\prime}} \in \dom (\ast) \),
    we have by the \aref{rule-of-inference:deduction-theorem}{Deduction Theorem},
    \begin{equation}
        \label{eq:e-star-eprime-defined-implies-equality}
        \tuple{e,e^{\prime}} \in \dom (\ast) \implies \mleft( e^{\prime} = e \mright)
    \end{equation}
    Similarly, if \( \tuple{e^{\prime},e} \in \dom (\ast) \), then
    \( \orderedpair{\tuple{e^{\prime},e}}{e^{\prime}} \in \ast \) and \( \orderedpair{\tuple{e^{\prime},e}}{e} \in \ast \)
    follow from \eqref{eq:eprime-star-e-is-defined-implies-eprime-star-e-equals-eprime} and
    \eqref{eq:eprime-star-e-is-defined-implies-eprime-star-e-equals-e} by \aref{rule-of-inference:modus-ponens}{Modus Ponens}.
    Since \( \ast \) is a function, \( e = e^{\prime} \)
    must hold by \cref{dfn:function-basic}. Applying the \aref{rule-of-inference:deduction-theorem}{Deduction Theorem} once again, we derive,
    \begin{equation}
        \label{eq:eprime-star-e-defined-implies-equality}
        \tuple{e^{\prime},e} \in \dom (\ast) \implies \mleft( e^{\prime} = e \mright)
    \end{equation}
    By \aref{rule-of-inference:constructive-dilemma}{Constructive Dilemma}, we can infer from \eqref{eq:e-star-eprime-defined-implies-equality}
    and \eqref{eq:eprime-star-e-defined-implies-equality} that,
    \begin{equation*}
        \big( \tuple{e,e^{\prime}} \in \dom (\ast) \,\,{\lor}\,\, \tuple{e^{\prime},e} \in \dom (\ast) \big) {\implies} \mleft( e^{\prime} = e \mright) \qedhere
    \end{equation*}
\end{proof}

\begin{cor}
    \label{cor:uniqueness-of-identity-element-for-total-operations}
    If \( \function{\ast}{A^{2}}{A} \), then the identity element (if it exists) of \( A \) with respect to \( \ast \) is always unique.
\end{cor}

\begin{thm}
    \label{thm:equality-of-left-and-right-identities}
    Let \( A \) be a non-empty set and \( \ast \) be a binary operation on \( A \). If \( l \in A \) and \( r \in A \) are left and right
    identites of \( A \) respectively and \( l \ast r \) is defined, then \( l = r \).
\end{thm}

\begin{proof}
    Let \( l \in A \) and \( r \in A \) be left and right identities of \( A \) respectively. Now suppose that
    \( \tuple{l,r} \in \dom (\ast) \). Since \( r \in A \) and \( l \) is a left identity of \( A \), it follows
    from \cref{dfn:identity-element-of-binary-operations} that \( \orderedpair{\tuple{l,r}}{r} \in \ast \). And
    since \( l \in A \) and \( r \) is a right identity of \( A \), we can similarly conclude from \cref{dfn:identity-element-of-binary-operations}
    that \( \orderedpair{\tuple{l,r}}{l} \in \ast \). We know that \( \ast \) is a function, so \cref{dfn:function-basic} implies \( l = r \).
\end{proof}

\begin{dfn}
    \label{dfn:left-and-right-inverses-of-binary-operations}
    Let \( A \) be a non-empty set and \( \ast \) be a binary operation with \( e \in A \) being an identity of \( A \) with respect to \( \ast \).
    Let \( \alpha \) be an element of \( A \). An element \( i \in A \) is:
    \begin{enumerate}[topsep=2pt,noitemsep,label=\textup{\texttt{(\alph*)}}]
        \item a \keyword{left inverse} of \( \alpha \) if and only if \( \orderedpair{\tuple{i,\alpha}}{e} \in \ast \).
        \item a \keyword{right inverse} of \( \alpha \) if and only if \( \orderedpair{\tuple{\alpha,i}}{e} \in \ast \).
    \end{enumerate}
\end{dfn}

\begin{notation}
    When the context makes it clear that \( \tuple{i,\alpha} \in \dom (\ast) \) and \( \tuple{\alpha,i} \in \dom (\ast) \),
    we can express the sentences \( \orderedpair{\tuple{i,\alpha}}{e} \in \ast \) and \( \orderedpair{\tuple{\alpha,i}}{e} \in \ast \) in
    infix notation: \( i \ast \alpha = e \) and \( \alpha \ast i = e \).
\end{notation}

\begin{dfn}
    \phantomsection
    \label{dfn:inverse-element-of-binary-operations}
    Let \( A \) be a non-empty set and \( \ast \) be a binary operation on \( A \). An element \( i \in A \) is an \keyword{inverse}
    of another element \( \alpha \in A \) if and only if it is both the left inverse and the right inverse of \( \alpha \).
\end{dfn}

\begin{thm}
    \label{thm:uniqueness-of-inverse-element-for-total-associative-binary-operations}
    Let \( A \) be a non-empty set and \( \function{\ast}{A^{2}}{A} \) be an associative binary operation on \( A \) with an identity element.
    If \( \alpha \in A \) has an inverse with respect to \( \ast \), then the inverse of \( \alpha \)
    is necessarily unique.
\end{thm}

\begin{proof}
    Let \( i \) and \( j \) be inverses of \( A \) with respect to \( \ast \). Since \( \ast \) is a total function on \( A^{2} \),
    it follows from \cref{cor:uniqueness-of-identity-element-for-total-operations}
    that the identity of \( A \) with respect to \( \ast \) is unique.
    Thus, by \cref{dfn:left-and-right-inverses-of-binary-operations,dfn:inverse-element-of-binary-operations},
    \begin{gather}
        i \ast \alpha = \alpha \ast i = e \label{eq:i-star-alpha-equals-alpha-star-i-equals-ei}\\
        j \ast \alpha = \alpha \ast j = e \label{eq:j-star-alpha-equals-alpha-star-i-equals-ej}
    \end{gather}
    where \( e \) denotes the unique identity of \( A \) with respect to \( \ast \).
    From \eqref{eq:i-star-alpha-equals-alpha-star-i-equals-ei}, we can immediately infer that
    \( \mleft( i \ast \alpha \mright) \ast j = e \ast j = j \).
    But \( \ast \) is associative, so we have
    \( \mleft( i \ast \alpha \mright) \ast j = i \ast \mleft( \alpha \ast j \mright) = i \ast e = i \). Therefore, necessarily \( i = j \).
\end{proof}

\begin{notation}
    \phantomsection \label{notation:inverse-element-of-binary-operations}
    If \( x \in A \) has a unique inverse with respect to a binary operation, the inverse is usually denoted as \( x^{-1} \) or \( -x \) depending on
    the context. The notation \( x^{-1} \) is typically preferred when the binary operation involved resembles multiplication
    whereas \( -x \) is preferred when the operation resembles addition.
\end{notation}

\begin{thm}
    \label{thm:inverse-of-x-star-y}
    If \( A \) is a non-empty set and \( \function{\ast}{A^{2}}{A} \) is an associative binary operation on \( A \) with an identity element,
    then for every \( x,y \in A \) such that \( x \) and \( y \) have inverses with respect to \( \ast \), \( \mleft( x \ast y \mright) \)
    also has an inverse and \( \mleft( x \ast y \mright)^{-1} = y^{-1} \ast x^{-1} \).
\end{thm}

\begin{proof}
    Let \( e \) denote the unique identity of \( A \) with respect to \( \ast \).
    Let \( x \) and \( y \) be arbitrary elements of \( A \). From \cref{dfn:inverse-element-of-binary-operations}, we know that
    \( \mleft( x \ast y \mright) \ast \mleft( x \ast y \mright)^{-1} = e \). Therefore,
    \begin{equation}
        \label{eq:x-inverse-star-e-equals-x-inverse}
        x^{-1} \ast \big[ \mleft( x \ast y \mright) \ast \mleft( x \ast y \mright)^{-1} \big] = x^{-1} \ast e = x^{-1}
    \end{equation}
    Since \( \ast \) is associative, we can simplify the right-hand side of \eqref{eq:x-inverse-star-e-equals-x-inverse} as follows:
    \begin{align}
        \big[ x^{-1} \ast \mleft( x \ast y \mright) \big] \ast \mleft( x \ast y \mright)^{-1}
            &= \big[ \mleft( x^{-1} \ast x \mright) \ast y \big] \ast \mleft( x \ast y \mright)^{-1} \notag \\
            &= \mleft( e \ast y \mright) \ast \mleft( x \ast y \mright)^{-1} \notag \\
            &= y \ast \mleft( x \ast y \mright)^{-1} \label{eq:starred-by-x}
    \end{align}
    From \eqref{eq:x-inverse-star-e-equals-x-inverse} and \eqref{eq:starred-by-x}, we infer that
    \( y \ast \mleft( x \ast y \mright)^{-1} = x^{-1} \) and thus,
    \begin{equation*}
        y^{-1} \ast \big[ y \ast \mleft( x \ast y \mright)^{-1} \big] = y^{-1} \ast x^{-1}
    \end{equation*}
    Again, by associativity of \( \ast \), we can simplify the left-hand side:
    \begin{equation*}
        \mleft( y^{-1} \ast y \mright) \ast \mleft( x \ast y \mright)^{-1} = e \ast \mleft( x \ast y \mright)^{-1}
        = \mleft( x \ast y \mright)^{-1}
    \end{equation*}
    Thus, \( \mleft( x \ast y \mright)^{-1} = y^{-1} \ast x^{-1} \).
\end{proof}

\begin{notation}[Infix]
    \phantomsection \label{notation:generalized-infix-notation}
    If \( x_{0}, \dots, x_{k-1} \) are \( k \) elements of \( A \) and \( \ast \) is a binary operation on \( A \),
    then the notation \( x_{0} \ast \cdots \ast x_{k-1} \) is interpreted from \emph{left to right}. As in:
    \begin{equation*}
        x_{0} \ast \cdots \ast x_{k-1} = \big( x_{0} \ast \big( x_{1} \ast \cdots \ast \mleft( x_{k-2} \ast x_{k-1} \mright) \cdots \big) \big)
    \end{equation*}
\end{notation}

\begin{thm}[Generalized Associative Law]
    \phantomsection \label{thm:generalized-associative-law}
    If \( A \) is a non-empty set and \( \function{\ast}{A^{2}}{A} \) is an associative, binary operation on \( A \),
    then for any \( x_{0}, \dots, x_{k-1} \in A \) where \( k \in \naturalset \) and \( k > 2 \), consecutive applications
    of \( \ast \) to \( x_{0}, \dots, x_{k-1} \) always yields \( x_{0} \ast \cdots \ast x_{k-1} \) regardless of how the
    elements are grouped by parentheses.
\end{thm}

\begin{thm}[Generalized Commutative Law]
    \phantomsection \label{thm:generalized-commutative-law}
    If \( A \) is a non-empty set and \( \function{\ast}{A^{2}}{A} \) is an associative and commutative binary operation on \( A \),
    then for any \( x_{0}, \dots, x_{k-1} \in A \) where \( k \in \naturalset \) and \( k > 2 \), consecutive applications
    of \( \ast \) to \( x_{0}, \dots, x_{k-1} \) always yields the same result regardless of the order in which the elements are arranged.
\end{thm}

\begin{thm}[Generalized Distributive Law]
    \phantomsection \label{thm:generalized-distributive-law}
    If \( A \) is a non-empty set and \( \function{\ast}{A^{2}}{A} \) and \( \function{+}{A^{2}}{A} \) are both binary operations on \( A \)
    such that \( + \) is associative and \( \ast \) distributes over \( + \), then for all \( x,y_{0},\dots,y_{k-1} \in A \) where \( k \in \naturalset \)
    and \( k > 2\),
    \begin{align*}
        x \ast \mleft( y_{0} + \cdots + y_{k-1} \mright)
        &= \mleft( x \ast y_{0} \mright) + \cdots + \mleft( x \ast y_{k-1} \mright)\\
        \mleft( y_{0} + \cdots + y_{k-1} \mright) \ast x
        &= \mleft( y_{0} \ast x \mright) + \cdots + \mleft( y_{k-1} \ast x \mright)
    \end{align*}
    
\end{thm}

\newpage
\section{Structures}
It has been established in \cref{sec:preliminary-remarks-for-mathematical-structures} that, conceptually,
a mathematical structure is the aggregation of a non-empty set and the pieces of information that describe
the mathematical properties of the set's elements. The descriptive information of a structure
typically falls into these categories:
\begin{enumerate*}[label=(\roman*)]
    \item structures,
    \item n-ary relations,
    \item n-ary operations,
    \item functions.
\end{enumerate*}
This leads us to the following definition.

\begin{dfn}[Mathematical Structure]
    \label{dfn:general-mathematical-structures}
    Let \( A \) be a non-empty set.
    A mathematical structure \( \mathfrak{A} \) on \( A \) is the tuple
    \( \tuple{A, S, R, F, G} \) where:
    \begin{enumerate}[topsep=2pt,noitemsep,label=\textup{\texttt{(\alph*)}}]
        \item \( S \) is a tuple of the structures associated with \( \mathfrak{A} \).
        \item \( R \) is a tuple of n-ary relations on \( A \).
        \item \( F \) is a tuple of n-ary operations on \( A \).
        \item \( G \) is a tuple of functions associated with \( \mathfrak{A} \) that aren't n-ary operations on \( A \).
    \end{enumerate}
    The set \( A \) is commonly referred to as the \keyword{universe} or \keyword{underlying set} of the structure \( \mathfrak{A} \).
    The elements collected in \( S \), \( R \), \( F \), \( G \) form the \keyword{parameters} of \( \mathfrak{A} \).
\end{dfn}

\begin{rem}
    For any set \( A \), the simplest possible mathematical structure is \( \tuple{A} \).
    In practice, it is customary to identify the 1-tuple \( \tuple{A} \) with the set \( A \) itself.
    In other words, we can think of sets as degenerate structures
    without any relations, operations, structures, and functions associated with them.
\end{rem}

\begin{notation}
    \phantomsection \label{notation:individual-components-of-mathematical-structures}
    The following notational convention will be used
    to denote the individual components of a structure \( \mathfrak{A} \):
    \begin{enumerate}[topsep=2pt,noitemsep,label=\textup{\texttt{(\alph*)}}]
        \item \( \abs{\mathfrak{A}} \) denotes to the universe of \( \mathfrak{A} \).
        \item \( S^{\mathfrak{A}} \) denotes the tuple of structures associated with \( \mathfrak{A} \).
        \item \( R^{\mathfrak{A}} \) denotes the tuple of n-ary relations on \( \abs{\mathfrak{A}} \).
        \item \( F^{\mathfrak{A}} \) denotes the tuple of n-ary operations on \( \abs{\mathfrak{A}} \).
        \item \( G^{\mathfrak{A}} \) denotes the tuple of functions associated with \( \mathfrak{A} \).
    \end{enumerate}
\end{notation}

\vspace{0.1cm}
\noindent
The astute reader may have noticed the recursive nature of \cref{dfn:general-mathematical-structures}; that is, a structure can have other structures embedded within it.
Recursive structures are fairly common in mathematics. A common example of such a structure is a \aref{dfn:vector-space}{vector space},
which is defined in terms of a second structure.
\begin{equation*}
    \mathcal{V} = \tuple{V, \tuple{\tuple{F, \tuple{}, \tuple{}, \tuple{+, \cdot}, \tuple{}}}, \tuple{}, \tuple{\oplus}, \tuple{\odot} }
\end{equation*}
The inner structure \( \tuple{F, \tuple{}, \tuple{}, \tuple{+, \cdot}, \tuple{}} \) is commonly known as a \aref{dfn:field}{field}.
The reader has undoubtedly noticed the clunkiness of the notation for nested structures.
Naturally, the reason we nest tuples
in \cref{dfn:general-mathematical-structures} is to distinguish between the symbols for the structures, relations, operations, and functions
that are associated with \( \mathfrak{A} \).
Consider, for instance, the structure \( \mathfrak{A} = \tuple{A,*,\&,\#, \Gamma, \Delta} \).
Without the inner tuples to group together the symbols of the same type, it becomes impossible
to interpret the meanings of the symbols in the absence of more information.
Does the symbol \( * \) refer to an operation on \( A \)? A relation?
Conversely, if a more explicit notation is used,
the meaning of the symbols is clear from the get-go.
E.g., for the structure \( \tuple{A, \tuple{}, \tuple{*,\&}, \tuple{\#}, \tuple{\Gamma, \Delta}} \),
we can interpret the symbols as follows:
\( * \) and \( \& \) represent n-ary relations on \( A \);
\( \# \) an n-ary operation on \( A \);
\( \Gamma \) and \( \Delta \) functions associated with \( \mathfrak{A} \).

\vspace{0.3cm}
\noindent
Fortunately, for many structures,
the meaning of the symbols is usually clear from the context,
in which case there is no need to organize the symbols using nested
inner tuples.
For example, if we interpret \( + \) and \( \cdot \) as binary operations on \( F \), \( \oplus \) as a binary operation on \( V \), and \( \odot \) as the
function \( \function{\odot}{\prod \tuple{F, V}}{V} \), then we can simplify the notation for the vector space \( \mathcal{V} \) by omitting all inner tuples:
\begin{equation*}
    \mathcal{V} = \tuple{V, F, +, \cdot, \oplus, \odot}
\end{equation*}

\noindent
We now introduce several concepts that are essential to the study of mathematical structures.

\begin{dfn}[Signature]
    \label{dfn:signature-of-a-structure}
    Given any structure \( \mathfrak{A} \) where \( p = \dom S^{\mathfrak{A}} \), \( n = \dom R^{\mathfrak{A}} \), and \( m = \dom F^{\mathfrak{A}} \)
    for some \( m,n,p \in \naturalset \),
    the \keyword{signature} of \( \mathfrak{A} \) is given by,
    \begin{equation*}
        \sigma \mleft[ \mathfrak{A} \mright] =
        \tuple{
            \tuple{\sigma \mleft[ S^{\mathfrak{A}}_{0} \mright],\dots,\sigma \mleft[ S^{\mathfrak{A}}_{p-1} \mright]},
            \tuple{r_{0}, \dots, r_{n-1}},
            \tuple{f_{0}, \dots, f_{m-1}},
            \dom G^{\mathfrak{A}}
        }
    \end{equation*}
    where \( r_{i} \) and \( f_{j} \) are the arities of \( R^{\mathfrak{A}}_{i} \) and \( F^{\mathfrak{A}}_{j} \) respectively.
\end{dfn}

\begin{dfn}[Isomorphism]
    \label{dfn:isomorphism}
    Two structures \( \mathfrak{A} \) and \( \mathfrak{A}^{\prime} \) are \keyword{isomorphic}, denoted as
    \( \mathfrak{A} \isomorphic \mathfrak{A^\prime} \),
    if and only if the following conditions hold:
    \begin{enumerate}[topsep=2pt,noitemsep,label=\textup{\texttt{(\alph*)}}]
        \item \( \sigma \mleft[ \mathfrak{A} \mright] = \sigma \mleft[ \mathfrak{A^{\prime}} \mright] \)
        \item \( S^{\mathfrak{A}}_{i} \isomorphic S^{\mathfrak{A}^{\prime}}_{i} \) for all \( i \in \dom S^{\mathfrak{A}} \)
        \item There exists a bijection \( \bijection{h}{\abs{\mathfrak{A}}}{\abs{\mathfrak{A^{\prime}}}} \) such that
            for all \( i \in \dom R^{\mathfrak{A}} \) and \( j \in \dom F^{\mathfrak{A}} \),
            \begin{align*}
                \tuple{x_{0}, \dots, x_{r_{i}-1}} \in R^{\mathfrak{A}}_{i} &\iff \tuple{h(x_{0}), \dots, h(x_{r_{i} - 1})} \in R^{\mathfrak{A^{\prime}}}_{i}\\
                \tuple{x_{0},\dots,x_{f_{j}-1}} \in \dom F^{\mathfrak{A}}_{j} &\iff \tuple{h(x_{0}),\dots,h(x_{f_{j}-1})} \in \dom F^{\mathfrak{A^{\prime}}}_{j}\\
                F^{\mathfrak{A}}_{j}(x_{0},\dots,x_{f_{j} - 1}) &= F^{\mathfrak{A}^{\prime}}_{j}(h(x_{0}),\dots,h(x_{f_{j} - 1}))
            \end{align*}
        \item For all \( k \in \dom G^{\mathfrak{A}} \), there exists a pair of bijections
            \( \bijection{\delta_{k}}{\dom G^{\mathfrak{A}}_{k}}{\dom G^{\mathfrak{A}^{\prime}}_{k}} \) and
            \( \bijection{\rho_{k}}{\ran G^{\mathfrak{A}}_{k}}{\ran G^{\mathfrak{A}^{\prime}}_{k}} \) such that
            for any \( x \in \dom G^{\mathfrak{A}}_{k} \),
            \begin{equation*}
                G^{\mathfrak{A}^{\prime}}_{k}(\delta_{k}(x)) = \rho_{k}(G^{\mathfrak{A}}_{k}(x))
            \end{equation*}
    \end{enumerate}
    The tuple \( I = \tuple{h,
                        \tuple{I_{0}, \dots, I_{\dom S^{\mathfrak{A}} - 1}},
                        \tuple{ \tuple{\delta_{0}, \rho_{0}}, \dots, \tuple{ \delta_{\dom G^{\mathfrak{A}} - 1}, \rho_{\dom G^{\mathfrak{A}} - 1}}}
              }\)
    is called an \keyword{isomorphism} between \( \mathfrak{A} \) and \( \mathfrak{A}^{\prime} \).
\end{dfn}

\noindent
Two structures are isomorphic if they ``mirror'' each other exactly:
\begin{enumerate*}[label=(\roman*)]
    \item every component of one structure has a counterpart in the other,
    \item the properties and interrelationship between the components and/or the elements of the
          underlying set are the same in both structures.
\end{enumerate*}
Our goal in this chapter to prove that isomorphic structures share the same mathematical properties.
However, doing so requires model theoretic concepts that we have yet to introduce,
so we will defer the presentation of the proof to the end of the chapter.

\begin{dfn}[Substructure]
    \phantomsection \label{dfn:substructure}
    A structure \( \mathfrak{A} \) is a \keyword{substructure} of another structure \( \mathfrak{B} \) if and only if the following
    conditions hold:
    \begin{enumerate}[topsep=2pt,noitemsep,label=\textup{\texttt{(\alph*)}}]
        \item \( \abs{\mathfrak{A}} \subseteq \abs{\mathfrak{B}} \)
        \item \( \sigma \mleft[ \mathfrak{A} \mright] = \sigma \mleft[ \mathfrak{B} \mright] \)
        \item \( S^{\mathfrak{A}}_{i} \) is a substructure of \( S^{\mathfrak{B}}_{i} \) for all \( i \in \dom S^{\mathfrak{A}} \)
        \item \( R^{\mathfrak{A}}_{i} = R^{\mathfrak{B}}_{i} \intersect \abs{\mathfrak{A}}^{\arity{R^{\mathfrak{A}}_{i}}} \)
            for all \( i \in \dom R^{\mathfrak{A}} \)
        \item \( F^{\mathfrak{A}}_{i} = \restrictedfunction{F^{\mathfrak{B}}_{i}}{D_{i}} \)
            for all \( i \in \dom F^{\mathfrak{A}} \) where \( D_{i} = \abs{\mathfrak{A}}^{\arity{F^{\mathfrak{A}}_{i}}} \)
        \item \( G^{\mathfrak{A}}_{i} \subseteq G^{\mathfrak{B}}_{i} \) for all \( i \in \dom G^{\mathfrak{A}} \).
    \end{enumerate}
\end{dfn}

\begin{notation}
    \phantomsection \label{notation:substructure}
    If \( \mathfrak{A} \) and \( \mathfrak{B} \) are mathematical structures and \( \mathfrak{A} \) is a substructure
    of \( \mathfrak{B} \), then we write \( \mathfrak{A} \leq \mathfrak{B} \).
\end{notation}

\begin{cor}
    \phantomsection \label{cor:every-structure-is-a-substructure-of-itself}
    For any structure \( \mathfrak{A} \), \( \mathfrak{A} \leq \mathfrak{A} \).
\end{cor}

\subsection*{Classification of Structures}
Mathematical structures are classified based on their signatures and properties. By properties, we mean a finite set of axioms (sentences of some first-order language)
that the parameters of the structure must satisfy. There are six main categories of structures in modern mathematics:
\begin{enumerate*}[label=(\roman*)]
    \item ordering,
    \item algebraic,
    \item metric,
    \item topological,
    \item measures,
    \item graphs.
\end{enumerate*}

\subsection{Orderings}
An ordering is any structure equipped with a binary relation that is either:
\begin{enumerate*}[label=(\roman*)]
    \item reflexive, anti-symmetric, and transitive; or
    \item asymmetric and transitive.
\end{enumerate*}
In other words, it is either a poset or a strict poset. Thus, in its simplest
form, an ordering is the structure,
\begin{equation*}
    \mathfrak{A} = \tuple{P, \tuple{  }, \tuple{<}, \tuple{  }, \tuple{  }}
\end{equation*}
where \( P \) is a non-empty set and \( < \) is a binary relation on \( P \) such that \( \mathfrak{A} \)
satisfies the following first-order sentences:
\begin{enumerate}[topsep=2pt,label=\textup{\texttt{(\alph*)}}]
    \item \( \forall x\,{\in}\,P \, \forall y\,{\in}\,P \, \big[ x < y \implies \lnot \, \mleft( y < x \mright) \big] \) (asymmetry)
    \item \( \forall x\,{\in}\,P \, \forall y\,{\in}\,P \, \forall z\,{\in}\,P \, \big[ x < y \land y < z \implies x < z \big] \) (transitivity)
\end{enumerate}
Since the symbols \( < \) and \( \leq \) are universally understood in mathematical contexts to represent order relations, we can simplify the notation:
\begin{equation*}
    \mathfrak{A} = \tuple{P, <}
\end{equation*}
Structures of this form is called a poset. The different variants of ordered structures have already been covered in \cref{ssec:order-relations}, so we will not repeat
them here. The key takeaway here is that it is the set of axioms that determines the type of ordering---or the type of any structure, for that matter.
E.g., for \( \mathfrak{A} \) to be a toset, \( \mathfrak{A} \) must satisfy one additional axiom apart from the other two mentioned above:
\begin{equation*}
    \forall x\,{\in}\,P \, \forall y\,{\in}\,P \, \big[ x \neq y \lor x < y \lor y < x \big]
\end{equation*}
Ordered structures are studied extensively in order theory and lattice theory.

\newpage
\subsection{Algebraic Structures}
A structure is \keyword{algebraic} if and only if it is equipped with operations that satisfy
a finite set of \emph{algebraic identities}---statements that resemble the laws of arithmetic and
elementary algebra such as commutativity, associativity, distributivity, etc.
Algebraic structures capture the forms and patterns behind mathematical ideas, thus allowing the gaps
between seemingly unrelated branches of mathematics to be bridged.
A full exposition of algebraic structures is obviously beyond the scope of this book.
Thus, we will outline only the most common algebraic structures and refer the interested reader to texts on abstract algebra
for an in-depth treatment on the subject.

\subsubsection{Groups}
\begin{dfn}
    \phantomsection
    \label{dfn:group}
    Let \( G \) be a non-empty set and \function{\ast}{G^{2}}{G} be a binary operation on \( G \). The structure
    \( \tuple{G, \ast} \) is a \keyword{group} if and only if the following first-order sentences
    are satisfied:
    \begin{enumerate}[topsep=2pt,noitemsep,label=\textup{\texttt{(\alph*)}}]
        \item \( \ast \) is associative.
            \begin{itemize}
                \item \( \forall x\,{\in}\,G \, \forall y\,{\in}\,G \, \forall z\,{\in}\,G \,
                    \big[ \mleft( x \ast y \mright) \ast z = x \ast \mleft( y \ast z \mright) \big] \)
            \end{itemize}
            \label{enum:group-axioms-associativity}
        \item \( G \) has an identity with respect to \( \ast \) and every element of \( G \) has an inverse.
            \begin{itemize}
                \item \( \exists z\,{\in}\,G \, \forall x\,{\in}\,G \, \big[ \mleft( x \ast z = z \ast x = x \mright)
                        \,\,{\land}\,\,
                        \exists y\,{\in}\,G \, \mleft( x \ast y = y \ast x = z \mright) \big] \)
            \end{itemize}
            \label{enum:group-axioms-identity-and-inverse}
    \end{enumerate}
    The sentences \ref{enum:group-axioms-associativity}, \ref{enum:group-axioms-identity-and-inverse} are collectively known as the \keyword{group axioms}.
\end{dfn}

\begin{dfn}
    \phantomsection
    \label{dfn:abelian-group}
    A group \( \tuple{G, \ast} \) is \keyword{abelian} or \keyword{commutative} if and only if it also satisfies the sentence
    \( \forall x\,{\in}\,G \, \forall y\,{\in}\,G \, \big[ x \ast y = y \ast x \big] \).
\end{dfn}

\begin{thm}
    \label{thm:-group-uniqueness-of-identity}
    For any group \( \tuple{G, \ast} \), \( G \) has a unique identity with respect to \( \ast \).
\end{thm}

\begin{proof}
    This follows readily from \cref{cor:uniqueness-of-identity-element-for-total-operations} because
    \( \ast \) is a total function on \( G^{2} \).
\end{proof}

\begin{thm}
    \label{thm:group-uniqueness-of-inverses}
    For any group \( \tuple{G, \ast} \), the inverse of every element of \( G \) is unique.
\end{thm}

\begin{proof}
    Let \( \alpha \) be any element of \( G \).
    Since \( \tuple{G, \ast} \) is a group, it follows from the \aref{dfn:group}{group axioms} that \( \alpha \) has an inverse
    with respect to \( \ast \) and that \( \ast \) is associative.
    In other words, \cref{thm:uniqueness-of-inverse-element-for-total-associative-binary-operations} applies
    and the inverse of \( \alpha \) is necessarily unique.
\end{proof}

\begin{thm}
    \label{thm:group-basic-properties}
    If \( \tuple{G, \ast} \) is a group and \( e \) is the identity of \( G \) with respect to \( \ast \),
    then the following properties hold for all \( x,y \in G \):
    \begin{enumerate}[topsep=2pt,noitemsep,label=\textup{\texttt{(\alph*)}}]
        \item \( \mleft( x^{-1} \mright)^{-1} = x \) \label{enum:group-inverse-of-inverse-is-itself}
        \item \( \mleft( x \ast y \mright)^{-1} = y^{-1} \ast x^{-1} \) \label{enum:group-inverse-of-product}
        \item \( \mleft( x^{-1} = y^{-1} \mright) {\implies} \mleft( x = y \mright) \) \label{enum:group-equal-inverses-implies-equal-elements}
        \item \( \mleft(x \ast x = x\mright) \iff \mleft( x = e \mright) \) \label{enum:group-idempotent-element-iff-element-equals-identity}
    \end{enumerate}
\end{thm}

\begin{proof}
    Let \( \tuple{G,\ast} \) be a group and \( e \) be the identity of \( G \) with respect to \( \ast \).
    Let \( x \) be an arbitrary element of \( G \) with the unique inverse \( x^{-1} \). By \cref{dfn:inverse-element-of-binary-operations},
    it follows that \( x^{-1} \) is both a left and right inverse of \( x \). In other words, \( x \ast x^{-1} = x^{-1} \ast x = e \). It follows
    from \cref{dfn:inverse-element-of-binary-operations} once again that \( x \) is both the left and right inverse of \( x^{-1} \). Therefore,
    \( \mleft( x^{-1} \mright)^{-1} = x \).

    \vspace{0.3cm}
    \noindent
    To prove property \ref{enum:group-inverse-of-product}, we note the fact that \( \function{\ast}{G^{2}}{G} \) is associative and also the fact that
    \( G \) has an identity with respect to \( \ast \). In other words, \cref{thm:inverse-of-x-star-y} applies and \( \mleft( x \ast y \mright)^{-1} = y^{-1} \ast x^{-1} \)
    for any \( x, y \in G \).

    \vspace{0.3cm}
    \noindent
    To prove property \ref{enum:group-equal-inverses-implies-equal-elements}, suppose that
    \( x^{-1} = y^{-1} \) for some arbitrary \( x,y \in G \). This means that there exists
    some \( z \in G \) such that \( z \) is the inverse of both \( x \) and \( y \) with
    respect to \( \ast \). By \cref{dfn:inverse-element-of-binary-operations}, it follows that,
    \begin{gather*}
        x \ast z = z \ast x = e\\
        y \ast z = z \ast y = e
    \end{gather*}
    Since \( z \in G \) and \( \tuple{G,\ast} \) forms a group,
    it follows from the \aref{dfn:group}{group axioms} that \( z \) has a unique inverse \( z^{-1} \) with respect to \( \ast \).
    Therefore,
    \begin{gather*}
        \mleft( x \ast z \mright) \ast z^{-1} = e \ast z^{-1} = z^{-1}\\
        \mleft( y \ast z \mright) \ast z^{-1} = e \ast z^{-1} = z^{-1}
    \end{gather*}
    The associativity of \( \ast \) then allows us to simplify the left-hand side of both equations:
    \begin{gather*}
        \mleft( x \ast z \mright) \ast z^{-1} = x \ast \mleft( z \ast z^{-1} \mright) = x \ast e = x\\
        \mleft( y \ast z \mright) \ast z^{-1} = y \ast \mleft( z \ast z^{-1} \mright) = y \ast e = y
    \end{gather*}
    This allows us to conclude that \( x = z^{-1} = y \).

    \vspace{0.3cm}
    \noindent
    To prove property \ref{enum:group-idempotent-element-iff-element-equals-identity}, we let \( x \) be an arbitrary
    element of \( G \) such that \( x \ast x = x \). It follows from the \aref{dfn:group}{group axioms} and \cref{thm:group-uniqueness-of-inverses}
    that \( x \) has a unique inverse \( x^{-1} \) with respect to \( \ast \). Since \( x \ast x = x \), it follows
    that \( \mleft( x \ast x \mright) \ast x^{-1} = x \ast x^{-1} = e \). But the \aref{dfn:group}{group axioms} state that \( \ast \) is associative.
    Thus, \( \mleft( x \ast x \mright) \ast x^{-1} = x \ast \mleft( x \ast x^{-1} \mright) = x \ast e = x \). Clearly, \( x = e \). The converse follows
    immediately from \cref{dfn:identity-element-of-binary-operations}. If \( x = e \), then obviously \( x \ast x = x \).
\end{proof}

\subsubsection{Rings}
\begin{dfn}
    \phantomsection
    \label{dfn:ring}
    Let \( R \) be a non-empty set. Let \function{+}{R^{2}}{R} and \function{\ast}{R^{2}}{R} be binary operations on \( R \).
    The structure \( \tuple{R, +, \ast} \) is a \keyword{ring} if and only if the following first-order sentences are satisfied:
    \begin{enumerate}[topsep=2pt,noitemsep,label=\textup{\texttt{(\alph*)}}]
        \item \( \tuple{R, +} \) is an abelian group.
            \begin{itemize}
                \setlength{\itemsep}{0.1em}
                \item \( \forall x\,{\in}\,R \, \forall y\,{\in}\,R \, \forall z\,{\in}\,R \, \big[ \mleft( x + y \mright) + z = x + \mleft( y + z \mright) \big] \)
                \item \( \exists z\,{\in}\,R \, \forall x\,{\in}\,R \, \big[ \mleft( x + z = z + x = x \mright)
                       \,\,{\land}\,\,
                       \exists y\,{\in}\,R \, \mleft( x + y = y + x = z \mright) \big] \)
                \item \( \forall x\,{\in}\,R \, \forall y\,{\in}\,R \, \big( x + y = y + x \big) \)
            \end{itemize} \label{enum:ring-R-plus-forms-abelian-group}
        \item \( \ast \) is associative.
            \begin{itemize}
                \item \( \forall x\,{\in}\,R \, \forall y\,{\in}\,R \, \forall z\,{\in}\,R \, \big[ \mleft( x \ast y \mright) \ast z = x \ast \mleft( y \ast z \mright) \big] \)
            \end{itemize} \label{enum:ring-star-is-associative}
        \item \( \ast \) is distributive over \( + \).
            \begin{itemize}
                \setlength{\itemsep}{0.1em}
                \item \( \forall x\,{\in}\,R \, \forall y\,{\in}\,R \, \forall z\,{\in}\,R \,
                         \big[ x \ast \mleft( y + z \mright) = \mleft( x \ast y \mright) + \mleft( x \ast z \mright) \big]\)
                \item \( \forall x\,{\in}\,R \, \forall y\,{\in}\,R \, \forall z\,{\in}\,R \,
                         \big[ \mleft( x + y \mright) \ast z = \mleft( x \ast z \mright) + \mleft( y \ast z \mright) \big]\)
            \end{itemize} \label{enum:ring-star-distributes-over-plus}
    \end{enumerate}
    The sentences \crefrange{enum:ring-R-plus-forms-abelian-group}{enum:ring-star-distributes-over-plus} are collectively known as the \keyword{ring axioms}.
\end{dfn}

\noindent
There are no requirements in the ring axioms that an identity of R with respect to \( \ast \) exists. In cases where an identity
with respect to \( \ast \) \emph{does} exist, then the ring is usually called a \keyword{ring with identity}.

\begin{thm}
    \phantomsection
    \label{thm:ring-uniqueness-of-identity}
    If \( \tuple{R,+,\ast} \) is a ring with identity, then the identity of \( R \) with respect to \( \ast \) is unique.
\end{thm}

\begin{proof}
    Since \( \function{\ast}{R^{2}}{R} \) is a total function on \( R^{2} \), the theorem follows immediately from
    \cref{cor:uniqueness-of-identity-element-for-total-operations}.
\end{proof}

\begin{thm}
    \label{thm:ring-uniqueness-of-star-inverse}
    If \( \tuple{R,+,\ast} \) is a ring with identity and \( x \in R \) has an inverse with respect to \( \ast \), then the inverse is unique.
\end{thm}

\begin{proof}
    Since \( \tuple{R,+,\ast} \) is a ring, it follows from the ring axioms that \( \ast \) is associative. Moreover, \( R \) has an identity
    with respect to \( \ast \), so \cref{thm:uniqueness-of-inverse-element-for-total-associative-binary-operations} applies and every element of
    \( R \) has at most one inverse.
\end{proof}

\begin{notation}
    \label{notation:ring-denote-identities-and-inverses-using-arithmetic-symbols}
    Rings are a generalization of the property of integers.
    As such, it is customary to denote the identites of a ring using familiar arithmetic symbols.
    Thus, for any ring \( \tuple{R, +, \ast} \) and \( x \in R \):
    \begin{enumerate}[topsep=2pt,noitemsep,label=\textup{\texttt{(\alph*)}}]
        \item \( 0 \) denotes the identity of \( R \) with respect to \( + \).
        \item \( 1 \) denotes the identity of \( R \) with respect to \( \ast \).
        \item \( -x \) denotes the inverse of \( x \) with respect to \( + \).
        \item \( x^{-1} \) denotes the inverse of \( x \) with respect to \( \ast \).
    \end{enumerate}
\end{notation}

\begin{rem}
    One should never assume that \( 1 \neq 0 \). Even though they are different symbols, \( 1 \) and \( 0 \) could very well refer to
    the same element of \( R \).
\end{rem}

\begin{dfn}[Commutative Ring]
    \label{dfn:ring-commutative}
    A ring \( \tuple{R,+,\ast} \) is \keyword{commutative} if and only if \( \ast \) is commutative; that is,
    \( \tuple{R,+,\ast} \) satisfies the following first-order sentence:
    \begin{equation*}
        \forall x\,{\in}\,R \, \forall y\,{\in}\,R \, \big[ x \ast y = y \ast x \big]
    \end{equation*}
\end{dfn}

\begin{dfn}[Division Ring]
    \label{dfn:ring-division}
    A ring \( \tuple{R,+,\ast} \) is a \keyword{division ring} if and only if the following conditions hold:
    \begin{enumerate}[topsep=2pt,noitemsep,label=\textup{\texttt{(\alph*)}}]
        \item \( R \) has an identity \( 1 \) with respect to \( \ast \).
        \item \( 1 \neq 0 \).
        \item \( \forall x\,{\in}\,R \, \big[ \mleft( x \neq 0 \mright) \implies \exists y\,{\in}\,R \, \mleft( x \ast y = y \ast x = 1 \mright) \big] \).
    \end{enumerate}
\end{dfn}

\begin{dfn}[Boolean Ring]
    \label{dfn:ring-boolean}
    A ring \( \tuple{R,+,\ast} \) is \keyword{boolean} if and only if is also satisfies the following first-order sentence:
    \(  \forall x\,{\in}\,R \, \big[ x \ast x = x \big] \).
\end{dfn}

\begin{rem}
    If \( \tuple{R,+,\ast} \) forms a boolean ring, then \( R \) consists only of elements that are
    \aref{dfn:idempotence-of-element-under-binary-operations}{idempotent} under \( \ast \).
\end{rem}

\begin{thm}
    \label{thm:ring-basic-properties}
    If \( \tuple{R,+,\ast} \) is a ring, then for all \( x,y \in R \):
    \begin{enumerate}[topsep=2pt,noitemsep,label=\textup{\texttt{(\alph*)}}]
        \item \( x \ast 0 = 0 \ast x \) \label{enum:ring-product-with-zero}
        \item \( \mleft( -x \mright) \ast y = x \ast \mleft( -y \mright) = - \mleft( x \ast y \mright) \)
                \label{enum:ring-single-negative-sign-product}
        \item \( \mleft( -x \mright) \ast \mleft( -y \mright) = x \ast y \)
                \label{enum:ring-double-negative-signs-product}
        \item \( -x = \mleft( -1 \mright) \ast x \)
                \label{enum:ring-negative-element-is-negative-one-times-element}
    \end{enumerate}
\end{thm}

\begin{proof}
    Conforming to \cref{notation:ring-denote-identities-and-inverses-using-arithmetic-symbols},
    we let \( 0 \) denote the identity of \( R \) with respect to \( + \) and \( -x \)
    denote the inverse of \( x \) with respect to \( + \) for any \( x \in R \).

    \vspace{0.3cm}
    \noindent
    \textsc{Property} \ref{enum:ring-product-with-zero}:
    Let \( x \in R \) be arbitrary. Since \( 0 \) is the identity of \( R \) with respect to \( + \),
    it follows from \cref{dfn:identity-element-of-binary-operations} that \( 0 + 0 = 0 \). Clearly,
    \begin{gather*}
        x \ast 0 = x \ast \mleft( 0 + 0 \mright)\\
        0 \ast x = \mleft( 0 + 0 \mright) \ast x
    \end{gather*}
    By the distributivity of \( \ast \) over \( + \), we have,
    \begin{gather}
        x \ast 0 = x \ast 0 + x \ast 0 \label{eq:x-star-0-distribute} \\
        0 \ast x = 0 \ast x + 0 \ast x \label{eq:0-star-x-distribute}
    \end{gather}
    From the \aref{dfn:ring}{ring axioms}, we know that \( \tuple{R, +} \) is a group, so it follows from the \aref{dfn:group}{group axioms} that there exist
    some \( y,z \in R \) such that \( y \) is the inverse of \( x \ast 0 \) with respect to \( + \) and \( z \) is the inverse
    of \( 0 \ast x \) with respect to \( + \). From \eqref{eq:x-star-0-distribute} and \eqref{eq:0-star-x-distribute}, we have,
    \begin{gather*}
        x \ast 0 + y = \mleft( x \ast 0 + x \ast 0 \mright) + y\\
        0 \ast x + z = \mleft( 0 \ast x + 0 \ast x \mright) + z
    \end{gather*}
    But \( + \) is associative (because \( \tuple{R,+} \) is a group), so these equations become,
    \begin{gather*}
        x \ast 0 + y = x \ast 0 + \mleft( x \ast 0 + y \mright)\\
        0 \ast x + z = 0 \ast x + \mleft( 0 \ast x + z \mright)
    \end{gather*}
    But \( y \) and \( z \) are the inverse of \( x \ast 0 \) and \( 0 \ast x \) respectively, so necessarily \( x \ast 0 + y = 0 \)
    and \( 0 \ast x + z = 0 \). Therefore,
    \begin{gather*}
        x \ast 0 + \mleft( x \ast 0 + y \mright) = x \ast 0 + 0 = x \ast 0 = 0\\
        0 \ast x + \mleft( 0 \ast x + z \mright) = 0 \ast x + 0 = 0 \ast x = 0
    \end{gather*}
    \noindent
    \textsc{Property} \ref{enum:ring-single-negative-sign-product}:
    Let \( x,y \in R \) be arbitrary. Since \( \tuple{R,+} \) is a group, \( x \) and \( y \) have unique inverses with respect to \( R \),
    denoted as \( -x \) and \( -y \). It follows from the distributivity of \( \ast \) over \( + \) and property \ref{enum:ring-product-with-zero} that,
    \begin{gather*}
        \mleft( -x \ast y \mright) + \mleft( x \ast y \mright) = \mleft( -x + x \mright) \ast y = 0 \ast y = 0 \\
        \mleft( x \ast y \mright) + \mleft( -x \ast y \mright) = \mleft( x + -x \mright) \ast y = 0 \ast y = 0 \\
        \mleft( x \ast -y \mright) + \mleft( x \ast y \mright) = x \ast \mleft( -y + y \mright) = x \ast 0 = 0 \\
        \mleft( x \ast y \mright) + \mleft( x \ast -y \mright) = x \ast \mleft( y + -y \mright) = x \ast 0 = 0
    \end{gather*}
    By \cref{dfn:inverse-element-of-binary-operations}, this implies that \( \mleft( -x \ast y \mright) \) and \( \mleft( x \ast -y \mright) \) are both the inverses of
    \( x \ast y \) with respect to \( + \). But \( \tuple{R, +} \) is a group, so \cref{thm:group-uniqueness-of-inverses} applies and the inverse of every element of \( R \)
    with respect to \( R \) must be unique. Therefore, \( \mleft( -x \ast y \mright) = \mleft( x \ast -y \mright) = - \mleft( x \ast y \mright) \).

    \vspace{0.3cm}
    \noindent
    \textsc{Property} \ref{enum:ring-double-negative-signs-product}:
    Let \( x,y \in R \) be arbitrary. By properties, \ref{enum:ring-product-with-zero}, \ref{enum:ring-single-negative-sign-product},
    and also the distributivity of \( \ast \) over \( + \), it follows that
    \( \mleft[ \mleft( -x \mright) \ast \mleft( -y \mright) \mright] + \mleft[ - \mleft( x \ast y \mright) \mright]
       = \mleft[ \mleft( -x \mright) \ast \mleft( -y \mright) \mright] + \mleft[ \mleft( -x \mright) \ast y \mright]
       = \mleft( -x \mright) \ast \mleft( -y + y \mright) = -x \ast 0 = 0 \).
    Similarly, \( -\mleft( x \ast y \mright) + \mleft[ \mleft( -x \mright) \ast \mleft( -y \mright) \mright]
        = \mleft[ x \ast \mleft( -y \mright) \mright] + \mleft[ \mleft( -x \mright) \ast \mleft( -y \mright) \mright]
        = \mleft[ x + \mleft( -x \mright) \mright] \ast \mleft( -y \mright)
        = 0 \ast \mleft( -y \mright)
        = 0
    \).
    By \cref{dfn:inverse-element-of-binary-operations} and \cref{thm:group-basic-properties}, it follows that
    \( \mleft( -x \mright) \ast \mleft( -y \mright) = - \mleft[ - \mleft( x \ast y \mright) \mright] = x \ast y \).

    \vspace{0.3cm}
    \noindent
    \textsc{Property} \ref{enum:ring-negative-element-is-negative-one-times-element}:
    Suppose that \( R \) has an identity with respect to \( \ast \).
    We will follow the convention introduced in \cref{notation:ring-denote-identities-and-inverses-using-arithmetic-symbols}
    and denote the identity of \( R \) with respect to \( \ast \) as \( 1 \).
    Let \( x \in R \) be arbitrary. Since \( 1 \) is the identity of \( R \) with respect to \( \ast \), it follows
    from \cref{dfn:identity-element-of-binary-operations} that \( x = 1 \ast x = x \ast 1 \).
    Thus,
    \begin{align*}
        \mleft( -1 \ast x \mright) + x
        &= \mleft( -1 \ast x \mright) + \mleft( 1 \ast x \mright)
        = \mleft( -1 + 1 \mright) \ast x
        = 0 \ast x
        = 0\\
        x + \mleft( -1 \ast x \mright)
        &= \mleft( 1 \ast x \mright) + \mleft( -1 \ast x \mright)
        = \mleft( 1 + \mleft( -1 \mright) \mright) \ast x
        = 0 \ast x
        = 0
    \end{align*}
    By \cref{dfn:inverse-element-of-binary-operations}, \( \mleft( -1 \ast x \mright) \) is a left and right inverse of \( x \) with respect to \( + \).
    In other words, \( \mleft( -1 \ast x \mright) = -x \).
\end{proof}

\begin{thm}[Generalized Distributivity]
    \phantomsection \label{thm:ring-generalized-distributive-law}
    If \( \tuple{R, +, \ast} \) is a ring, then for any \( k \in \naturalset \) where \( k > 1 \),
    \begin{equation*}
        \forall x\,{\in}\,R \, \forall y_{0}\,{\in}\,R \, \dots \forall y_{k-1}\,{\in}\,R \,
        \big[ x \ast \mleft( y_{0} + \cdots y_{k-1} \mright)
        = x \ast y_{0} + \cdots + x \ast y_{k-1}
        \big]
    \end{equation*}
\end{thm}

\begin{thm}
    \label{thm:ring-boolean-ring-additive-inverse-of-an-element-is-the-element-itself}
    If \( \tuple{R,+,\ast} \) is a boolean ring, then \( \forall x\,{\in}\,R \, \mleft[ x = -x \mright] \).
\end{thm}

\begin{proof}
    Let \( x \in R \) be arbitrary. Since \( \tuple{R,+,\ast} \) is a boolean ring, we have by \cref{dfn:ring-boolean}
    \( \mleft( -x \mright) \ast \mleft( -x \mright) = -x \). But \cref{thm:ring-basic-properties} implies that
    \( \mleft( -x \mright) \ast \mleft( -x \mright) = x \ast x = x \). Therefore, \( x = -x \).
\end{proof}

\begin{thm}
    \label{thm:ring-commutativity-of-boolean-rings}
    Every boolean ring is commutative.
\end{thm}

\begin{proof}
    Let \( \tuple{R,+,\ast} \) be an arbitrary boolean ring. Let \( x \) and \( y \) be arbitrary elements of \( R \).
    Now consider the operation \( \mleft( x + y \mright) \ast \mleft( x + y \mright) \). We know that \( \tuple{R,+,\ast} \)
    is a boolean ring, so the following must be true by \cref{dfn:ring-boolean} and the distributivity of \( \ast \) over \( + \):
    \begin{gather*}
        \mleft( x + y \mright) \ast \mleft( x + y \mright) = \mleft( x + y \mright)\\
        \mleft( x + y \mright) \ast \mleft( x + y \mright) = x \ast \mleft( x + y \mright) + y \ast \mleft( x + y \mright)
    \end{gather*}
    Combining the two equations, it follows that \( \mleft( x + y \mright) = x \ast \mleft( x + y \mright) + y \ast \mleft( x + y \mright) \).
    Applying the distributivity of \( \ast \) over \( + \) once again, we have:
    \begin{equation*}
        \mleft( x + y \mright) = \big[ x \ast x + x \ast y \big] + \big[ y \ast x + y \ast y \big] = \mleft[ x + x \ast y \mright] + \mleft[ y \ast x + y \mright]
    \end{equation*}
    We know that \( \tuple{R,+} \) is an abelian group, so \( + \) is associative and commutative. The right-hand side of the equation can therefore be rearranged to
    \( \mleft[ x + y \mright] + \mleft[ x \ast y + y \ast x \mright] \). Thus,
    \begin{equation*}
        \mleft( x + y \mright) = \mleft( x + y \mright) + \mleft( x \ast y + y \ast x \mright)
    \end{equation*}
    By \cref{dfn:identity-element-of-binary-operations}, it follows that \( \mleft( x \ast y + y \ast x \mright) \) is the identity of \( R \) with respect to \( + \),
    which is unique by \cref{thm:ring-uniqueness-of-identity}. Using the convention introduced in \cref{notation:ring-denote-identities-and-inverses-using-arithmetic-symbols},
    we can express this fact as:
    \begin{equation*}
        \mleft( x \ast y + y \ast x \mright) = 0
    \end{equation*}
    By \cref{dfn:left-and-right-inverses-of-binary-operations,dfn:inverse-element-of-binary-operations},
    it follows that \( x \ast y \) and \( y \ast x \) are inverses of each other with respect to \( + \).
    That is,
    \begin{equation*}
        y \ast x = - \mleft( x \ast y \mright)
    \end{equation*}
    But \( \tuple{R,+,\ast} \) is a boolean ring, so
    \( x \ast y = - \mleft( x \ast y \mright) \) by \cref{thm:ring-boolean-ring-additive-inverse-of-an-element-is-the-element-itself}.
    Therefore, necessarily \( x \ast y = - \mleft( x \ast y \mright) = y \ast x \). Since \( x \) and \( y \) are arbitrary, the theorem
    follows from \aref{rule-of-inference:universal-generalization}{Universal Generalization}.
\end{proof}

\begin{thm}
    \label{thm:ring-power-set-of-a-set-forms-a-boolean-ring}
    Let \( A \) be a non-empty set where \( \function{+}{\powerset{A}^{2}}{\powerset{A}} \)
    and \( \function{\ast}{\powerset{A}^{2}}{\powerset{A}} \) are binary operations on \( \powerset{A} \)
    with the following definitions:
    \begin{gather*}
        +
        =
        \setbuilderBig{\orderedpair{\tuple{x,y}}{z}}{x \in \powerset{A} \,\,{\land}\,\, y \in \powerset{A} \,\,{\land}\,\, z = x \symdiff y }\\
        \ast
        =
        \setbuilderBig{\orderedpair{\tuple{x,y}}{z}}{x \in \powerset{A} \,\,{\land}\,\, y \in \powerset{A} \,\,{\land}\,\, z = x \intersect y }
    \end{gather*}
    The structure \( \tuplebig{\powerset{A},+,\ast} \) is a boolean ring.
\end{thm}

\begin{proof}
    According to \cref{dfn:ring,dfn:ring-boolean}, we have to prove each of the following claims:
    \begin{enumerate*}[label=(\roman*)]
        \item \( \tuplebig{\powerset{A},+} \) is an abelian group, \label{enum:ring-power-set-forms-a-boolean-ring-claim1}
        \item \( \ast \) is associative, \label{enum:ring-power-set-forms-a-boolean-ring-claim2}
        \item \( \ast \) is distributive over \( + \), \label{enum:ring-power-set-forms-a-boolean-ring-claim3}
        \item \( \forall x\,{\in}\,\powerset{A} \, \big[ x \ast x = x \big]\). \label{enum:ring-power-set-forms-a-boolean-ring-claim4}
    \end{enumerate*}

    \vspace{0.5cm}
    \noindent
    \textsc{Claim} \ref{enum:ring-power-set-forms-a-boolean-ring-claim1}:
    The goal is to prove that the structure \( \tuplebig{\powerset{A},+} \) satisfies the \aref{dfn:group}{group axioms}.
    Let \( x, y, z \in \powerset{A} \) be arbitrary. Thus, \( x + \mleft( y + z \mright) = x \symdiff \mleft( y \symdiff z \mright) \)
    and \( \mleft( x + y \mright) + z = \mleft( x \symdiff y \mright) \symdiff z \) by definition.
    But \( x \symdiff \mleft( y \symdiff z \mright) = \mleft( x \symdiff y \mright) \symdiff z \) by \cref{set-identities:associativity-of-symmetric-difference}.
    Therefore \( x + \mleft( y + z \mright) = \mleft( x + y \mright) + z \); in other words, \( + \) is associative.
    By \cref{set-identities:commutativity-of-symmetric-difference}, we also have \( x \symdiff y = y \symdiff x \) for any \( x,y \in \powerset{A} \), so
    \( + \) is necessarily commutative.
    Next, we note that \( \emptyset \in \powerset{A} \) and,
    \begin{gather*}
        x \symdiff \emptyset = \mleft( x \union \emptyset \mright) - \mleft( x \intersect \emptyset \mright) = x - \emptyset = x\\
        \emptyset \symdiff x = \mleft( \emptyset \union x \mright) - \mleft( \emptyset \intersect x \mright) = x - \emptyset = x
    \end{gather*}
    In other words, \( \emptyset \) is both the left and right identities of \( \powerset{A} \) with respect to \( + \).
    Thus, \( \emptyset \) is the identity of \( \powerset{A} \) by \cref{dfn:identity-element-of-binary-operations}.
    Now let \( x \in \powerset{A} \) be arbitrary.
    Clearly, \( x + x = x \symdiff x = \mleft( x \union x \mright) - \mleft( x \intersect x \mright) = x - x = \emptyset \).
    Since \( \emptyset \) is the identity of \( \powerset{A} \) with respect to \( + \), it follows from \cref{dfn:inverse-element-of-binary-operations}
    that \( x \) is the inverse of \( x \) with respect to \( + \) for any \( x \in \powerset{A} \); in other words, every \( x \in \powerset{A} \)
    has an inverse with respect to \( + \). Having established that \( + \) is associative, commutative, \( \powerset{A} \) has an identity with
    respect to \( + \), and that every element of \( \powerset{A} \) has an inverse with respect to \( + \), we conclude that \( \tuplebig{\powerset{A},+} \)
    does indeed form an abelian group.

    \vspace{0.5cm}
    \noindent
    \textsc{Claim} \ref{enum:ring-power-set-forms-a-boolean-ring-claim2}:
    Let \( x, y, z \in \powerset{A} \) be arbitrary.
    By definition,
    \( x \ast \mleft( y \ast z \mright) = x \intersect \mleft( y \intersect z \mright) \). But by \cref{set-identities:associativity-of-intersection},
    \( x \intersect \mleft( y \intersect z \mright) = \mleft( x \intersect y \mright) \intersect z = \mleft( x \ast y \mright) \ast z \). Therefore,
    \( x \ast \mleft( y \ast z \mright) = \mleft( x \ast y \mright) \ast z \). By \aref{rule-of-inference:universal-generalization}{Universal Generalization}, we conclude that \( \ast \) is indeed
    an associative operation.

    \vspace{0.5cm}
    \noindent
    \textsc{Claim} \ref{enum:ring-power-set-forms-a-boolean-ring-claim3}:
    This is a direct consequence of \cref{set-identities:distributivity-of-intersection-over-symmetric-difference}, which states that
    \( \intersect \) is distributive over \( \symdiff \).

    \vspace{0.5cm}
    \noindent
    \textsc{Claim} \ref{enum:ring-power-set-forms-a-boolean-ring-claim4}:
    This follows immediately form the definition of \( \ast \) since \( x \ast x = x \intersect x = x \) for any \( x \in \powerset{A} \).

    \vspace{0.5cm}
    \noindent
    Having proven claims \crefrange{enum:ring-power-set-forms-a-boolean-ring-claim1}{enum:ring-power-set-forms-a-boolean-ring-claim4}, we
    conclude that the structure \( \tuplebig{\powerset{A},+,\ast} \) is indeed a boolean ring.
\end{proof}

\begin{thm}
    \label{thm:ring-system-of-sets-forms-a-ring-iff-it-contains-the-empty-set-and-is-closed-under-union-and-set-diff}
    If \( S \) is a non-empty system of sets, then \( \tuple{S, \symdiff, \intersect} \) is a ring if and only if the following conditions hold:
    \begin{enumerate}[topsep=2pt,noitemsep,label=\textup{\texttt{(\alph*)}}]
        \item \( \emptyset \in S \) \label{enum:ring-empty-set-is-in-the-system-of-sets}
        \item \( S \) is closed under set difference
            \begin{itemize}
                \item \( \forall x\,{\in}\,S \, \forall y\,{\in}\,S \, \big[ x - y \in S \big] \)
            \end{itemize} \label{enum:ring-closure-under-set-diff}
        \item \( S \) is closed under binary union
            \begin{itemize}
                \item \( \forall x\,{\in}\,S \, \forall y\,{\in}\,S \, \big[ x \union y \in S \big] \)
            \end{itemize} \label{enum:ring-closure-under-binary-union}
    \end{enumerate}
\end{thm}

\begin{proof}
    Let \( S \) be a non-empty system of sets such that \( \tuple{S, \symdiff, \intersect} \) forms a ring.
    Since \( \tuple{S, \symdiff, \intersect} \) is a ring, it is thus required by the \aref{dfn:ring}{ring axioms} that
    \( \symdiff \) and \( \intersect \) be \emph{total} functions on \( S^{2} \); that is, \( S \) is closed under \( \symdiff \) and \( \intersect \).
    Let \( x \in S \) be arbitrary.
    By the \aref{set-identities:symmetric-difference-of-any-set-with-itself-is-empty}{nilpotence of \( \symdiff \)},
    it follows that \( x \symdiff x = \emptyset \). But \( S \) is closed under \( \symdiff \), so necessarily \( \emptyset = x \symdiff x \in S \), proving
    \ref{enum:ring-empty-set-is-in-the-system-of-sets}.

    \vspace{0.3cm}
    \noindent
    Next, we let \( x, y \in S \) be arbitrary.
    Using \cref{set-identities:union-in-terms-of-symmetric-difference-and-intersection,set-identities:set-difference-in-terms-of-symmetric-difference-and-intersection},
    we can express \( x \union y \) and \( x - y \) as follows:
    \begin{align*}
        x \union y &= \mleft( x \symdiff y \mright) \symdiff \mleft( x \intersect y \mright)\\
        x - y &= x \symdiff \mleft( x \intersect y \mright)
    \end{align*}
    Since \( x \) and \( y \) are both elements of \( S \) and \( S \) is closed under \( \symdiff \) and \( \intersect \), it follows that \( x \symdiff y \in S \)
    and \( x \intersect y \in S \).
    Therefore, \( x \union y = \mleft( x \symdiff y \mright) \symdiff \mleft( x \intersect y \mright) \) and \( x - y = x \symdiff \mleft( x \intersect y \mright) \)
    must also be elements of \( S \), which proves \ref{enum:ring-closure-under-set-diff} and \ref{enum:ring-closure-under-binary-union}. In other words, the fact
    that \( \tuple{S, \symdiff, \intersect} \) forms a ring implies \crefrange{enum:ring-empty-set-is-in-the-system-of-sets}{enum:ring-closure-under-binary-union}.

    \vspace{0.3cm}
    \noindent
    To prove the converse, we start with the hypothesis that \( S \) is a non-empty system of sets such that
    \crefrange{enum:ring-empty-set-is-in-the-system-of-sets}{enum:ring-closure-under-binary-union} holds.
    Then we proceed to show that \( \tuple{S, \symdiff, \intersect} \) does indeed satisfy the \aref{dfn:ring}{ring axioms}:
    \begin{enumerate*}[label=(\roman*)]
        \item \( S \) is closed under \( \symdiff \) and \( \intersect \), \label{enum:S-is-closed-under-symdiff-and-intersection}
        \item \( \tuple{S, \symdiff} \) forms an \aref{dfn:abelian-group}{abelian group}, \label{enum:S-symdiff-forms-abelian-group}
        \item \( \intersect \) is associative, \label{enum:intersect-is-associative}
        \item \( \intersect \) is distributive over \( \symdiff \). \label{enum:intersect-distributes-over-symdiff}
    \end{enumerate*}

    \vspace{0.3cm}
    \noindent
    \textsc{Claim} \ref{enum:S-is-closed-under-symdiff-and-intersection}: \( S \) is closed under \( \symdiff \) and \( \intersect \).

    \vspace{0.1cm}
    \noindent
    Let \( x, y \in S \) be arbitrary. By \cref{set-identities:symmetric-difference-as-the-union-of-differences},
    we have \( x \symdiff y = \mleft( x - y \mright) \union \mleft( y - x \mright) \). Since \( x, y \in S \), it
    follows from \ref{enum:ring-closure-under-set-diff} that \( x - y \in S \) and \( y - x \in S \),
    which in turn imply \( \mleft( x - y \mright) \union \mleft( y - x \mright) \in S \) due to \ref{enum:ring-closure-under-binary-union}.
    In other words, \( S \) is closed under \( \symdiff \).
    To prove that \( S \) is closed under \( \intersect \), observe that \( x \intersect y = \mleft( x \union y \mright) - \mleft( x \symdiff y \mright) \).
    Since \( S \) is closed under \( \union \) by \ref{enum:ring-closure-under-binary-union}, it follows that \( x \union y \in S \). It has already been proven
    above that \( S \) is also closed under \( \symdiff \), so \( x \symdiff y \in S \) is also true. Moreover, \ref{enum:ring-closure-under-set-diff} states that
    \( S \) is closed under \( - \), so \( x \intersect y = \mleft( x \union y \mright) - \mleft( x \symdiff y \mright) \in S \) is necessarily true.

    \vspace{0.3cm}
    \noindent
    \textsc{Claim} \ref{enum:S-symdiff-forms-abelian-group}: \( \tuple{S, \symdiff} \) forms an \aref{dfn:abelian-group}{abelian group}.

    \vspace{0.1cm}
    \noindent
    First, we must show that \( \tuple{S, \symdiff} \) is a \aref{dfn:group}{group} and that \( \symdiff \) is commutative. The associativity
    and commutativity of \( \symdiff \) follows trivially from
    \cref{set-identities:associativity-of-symmetric-difference,set-identities:commutativity-of-symmetric-difference}.
    Also, \( \emptyset \in S \) by assumption \ref{enum:ring-empty-set-is-in-the-system-of-sets} and
    \( x \symdiff \emptyset = \emptyset \symdiff x = x \) for any \( x \in S \), so \( \emptyset \) is the identity
    of \( S \) with respect to \( \symdiff \). Moreover, \( x \symdiff x = \emptyset \) for any \( x \in S \), so every element of \( S \) has an
    inverse with respect to \( \symdiff \). Taking all of these facts into account, it follows that \( \tuple{S, \symdiff} \) is indeed an
    \aref{dfn:abelian-group}{abelian group}.

    \vspace{0.5cm}
    \noindent
    \textsc{Claims} \ref{enum:intersect-is-associative} and \ref{enum:intersect-distributes-over-symdiff}:
    \( \intersect \) is associative and \( \intersect \) distributes over \( \symdiff \).

    \vspace{0.1cm}
    \noindent
    The associativity of \( \intersect \) follows readily from \cref{set-identities:associativity-of-intersection}
    whereas the distributivity of \( \intersect \) over \( \symdiff \) has already been proven in
    \cref{set-identities:distributivity-of-intersection-over-symmetric-difference}.

    \vspace{0.3cm}
    \noindent
    Having established all four claims, we can now conclude that if \crefrange{enum:ring-empty-set-is-in-the-system-of-sets}{enum:ring-closure-under-binary-union}
    hold, then \( \tuple{S, \symdiff, \intersect} \) must form a ring.
\end{proof}

\noindent
\cref{thm:ring-system-of-sets-forms-a-ring-iff-it-contains-the-empty-set-and-is-closed-under-union-and-set-diff} shows that
there is an equivalent formulation for the \aref{dfn:ring}{ring axioms} if the universe of the structure is a system of sets,
which leads us to the following definition.

\begin{dfn}[Ring of Sets]
    \label{dfn:ring-rings-of-sets}
    A structure \( \tuple{S, \symdiff, \intersect} \) is a \keyword{ring of sets} if and only if the following conditions hold:
    \begin{enumerate}[topsep=2pt,noitemsep,label=\textup{\texttt{(\alph*)}}]
        \setlength{\itemsep}{0.1em}
        \item \( \emptyset \in S \)
        \item \( \forall x\,{\in}\,S \, \forall y\,{\in}\,S \, \big[ x - y \in S \big] \)
        \item \( \forall x\,{\in}\,S \, \forall y\,{\in}\,S \, \big[ x \union y \in S \big] \)
    \end{enumerate}
\end{dfn}

\begin{rem}
    Since the structure \( \tuple{S,\symdiff, \intersect} \) forms a ring when the three conditions are met
    and since all elements of \( S \) are sets, the structure \( \tuple{S, \symdiff, \intersect} \) is aptly called a ring of sets.
\end{rem}

\subsubsection{Fields}
\begin{dfn}
    \label{dfn:field}
    Let \( F \) be a non-empty set and \( \function{+}{F^{2}}{F} \) and \( \function{\ast}{F^{2}}{F} \) be binary operations on \( F \).
    The structure \( \tuple{F, +, \ast} \) is a \keyword{field} if and only if the following first-order sentences are satisfied:
    \begin{enumerate}[topsep=2pt,noitemsep,label=\textup{\texttt{(\alph*)}}]
        \item \( + \) and \( \ast \) are associative \phantomsection \label{enum:field-plus-star-associative}
            \begin{itemize}
                \setlength{\itemsep}{0.2em}
                \item \( \forall x\,{\in}\,F \, \forall y\,{\in}\,F \, \forall z\,{\in}\,F \, \big[ x + \mleft( y + z \mright) = \mleft( x + y \mright) + z \big] \)
                \item \( \forall x\,{\in}\,F \, \forall y\,{\in}\,F \, \forall z\,{\in}\,F \, \big[ x \ast \mleft( y \ast z \mright) = \mleft( x \ast y \mright) \ast z \big] \)
            \end{itemize}
        \item \( + \) and \( \ast \) are commutative \phantomsection \label{enum:field-plus-star-commutative}
            \begin{itemize}
                \setlength{\itemsep}{0.2em}
                \item \( \forall x\,{\in}\,F \, \forall y\,{\in}\,F \, \big[ x + y = y + x \big] \)
                \item \( \forall x\,{\in}\,F \, \forall y\,{\in}\,F \, \big[ x \ast y = y \ast x \big] \)
            \end{itemize}
        \item There exist distinct elements \( 0 \) and \( 1 \) such that \( 0 \) is the identity of \( F \) with respect to \( + \)
              and \( 1 \) is the identity of \( F \) with respect to \( \ast \) \phantomsection \label{enum:field-plus-star-have-identities}
            \begin{itemize}
                \item \( \exists i\,{\in}\,F \, \exists j\,{\in}\,F \,
                    \big[
                        i \neq j
                        \,\,{\land}\,\,
                        \forall x\,{\in}\,F \, \big( x + i = i + x = x \big)
                        \,\,{\land}\,\,
                        \forall x\,{\in}\,F \, \big( x \ast j = j \ast x = x \big)
                    \big] \)
            \end{itemize}
        \item Every element of \( F \) has an inverse with respect to \( + \) \phantomsection \label{enum:field-inverse-with-respect-to-plus}
            \begin{itemize}
                \item \( \forall x\,{\in}\,F \, \exists y\,{\in}\,F \, \big[ x + y = y + x = 0 \big] \)
            \end{itemize}
        \item Every element of \( F \) that is not \( 0 \) has an inverse with respect to \( \ast \) \phantomsection \label{enum:field-inverse-with-respect-to-star-non-zero}
            \begin{itemize}
                \item \( \forall x\,{\in}\,F \, \big[ x \neq 0 {\implies} \exists y\,{\in}\,F \, \mleft( x \ast y = y \ast x = 1 \mright) \big] \)
            \end{itemize}
        \item \( \ast \) is distributive over \( + \) \phantomsection \label{enum:field-star-distributes-over-plus}
            \begin{itemize}
                \setlength{\itemsep}{0.2em}
                \item \( \forall x\,{\in}\,F \, \forall y\,{\in}\,F \, \forall z\,{\in}\,F \, \big[ x \ast \mleft( y + z \mright) = x \ast y + x \ast z \big] \)
                \item \( \forall x\,{\in}\,F \, \forall y\,{\in}\,F \, \forall z\,{\in}\,F \, \big[ \mleft( x + y \mright) \ast z = x \ast z + y \ast z \big] \)
            \end{itemize}
    \end{enumerate}
    These sentences \crefrange{enum:field-plus-star-associative}{enum:field-star-distributes-over-plus} are collectively known as the \keyword{field axioms}.
\end{dfn}

\begin{rem}
    The astute reader may have already noticed that a field is in fact a commutative, \aref{dfn:ring-division}{division ring}.
    This means that all the properties in \cref{thm:ring-basic-properties} apply to fields as well.
\end{rem}

\subsubsection{Vector Spaces} \phantomsection \label{sssec:vector-spaces}
\begin{dfn}
    \phantomsection
    \label{dfn:vector-space}
    Let \( \mathcal{F} = \tuple{F,+,\cdot} \) be a \aref{dfn:field}{field} and \( V \) be a non-empty set equipped with a binary operation
    \( \function{\oplus}{V^{2}}{V} \) and a binary function \( \function{\odot}{\prod \tuple{F,V}}{V} \). The structure
    \( \mathcal{V} = \tuple{V, \mathcal{F}, \oplus, \odot} \) or \( \mathcal{V} = \tuple{V,F,+,\cdot,\oplus,\odot} \) is a \keyword{vector space}
    over \( \mathcal{F} \) if and only if the following first-order sentences are satisfied:
    \begin{enumerate}[topsep=2pt,noitemsep,label=\textup{\texttt{(\alph*)}}]
        \item \( \tuple{V, \oplus} \) forms an \aref{dfn:group}{abelian group}.\phantomsection \label{enum:vector-space-V-and-addition-forms-abelian-group}
            \begin{itemize}
                \setlength{\itemsep}{0.1em}
                \item \( \forall \vectorvar{x}\,{\in}\,V \, \forall \vectorvar{y}\,{\in}\,V \, \forall \vectorvar{z}\,{\in}\,V \,
                    \big[ \vectorvar{x} \oplus \mleft( \vectorvar{y} \oplus \vectorvar{z} \mright)
                    = \mleft( \vectorvar{x} \oplus \vectorvar{y} \mright) \oplus \vectorvar{z} \big] \)
                \item \( \exists \vectorvar{z}\,{\in}\,V \, \forall \vectorvar{x}\,{\in}\,V \,
                    \big[ \mleft( \vectorvar{x} \oplus \vectorvar{z} = \vectorvar{z} \oplus \vectorvar{x} = \vectorvar{x} \mright)
                    \,\,{\land}\,\,
                    \exists \vectorvar{y}\,{\in}\,V \, \mleft( \vectorvar{x} \oplus \vectorvar{y} = \vectorvar{y} \oplus \vectorvar{x} = \vectorvar{z} \mright)
                    \big]\)
            \end{itemize}
        \item \( \forall x\,{\in}\,F \, \forall y\,{\in}\,F \, \forall \vectorvar{v}\,{\in}\,V \,
            \big[ x \cdot \mleft( y \odot \vectorvar{v} \mright) = \mleft( x \cdot y \mright) \odot \vectorvar{v} \big] \)
            \label{enum:vector-space-scalar-can-be-grouped-in-scalar-multiplication}
        \item \( \forall \vectorvar{v}\,{\in}\,V \, \big[ 1 \odot \vectorvar{v} = \vectorvar{v} \big] \)
            \label{enum:vector-space-scalar-multiplication-with-1-yields-the-original-vector}
        \item \( \odot \) distributes over \( \oplus \) and \( + \) \phantomsection \label{enum:vector-space-scalar-multiplication-distributes-over-vector-and-field-addition}
            \begin{itemize}
                \setlength{\itemsep}{0.1em}
                \item \( \forall x\,{\in}\,F \, \forall \vectorvar{u}\,{\in}\,V \, \forall \vectorvar{v}\,{\in}\,V \,
                    \big[ x \odot \mleft( \vectorvar{u} \oplus \vectorvar{v} \mright) = \mleft( x \odot \vectorvar{u} \mright)
                    \oplus \mleft( x \odot \vectorvar{v} \mright)\big] \)
                \item \( \forall x\,{\in}\,F \, \forall y\,{\in}\,F \, \forall \vectorvar{v}\,{\in}\,V \,
                    \big[ \mleft( x + y \mright) \odot \vectorvar{v} = \mleft( x \odot \vectorvar{v} \mright) + \mleft( y \odot \vectorvar{v} \mright)  \big]\)
            \end{itemize}
    \end{enumerate}
    The following terminology also applies:
    \begin{enumerate}[topsep=2pt,noitemsep,label=\textup{\texttt{(\arabic*)}}]
        \setlength{\itemsep}{0.1em}
        \item An element of \( F \) is called a \keyword{scalar}. \phantomsection \label{enum:vector-space-scalar}
        \item An element of \( V \) is called a \keyword{vector}. \phantomsection \label{enum:vector-space-vector}
        \item The binary operation \( \function{\oplus}{V^{2}}{V} \) is known as \keyword{vector addition}. \phantomsection \label{enum:vector-space-vector-addition}
        \item The binary function \( \function{\odot}{\prod \tuple{F,V}}{V} \) is known as \keyword{scalar multiplication}.
            \phantomsection \label{enum:vector-space-scalar-multiplication}
        \item For every \( \vectorvar{v} \in V \),\phantomsection \label{enum:vector-space-zero-vector}
            the unique \aref{dfn:identity-element-of-binary-operations}{identity} of \( V \) with respect to \( \oplus \)
            is known as the \keyword{zero vector} of \( \mathcal{V} \) and is denoted as \( \zerovector \).
        \item For every \( \vectorvar{v} \in V \), the unique \aref{dfn:inverse-element-of-binary-operations}{inverse} of
            \( \vectorvar{v} \) with respect to \( \oplus \) is called the \keyword{negative} of \( \vectorvar{v} \) and is denoted as \( -\vectorvar{v} \).
            \phantomsection \label{enum:vector-space-additive-inverse}
    \end{enumerate}
\end{dfn}

\begin{notation}
    \phantomsection \label{notation:simplified-notation-for-vector-spaces}
    To distinguish between vectors and scalars, variables and constants denoting vectors will be typeset using boldface fonts
    (\vectorvar{u},\vectorvar{v},\vectorvar{w}). This makes it possible to use the same symbols to denote different
    operations of \( \mathcal{V} \).
    For example, if the symbol `\( + \)' is used for both \( + \) and \( \oplus \), then there will be no confusion as to which
    operation is being performed in expressions like \( \vectorvar{u} + \vectorvar{v} \) and \( u + v \). Therefore,
    the following notational conventions will be used from now on for improved readability:
    \begin{enumerate*}[label=(\alph*)]
        \item \( \vectorvar{u} \oplus \vectorvar{v} \) will be denoted as \( \vectorvar{u} + \vectorvar{v} \),
        \item \( x \cdot y \) will be denoted as \( xy \),
        \item \( x \odot \vectorvar{v} \) will be denoted as \( x\vectorvar{v} \).
    \end{enumerate*}
\end{notation}

\begin{thm}
    \label{thm:vector-space-basic-properties}
    If \( \mathcal{V} \) is a vector space over a field \( \mathcal{F} \), then the following hold true for all \( x \in \abs{\mathcal{F}} \)
    and \( \vectorvar{v} \in \abs{\mathcal{V}} \):
    \begin{enumerate}[topsep=2pt,noitemsep,label=\textup{\texttt{(\alph*)}}]
        \item \( 0\vectorvar{v} = \zerovector \) \label{enum:vector-space-0-scalar-multiplication-yields-0-vector}
        \item \( x\zerovector = \zerovector \) \label{enum:vector-space-scalar-multiplication-with-0-vector-yields-zero-vector}
        \item \( \mleft( -1 \mright)\vectorvar{v} = -\vectorvar{v} \) \label{enum:vector-space-negative-1-times-v-yields-negative-v}
        \item \( \mleft( -x \mright)\vectorvar{v} = -x\vectorvar{v} \) \label{enum:vector-space-negative-x-times-v-yields-negative-xv}
        \item \( x\vectorvar{v} = \zerovector \iff \mleft( x = 0 \,\,{\lor}\,\, \vectorvar{v} = \zerovector \mright) \)
            \label{enum:vector-space-xv-equals-zero-vector-iff-x-equals-0-or-v-equals-zero-vector}
    \end{enumerate}
\end{thm}

\begin{proof}
    Let \( \vectorvar{v} \in \abs{\mathcal{V}} \) be an arbitrary vector. Since \( 0 \) is the identity of \( \abs{\mathcal{F}} \) with respect
    to \( + \), it follows that \( 0 = 0 + 0 \). Therefore, \( 0\vectorvar{v} = \mleft( 0 + 0 \mright)\vectorvar{v} = 0\vectorvar{v} + 0\vectorvar{v}\)
    by \cref{dfn:vector-space}. Since \( \tuple{\abs{\mathcal{V}}, +} \) is a \aref{dfn:group}{group} and \( 0\vectorvar{v} \in \abs{\mathcal{V}} \),
    \cref{thm:group-basic-properties}\ref{enum:group-idempotent-element-iff-element-equals-identity} states that
    \( 0\vectorvar{v} = 0\vectorvar{v} + 0\vectorvar{v} \) if and only if \( 0\vectorvar{v} \) is the identity of \( \abs{\mathcal{V}} \) with respect
    to vector addition. Since the zero vector \( \zerovector \) is the identity of \( \abs{\mathcal{V}} \) with respect to vector addition, it follows
    that \( 0\vectorvar{v} = \zerovector \), thus proving \ref{enum:vector-space-0-scalar-multiplication-yields-0-vector}.

    \vspace{0.3cm}
    \noindent
    To prove \ref{enum:vector-space-scalar-multiplication-with-0-vector-yields-zero-vector}, we utilize the fact that \( \zerovector = \zerovector + \zerovector \).
    Thus, \( x\zerovector = x \mleft( \zerovector + \zerovector \mright) = x\zerovector + x\zerovector \) by
    \cref{dfn:vector-space}\ref{enum:vector-space-scalar-multiplication-distributes-over-vector-and-field-addition}.
    Since \( \tuple{\abs{\mathcal{V}},+} \) is a group and \( x\zerovector \in \abs{\mathcal{V}} \),
    \cref{thm:group-basic-properties}\ref{enum:group-idempotent-element-iff-element-equals-identity} applies and \( x\zerovector = \zerovector \) because
    \( \zerovector \) is the identity of \( \abs{\mathcal{V}} \) with respect to vector addition.

    \vspace{0.3cm}
    \noindent
    To prove \ref{enum:vector-space-negative-1-times-v-yields-negative-v}, let \( \vectorvar{v} \in \abs{\mathcal{V}} \) be an arbitrary vector.
    By \cref{dfn:vector-space}\ref{enum:vector-space-scalar-multiplication-with-1-yields-the-original-vector} and
    \cref{dfn:vector-space}\ref{enum:vector-space-scalar-multiplication-distributes-over-vector-and-field-addition}, we have
    \( \vectorvar{v} + \mleft( -1 \mright)\vectorvar{v} = 1\vectorvar{v} + \mleft( -1 \mright)\vectorvar{v}
    = \mleft( 1 + \mleft( -1 \mright) \mright)\vectorvar{v} = 0\vectorvar{v} = \zerovector \). Similarly,
    \( \mleft( -1 \mright)\vectorvar{v} + \vectorvar{v} = \mleft( -1 \mright)\vectorvar{v} + 1\vectorvar{v}
    = \mleft( \mleft( -1 \mright) + 1 \mright)\vectorvar{v} = 0\vectorvar{v} = \zerovector \). In other words,
    \( \mleft( -1 \mright)\vectorvar{v} \) is the inverse of \( \vectorvar{v} \) with respect to
    vector addition. Since \( \tuple{\abs{\mathcal{V}}, +} \) is a group and \( \vectorvar{v} \in \abs{\mathcal{V}} \),
    it follows from \cref{thm:group-uniqueness-of-inverses} that the inverse of \( \vectorvar{v} \) is unique. Thus, we can denote
    it with the symbol \aref{notation:inverse-element-of-binary-operations}{\( -\vectorvar{v} \)}.
    In other words, \( \mleft( -1 \mright)\vectorvar{v} = -\vectorvar{v} \).

    \vspace{0.3cm}
    \noindent
    To prove \ref{enum:vector-space-negative-x-times-v-yields-negative-xv}, we note that
    \( x + \mleft( -x \mright) = \mleft( -x \mright) + x = 0 \) for all \( x \in \abs{\mathcal{F}} \)
    as implied by the \aref{dfn:ring}{ring axioms}. By \cref{dfn:vector-space}\ref{enum:vector-space-scalar-multiplication-distributes-over-vector-and-field-addition},
    we have \( \mleft( -x \mright)\vectorvar{v} + x\vectorvar{v} = \mleft( \mleft( -x \mright) + x \mright)\vectorvar{v} = 0\vectorvar{v} \).
    And by \ref{enum:vector-space-0-scalar-multiplication-yields-0-vector}, it follows that \( \mleft( -x \mright)\vectorvar{v} + x\vectorvar{v} = \zerovector \).
    Since \aref{enum:vector-space-V-and-addition-forms-abelian-group}{\( \tuple{\abs{\mathcal{V}},+} \) is an abelian group}, it follows that
    \( x\vectorvar{v} + \mleft( -x \mright)\vectorvar{v} = \mleft( -x \mright)\vectorvar{v} + x\vectorvar{v} = \zerovector \).
    By \cref{dfn:inverse-element-of-binary-operations}, it follows that \( \mleft( -x \mright)\vectorvar{v} \) is the inverse of \( x\vectorvar{v} \) with respect
    to vector addition. In other words, \( \mleft( -x \mright)\vectorvar{v} = -x\vectorvar{v} \).

    \vspace{0.3cm}
    \noindent
    To prove \ref{enum:vector-space-xv-equals-zero-vector-iff-x-equals-0-or-v-equals-zero-vector}, it suffices to show that
    \( x\vectorvar{v} = \zerovector \implies \mleft( x = 0 \,\,{\lor}\,\, \vectorvar{v} = \zerovector \mright) \) since
    the converse is an immediate consequence of \ref{enum:vector-space-0-scalar-multiplication-yields-0-vector}
    and \ref{enum:vector-space-scalar-multiplication-with-0-vector-yields-zero-vector}. First, suppose that \( x\vectorvar{v} = \zerovector \)
    for some arbitrary \( x \in \abs{\mathcal{F}} \) and \( \vectorvar{v} \in \abs{\mathcal{V}} \).
    By the \aref{rule-of-inference:law-of-excluded-middle}{Law of Excluded Middle}, we derive:
    \begin{equation}
        \label{eq:x-equals-0-or-x-neq-0}
        x = 0 \,\,{\lor}\,\, x \neq 0
    \end{equation}
    Suppose that \( x \neq 0 \). Since \( \mathcal{F} \) is a field and \( x \neq 0 \),
    the \aref{enum:field-inverse-with-respect-to-star-non-zero}{field axioms} guarantee
    that \( x \) has a unique inverse \( x^{-1} \) with respect to multiplication. Thus,
    \( x^{-1}\mleft( x\vectorvar{v} \mright) = x^{-1}\zerovector \). By \ref{enum:vector-space-0-scalar-multiplication-yields-0-vector},
    the righthand side reduces to \( \zerovector \). And by \cref{dfn:vector-space}\ref{enum:vector-space-scalar-can-be-grouped-in-scalar-multiplication},
    \( x^{-1}\mleft( x\vectorvar{v} \mright) = \mleft( x^{-1}x \mright)\vectorvar{v} = 1\vectorvar{v} = \vectorvar{v} \). Putting everything together, we have,
    \begin{gather*}
        x^{-1}\mleft( x\vectorvar{v} \mright) = \vectorvar{v} = x^{-1}\vectorvar{v} = \zerovector\\
        \therefore \, \vectorvar{v} = \zerovector
    \end{gather*}
    Since \( \vectorvar{v} = \zerovector \) is provable from the assumption that \( x \neq 0 \), we conclude by the \aref{rule-of-inference:deduction-theorem}{Deduction Theorem}:
    \begin{equation}
        \label{eq:x-neq-0-implies-v-is-zerovector}
        x \neq 0 \implies \mleft( \vectorvar{v} = \zerovector \mright)
    \end{equation}
    Note also that \( x = 0 {\implies} x = 0 \) is a tautology,
    so applying \aref{rule-of-inference:constructive-dilemma}{Constructive Dilemma} on \eqref{eq:x-equals-0-or-x-neq-0} and
    \eqref{eq:x-neq-0-implies-v-is-zerovector} yields \( x = 0 \,\,{\lor}\,\, \vectorvar{v} = \zerovector \).
\end{proof}

\begin{dfn}[Linear Combination]
    \phantomsection \label{dfn:vector-space-linear-combination}
    Let \( \mathcal{V} \) be a vector space over a field \( \mathcal{F} \) and \( k \in \naturalset \).
    A vector \( \vectorvar{v} \in \abs{\mathcal{V}} \) is a \keyword{linear combination} of
    the vectors \( \vectorvar{a}_0, \dots, \vectorvar{a}_{k-1} \in \abs{\mathcal{V}} \) if and only
    if \( \exists c_{0}\,{\in}\,\abs{\mathcal{F}} \, \dots, \exists c_{k-1}\,{\in}\,\abs{\mathcal{F}} \,
    \big[ \vectorvar{v} = c_{0}\vectorvar{a}_0 + \dots + c_{k-1}\vectorvar{a}_{k-1} \big] \).
    The scalars \( c_{0}, \dots, c_{k-1} \) are called the \keyword{coefficients} of the linear combination.
\end{dfn}

\begin{dfn}[Linear Dependence]
    \phantomsection \label{dfn:vector-space-linear-dependence}
    Let \( \mathcal{V} \) be a vector space over a field \( \mathcal{F} \).
    If \( S \) is a non-empty subset of \( \abs{\mathcal{V}} \), then the elements
    of \( S \) are \keyword{linearly dependent} if and only if there exists a finite, non-empty subset
    \( \mleft\{ \vectorvar{s}_{0},\dots,\vectorvar{s}_{k-1} \mright\} \) of \( S \) such that
    \( \sum_{i \in k} a_{i}\vectorvar{s}_{i} = \zerovector \) for some
    \( a_{0},\dots,a_{k-1} \in \abs{\mathcal{F}} \) where \( \exists i\,{\in}\,k \, \big[ a_{i} \neq 0 \big] \).
\end{dfn}

\begin{dfn}[Linear Independence]
    \phantomsection \label{dfn:vector-space-linear-independence}
    Let \( \mathcal{V} \) be a vector space over a field \( \mathcal{F} \).
    If \( S \) is a non-empty subset of \( \abs{\mathcal{V}} \), then the elements
    of \( S \) are \keyword{linearly independent} if and only if they are not linearly dependent.
\end{dfn}

\begin{thm}
    \phantomsection \label{thm:vector-space-linearly-independent-iff-every-finite-subset-can-only-form-0-vector-with-all-0-coefficients}
    If \( \mathcal{V} \) is a vector space over a field \( \mathcal{F} \) and \( S \) is a subset of \( \abs{\mathcal{V}} \), then the
    elements of \( S \) are linearly independent if and only if for every finite subset
    \( \mleft\{ \vectorvar{s}_{0}, \dots, \vectorvar{s}_{k-1} \mright\} \) of \( S \),
    \begin{equation*}
        \forall a_{0}\,{\in}\,\abs{\mathcal{F}} \, \dots \forall a_{k-1}\,{\in}\,\abs{\mathcal{F}} \,
        \bigg[ \sum_{i \in k} a_{i}\vectorvar{s}_{i} = \zerovector \iff \forall i\,{\in}\,k \, \big( a_{i} = 0 \big) \bigg]
    \end{equation*}
\end{thm}

\begin{proof}
    Observe that \cref{dfn:vector-space-linear-independence} is a negation of \cref{dfn:vector-space-linear-dependence}. Thus,
    we conclude by \aref{rule-of-inference:existential-definition}{Existential Definition} that
    the elements of \( S \) are linearly independent if and only if for every finite,
    non-empty subset \( \mleft\{ \vectorvar{s}_{0}, \dots, \vectorvar{s}_{k-1} \mright\} \) of \( S \),
    \begin{equation}
        \phantomsection \label{eq:negate-the-existential-statement}
        \lnot \, \exists a_{0}\,{\in}\,\abs{\mathcal{F}} \, \dots \exists a_{k-1}\,{\in}\,\abs{\mathcal{F}} \,
        \bigg[ \sum_{i \in k} a_{i}\vectorvar{s}_{i} = \zerovector \,\,{\land}\,\, \exists i\,{\in}\,k \, \mleft( a_{i} \neq 0 \mright) \bigg]
    \end{equation}
    Finally, we invoke \aref{rule-of-inference:existential-definition}{Existential Definition} on \eqref{eq:negate-the-existential-statement}
    to derive the needed formula:
    \begin{equation*}
        \forall a_{0}\,{\in}\,\abs{\mathcal{F}} \, \dots \forall a_{k-1}\,{\in}\,\abs{\mathcal{F}} \,
        \bigg[ \sum_{i \,\in\, k} a_{i}\vectorvar{s}_{i} = \zerovector \implies \forall i\,{\in}\,k \, \mleft( a_{i} = 0 \mright) \bigg] \qedhere
    \end{equation*}
\end{proof}

\begin{notation}
    \phantomsection \label{notation:linear-dependence-and-independence}
    We write \( \vectorvar{u} \parallel \vectorvar{v} \) to mean that the vectors \( \vectorvar{u} \) and \( \vectorvar{v} \) are linearly
    dependent on each other. Conversely, \( \vectorvar{u} \nparallel \vectorvar{v} \) means that \( \vectorvar{u} \) and \( \vectorvar{v} \)
    are linearly independent. The notation can also be extended to sets (finite or infinite) of vectors.
    \( {\parallel} S \) states that `\( S \) is a set of linearly dependent vectors' whereas \( {\nparallel}S \) states that `\( S \) is a set of linearly
    independent vectors.'
\end{notation}

\begin{dfn}[Subspace]
    \phantomsection \label{dfn:vector-subspace}
    If \( \mathcal{V} = \tuple{V, \mathcal{F}, +, \cdot} \) is a \aref{dfn:vector-space}{vector space} over \( \mathcal{F} \),
    then a structure \( \mathcal{W} \) is a vector \keyword{subspace} of \( \mathcal{V} \) if and only if:
    \begin{enumerate}[topsep=2pt,noitemsep,label=\textup{\texttt{(\alph*)}}]
        \item \( \abs{\mathcal{W}} \subseteq V \)
        \item \( \mathcal{W} \) is equipped with the binary operation \( \smallsubscript{+}{\abs{\mathcal{W}}} = \restrictedfunction{+}{\abs{\mathcal{W}}^{2}} \)
        \item \( \mathcal{W} \) is equipped with the binary function \( \smallsubscript{\cdot}{\abs{\mathcal{W}}} = \restrictedfunction{\cdot}{\prod \tuple{\abs{\mathcal{F}},\abs{\mathcal{W}}}}\)
        \item \( \mathcal{W} \) is a vector space over \( \mathcal{F} \)
    \end{enumerate}
\end{dfn}

\begin{notation}
    In general, given any \( S \subseteq \abs{\mathcal{V}} \), the restriction on the vector addition of \( \mathcal{V} \) to \( S^{2} \) is denoted as \( \smallsubscript{+}{S} \).
    Similarly, the restriction on the scalar multiplication of \( \mathcal{V} \) to \( \prod \tuple{\abs{\mathcal{F}}, S} \) is denoted as \( \smallsubscript{\cdot}{S} \).
\end{notation}

\begin{notation}
    The definition for vector subspaces is clearly consistent with the definition of \aref{dfn:substructure}{substructures}
    (i.e., a vector space \( \mathcal{W} \) is a vector subspace of \( \mathcal{V} \) if and only if \( \mathcal{W} \) is a \aref{dfn:substructure}{substructure}
    of \( \mathcal{V} \)). Thus, we can write \( \mathcal{W} \leq \mathcal{V} \) to mean `\( \mathcal{W} \) is a vector subspace of \( \mathcal{V} \)' as long as
    it is understood from the context that \( \mathcal{W} \) and \( \mathcal{V} \) denote vector spaces over the same field.
\end{notation}

\begin{dfn}
    \phantomsection \label{dfn:set-of-all-subsets-that-form-vector-subspaces}
    Given any vector space \( \mathcal{V} \) over \( \mathcal{F} \), \( \vectorsubspaces{\mathcal{V}} \) denotes the set of all non-empty subsets of \( \abs{\mathcal{V}} \)
    with which vector subspaces of \( \mathcal{V} \) can be formed by restricting the \( + \) and \( \cdot \) operations of \( \mathcal{V} \).
    In other words,
    \begin{equation*}
        \vectorsubspaces{\mathcal{V}}
        =
        \setbuilderBig{X \subseteq \abs{\mathcal{V}}}{X \neq \emptyset \,\,{\land}\,\,
        \tuple{X, \mathcal{F}, \smallsubscript{+}{X}, \smallsubscript{\cdot}{X}} \leq \mathcal{V} }
    \end{equation*}
\end{dfn}

\begin{thm}
    \phantomsection \label{thm:all-vector-subspaces-of-V-share-the-same-zero-vector-as-V}
    If \( \mathcal{V} \) is a vector space and \( \zerovector \) denotes the \aref{enum:vector-space-zero-vector}{zero vector} of \( \mathcal{V} \),
    then for any vector subspace \( \mathcal{W} \) of \( \mathcal{V} \):
    \begin{enumerate}[topsep=2pt,noitemsep,label=\textup{\texttt{(\alph*)}}]
        \item \( \zerovector \in \abs{\mathcal{W}} \)
        \item \( \zerovector \) is the identity of \( \abs{\mathcal{W}} \) with respect to \( \smallsubscript{+}{\abs{\mathcal{W}}} \)
    \end{enumerate}
\end{thm}

\begin{proof}
    Let \( \mathcal{V} = \tuple{V, \mathcal{F}, +, \cdot} \) be a vector space over \( \mathcal{F} \).
    Let \( \mathcal{W} = \tuple{W, \mathcal{F}, \smallsubscript{+}{W}, \smallsubscript{\cdot}{W} } \) be any
    vector subspace of \( \mathcal{V} \). By \cref{dfn:vector-subspace}, we know that \( W \subseteq V \) and that
    \( \smallsubscript{+}{W} \) is the function \( + \) restricted to \( W^{2} \). In other words,
    \begin{equation}
        \label{eq:plusW-is-plus-restricted-to-W2}
        \forall \vectorvar{x}\,{\in}\,W \, \forall \vectorvar{y}\,{\in}\,W \, \big[ \vectorvar{x} \, \smallsubscript{+}{W} \, \vectorvar{y} = \vectorvar{x} + \vectorvar{y} \big]
    \end{equation}
    Since \( \mathcal{W} \) is also a vector space over \( \mathcal{F} \),
    it follows from \aref{dfn:vector-space}{the vector space axioms} that \( \tuple{W, \smallsubscript{+}{W}} \)
    is an abelian group, which means that \( W \) must have a unique identity with respect to \( \smallsubscript{+}{W} \).
    Let \( \vectorvar{e} \) be the identity of \( W \) with respect to \( \smallsubscript{+}{W} \).
    Clearly,
    \begin{equation}
        \label{eq:e-is-the-identity-of-plusW}
        \vectorvar{e} \, \smallsubscript{+}{W} \, \vectorvar{e} = \vectorvar{e}
    \end{equation}
    Furthermore, since \( \vectorvar{e} \in W \), we can \aref{rule-of-inference:universal-instantiation}{instantiate} \eqref{eq:plusW-is-plus-restricted-to-W2}
    to obtain,
    \begin{equation}
        \label{eq:e-plusW-e-equals-e}
        \vectorvar{e} \, \smallsubscript{+}{W} \, \vectorvar{e} = \vectorvar{e} + \vectorvar{e}
    \end{equation}
    Substituting \eqref{eq:e-is-the-identity-of-plusW} into \eqref{eq:e-plusW-e-equals-e} then yields \( \vectorvar{e} + \vectorvar{e} = \vectorvar{e} \).
    But \( \tuple{V, +} \) is a group, so \cref{thm:group-basic-properties}\ref{enum:group-idempotent-element-iff-element-equals-identity} applies,
    which means that \( \vectorvar{e} \) must be the identity of \( V \) with respect to \( + \).
    In other words, \( \vectorvar{e} = \zerovector \).
\end{proof}

\begin{thm}
    \phantomsection \label{thm:the-insersection-of-the-set-of-all-vector-subspaces-is-the-singleton-with-only-empty-set-as-its-element}
    For any vector space \( \mathcal{V} \), \( \Intersect \vectorsubspaces{\mathcal{V}} = \mleft\{ \zerovector \mright\} \).
\end{thm}

\begin{proof}
    Let \( \mathcal{V} \) be a vector space over \( \mathcal{F} \) and \( S = \Intersect \vectorsubspaces{\mathcal{V}} \).
    Thus, by \cref{dfn:intersection-of-sets}, the following first-order sentence must hold true:
    \begin{equation}
        \label{eq:S-is-the-intersection-of-subV}
        \forall \vectorvar{x} \, \big[ \vectorvar{x} \in S \iff \forall Z\,{\in}\,\vectorsubspaces{\mathcal{V}} \, \mleft( \vectorvar{x} \in Z \mright) \big]
    \end{equation}
    Let \( W \) be an arbitrary element of \( \vectorsubspaces{\mathcal{V}} \). Thus, \( W \subseteq \abs{\mathcal{V}} \) and
    \( \tuple{W, \mathcal{F}, \smallsubscript{+}{W}, \smallsubscript{\cdot}{W}} \leq \mathcal{V} \) by \cref{dfn:set-of-all-subsets-that-form-vector-subspaces}.
    Since \( \zerovector \) is the identity of \( \abs{\mathcal{V}} \) with respect to \( + \)
    and \( \tuple{W, \mathcal{F}, \smallsubscript{+}{W}, \smallsubscript{\cdot}{W}} \) is a \aref{dfn:vector-subspace}{vector subspace} of \( \mathcal{V} \),
    it follows from \cref{thm:all-vector-subspaces-of-V-share-the-same-zero-vector-as-V} that \( \zerovector \in W \). Since \( W \in \vectorsubspaces{\mathcal{V}} \)
    is arbitrary, we have by \aref{rule-of-inference:universal-generalization}{Universal Generalization},
    \( \forall Z\,{\in}\,\vectorsubspaces{\mathcal{V}} \, \mleft[ \zerovector \in Z \mright] \).
    Thus, \( \zerovector \in S \) by \eqref{eq:S-is-the-intersection-of-subV}.
    In other words,
    \begin{equation}
        \label{eq:if-x-is-zerovector-then-x-must-be-an-element-of-S}
        \forall \vectorvar{x} \, \big[ \vectorvar{x} = \zerovector \implies \vectorvar{x} \in S \big]
    \end{equation}
    To prove the converse, it suffices to show that \( \mleft\{ \zerovector \mright\} \in \vectorsubspaces{\mathcal{V}} \).
    For convenience, we let \( T = \mleft\{ \zerovector \mright\} \). Suppose also that \( \function{\smallsubscript{+}{T}}{T^{2}}{T} \)
    and \( \function{\smallsubscript{(\cdot)}{T}}{\prod \tuple{\abs{\mathcal{F}}, T}}{T} \) are functions with the following definitions:
    \begin{equation*}
        \begin{aligned}
            \smallsubscript{+}{T} &= \big\{ \orderedpair{\tuple{\zerovector,\zerovector}}{\zerovector} \big\} \\
            \smallsubscript{(\cdot)}{T} &=
            \setbuilderbig{\orderedpair{x}{y}}{ \exists a\,{\in}\,\abs{\mathcal{F}} \, \mleft[ x = \tuple{a, \zerovector} \mright] \,\,{\land}\,\, y = \zerovector }
        \end{aligned}
    \end{equation*}
    It is not difficult to see that the operations \( \smallsubscript{+}{T} \) and \( + \) yields the same value for all elements of \( T \).
    After all, \( \zerovector \) is the sole element of \( T \) and \(  \zerovector + \zerovector = \zerovector = \zerovector \, \smallsubscript{+}{T} \, \zerovector \).
    The same is true for \( \smallsubscript{(\cdot)}{T} \) and \( (\cdot) \) since
    \( c \cdot \zerovector = \zerovector = c \, \smallsubscript{\cdot}{T} \, \zerovector \) for all \( c \in \abs{\mathcal{F}} \). In other words,
    \( \smallsubscript{+}{T} = \restrictedfunction{+}{T^2} \) and \( \smallsubscript{(\cdot)}{T} = \restrictedfunction{(\cdot)}{\prod \tuple{\abs{\mathcal{F}},T}} \).
    And since \( T \subseteq \abs{\mathcal{V}} \), it follows that the operations \( \smallsubscript{+}{T} \) and \( \smallsubscript{(\cdot)}{T} \) share
    the exact same properties with \( + \) and \( (\cdot) \) respectively.
    Simply put, the structure \( \tuple{T, \mathcal{F}, \smallsubscript{+}{T}, \smallsubscript{\cdot}{T} }\)
    is a vector subspace of \( \mathcal{V} \) and thus \( T = \mleft\{ \zerovector \mright\} \in \vectorsubspaces{\mathcal{V}} \)
    by \cref{dfn:set-of-all-subsets-that-form-vector-subspaces}.

    \vspace{0.3cm}
    \noindent
    Now let \( \vectorvar{x}_{0} \in S \) be arbitrary. By \eqref{eq:S-is-the-intersection-of-subV},
    it follows that \( \forall Z\,{\in}\,\vectorsubspaces{\mathcal{V}} \, \mleft( \vectorvar{x}_{0} \in Z \mright) \).
    But \( \mleft\{ \zerovector \mright\} \in \vectorsubspaces{\mathcal{V}} \), so
    \( \vectorvar{x}_{0} \in \mleft\{ \zerovector \mright\} \) must hold by \aref{rule-of-inference:universal-instantiation}{Universal Instantiation}.
    This in turn implies that \( \vectorvar{x}_{0} = \zerovector \) since \( \zerovector \) is the sole element of the singleton \( \mleft\{ \zerovector \mright\} \).
    By \aref{rule-of-inference:universal-generalization}{Universal Generalization}, we conclude that,
    \begin{equation}
        \label{eq:if-x-is-an-element-of-x-then-x-must-be-the-zerovector}
        \forall \vectorvar{x} \, \big[ \vectorvar{x} \in S \implies \vectorvar{x} = \zerovector \big]
    \end{equation}
    Applying \aref{rule-of-inference:biconditional-introduction}{Biconditional Introduction} to \eqref{eq:if-x-is-zerovector-then-x-must-be-an-element-of-S}
    and \eqref{eq:if-x-is-an-element-of-x-then-x-must-be-the-zerovector} then yields,
    \begin{equation*}
        \forall \vectorvar{x} \, \big[ \vectorvar{x} \in S \iff \vectorvar{x} = \zerovector \big]
    \end{equation*}
    In other words, \( S = \Intersect \vectorsubspaces{\mathcal{V}} = \mleft\{ \zerovector \mright\} \).
\end{proof}

\begin{thm}
    \phantomsection \label{thm:intersection-of-a-collection-of-vector-subspaces-is-also-a-vector-subspace}
    For any vector space \( \mathcal{V} \), the intersection of every non-empty subset of \( \vectorsubspaces{\mathcal{V}} \)
    also belongs to \( \vectorsubspaces{\mathcal{V}} \).
    \begin{equation*}
        \forall X \, \Big[ X \subseteq \vectorsubspaces{\mathcal{V}} \,\,{\land}\,\, X \neq \emptyset \implies \Intersect X \in \vectorsubspaces{\mathcal{V}} \Big]
    \end{equation*}
\end{thm}

\begin{proof}
    Let \( \mathcal{V} \) be a vector space over a field \( \mathcal{F} \).
    Let \( X \) be an arbitrary, non-empty subset of \( \vectorsubspaces{\mathcal{V}} \) and
    \( S \) be the intersection of \( X \). Thus, it follows from \cref{dfn:intersection-of-sets} that,
    \begin{equation}
        \phantomsection \label{eq:S-is-the-intersection-of-X}
        \forall x \, \big[ x \in S \iff \forall z\,{\in}\,X \, \big( x \in z \big) \big]
    \end{equation}

    \vspace{0.3cm}
    \noindent
    To prove that \( S \in \vectorsubspaces{\mathcal{V}} \), one must show that \( S \) satisfies all the conditions stated in
    \cref{dfn:set-of-all-subsets-that-form-vector-subspaces}; namely:
    \begin{enumerate*}[label=(\roman*)]
        \item \( S \neq \emptyset \), \phantomsection \label{enum:S-is-non-empty}
        \item \( \tuple{S, \mathcal{F}, \smallsubscript{+}{S}, \smallsubscript{\cdot}{S} } \) is a vector subspace of \( \mathcal{V} \).
            \phantomsection \label{enum:S-is-a-vector-subspace-of-V}
    \end{enumerate*}

    \vspace{0.3cm}
    \noindent
    {\scshape Claim} \ref{enum:S-is-non-empty}: \( S \neq \emptyset \)

    \vspace{0.1cm}
    \noindent
    Since \( \mathcal{V} \) is a vector space,
    we know from the \aref{dfn:vector-space}{axioms defining vector spaces} that \( \tuple{\abs{\mathcal{V}}, +} \)
    forms an abelian group. This means that \( \abs{\mathcal{V}} \) has a \aref{thm:-group-uniqueness-of-identity}{unique identity} with respect to \( + \).
    This unique element is in fact the \aref{enum:vector-space-zero-vector}{zero vector} \( \zerovector \) of \( \mathcal{V} \).

    \vspace{0.3cm}
    \noindent
    Now let \( W \) be an arbitrary element of \( X \). Since \( X \) is a subset of \( \vectorsubspaces{\mathcal{V}} \),
    necessarily \( W \in \vectorsubspaces{\mathcal{V}} \). By \cref{dfn:set-of-all-subsets-that-form-vector-subspaces},
    it follows that \( \tuple{W, \mathcal{F}, \smallsubscript{+}{W}, \smallsubscript{\cdot}{W} } \)
    is a \aref{dfn:vector-subspace}{vector subspace} of \( \mathcal{V} \);
    in other words, \( \tuple{W, \mathcal{F}, \smallsubscript{+}{W}, \smallsubscript{\cdot}{W} } \)
    is also a vector space over \( \mathcal{F} \) and \( W \subseteq \abs{\mathcal{V}} \). Moreover, \( \tuple{W, \smallsubscript{+}{W} } \)
    forms an abelian group, which means that there exists an \( \vectorvar{e} \in W \) such that \( \vectorvar{e} \) is the
    \aref{thm:-group-uniqueness-of-identity}{unique identity} of \( W \) with respect to \( \smallsubscript{+}{W} \). Thus,
    \begin{equation}
        \label{eq:e-is-the-identity-of-W-with-respect-to-vector-addition}
        \forall \vectorvar{x}\,{\in}\,W \, \big[ \vectorvar{x} \,\, \smallsubscript{+}{W} \,
        \vectorvar{e} = \vectorvar{e} \,\,\smallsubscript{+}{W} \, \vectorvar{x} = \vectorvar{x}  \big]
    \end{equation}
    Since \( \vectorvar{e} \in W \), we can \aref{rule-of-inference:universal-instantiation}{instantiate}
    \eqref{eq:e-is-the-identity-of-W-with-respect-to-vector-addition} to obtain \( \vectorvar{e} \,\, \smallsubscript{+}{W}\, \vectorvar{e} = \vectorvar{e} \).
    Note that \( \smallsubscript{+}{W} \) is the vector addition of \( \mathcal{V} \) with its domain \aref{dfn:restricted-functions}{restricted} to \( W^{2} \).
    This means that for all \( \vectorvar{x}, \vectorvar{y} \in W \), \( \vectorvar{x} + \vectorvar{y} \) is also defined and that
    \( \vectorvar{x} + \vectorvar{y} = \vectorvar{x} \,\, \smallsubscript{+}{W} \, \vectorvar{y} \).
    Since \( \vectorvar{e} \,\, \smallsubscript{+}{W} \, \vectorvar{e} = \vectorvar{e} \), we conclude that \( \vectorvar{e} + \vectorvar{e} = \vectorvar{e} \).
    Clearly, \( \vectorvar{e} \in \abs{\mathcal{V}} \) because \( \vectorvar{e} \in W \) and \( W \subseteq \abs{\mathcal{V}} \).
    Since \( \tuple{\abs{\mathcal{V}}, +} \) is a group and \( \vectorvar{e} \in \abs{\mathcal{V}} \), \aref{rule-of-inference:universal-instantiation}{instantiating}
    \cref{thm:group-basic-properties}\ref{enum:group-idempotent-element-iff-element-equals-identity} yields
    \( \vectorvar{e} + \vectorvar{e} = \vectorvar{e} {\iff} \vectorvar{e} = \zerovector \). But \( \vectorvar{e} + \vectorvar{e} = \vectorvar{e} \) has already been established
    above, so \( \vectorvar{e} = \zerovector \) follows by \aref{rule-of-inference:modus-ponens}{Modus Ponens}.
    Since \( \vectorvar{e} \) belongs to \( W \), \( \zerovector \) must also belongs to \( W \) because \( \vectorvar{e} \) and \( \zerovector\) are equal.
    In other words, we have just shown that the formula \( \zerovector \in W \) is provable from the initial assumption that \( W \in X \). Thus, we have by the
    \aref{rule-of-inference:deduction-theorem}{Deduction Theorem},
    \begin{equation*}
        W \in X \implies \zerovector \in W
    \end{equation*}
    Since \( W \in X \) is arbitrary, \aref{rule-of-inference:universal-generalization}{Universal Generalization} can be applied and we derive
    \( \forall z\,{\in}\,X \, \big( \zerovector \in z \big) \) as a result, which can then be applied to \eqref{eq:S-is-the-intersection-of-X}
    to yield \( \zerovector \in S \).
    Since the set \( S \) has the zero vector of \( \mathcal{V} \) as an element, it cannot possibly be empty.

    \vspace{0.5cm}
    \noindent
    {\scshape Claim} \ref{enum:S-is-a-vector-subspace-of-V}: \( \tuple{S, \mathcal{F}, \smallsubscript{+}{S}, \smallsubscript{\cdot}{S}} \)
        is a \aref{dfn:vector-subspace}{vector subspace} of \( \mathcal{V} \)

    \vspace{0.1cm}
    \noindent
    We begin by showing that \( S \subseteq \abs{\mathcal{V}} \). Suppose that \( \vectorvar{x}_{0} \) is an arbitrary element of \( S \).
    We know that \( X \) is a non-empty subset of \( \vectorsubspaces{\mathcal{V}} \), which means that \( X \) must contain at least one element.
    Let \( Z_{0} \) denote such an element; that is, \( Z_{0} \in X \).
    Since \( \vectorvar{x}_{0} \in S \), it follows from \eqref{eq:S-is-the-intersection-of-X} that \( \forall z\,{\in}\,X \, \mleft[ \vectorvar{x}_{0} \in z \mright] \).
    Since \( Z_{0} \in X \), we conclude by \aref{rule-of-inference:universal-instantiation}{Universal Instantiation} that \( \vectorvar{x}_{0} \in Z_{0} \).
    But \( X \subseteq \vectorsubspaces{\mathcal{V}} \), so \( Z_{0} \in X \) implies \( Z_{0} \in \vectorsubspaces{\mathcal{V}} \).
    Thus, it follows from \cref{dfn:set-of-all-subsets-that-form-vector-subspaces} that \( Z_{0} \subseteq \abs{\mathcal{V}} \),
    from which we can easilly infer \( \vectorvar{x}_{0} \in \abs{\mathcal{V}} \) given that \( \vectorvar{x}_{0} \in Z_{0} \).
    And since \( \vectorvar{x}_{0} \in \abs{\mathcal{V}} \) is provable from the assumption that \( \vectorvar{x}_{0} \in S \), we have by the
    \aref{rule-of-inference:deduction-theorem}{Deduction Theorem} \( \vectorvar{x}_{0} \in S {\implies} \vectorvar{x}_{0} \in \abs{\mathcal{V}} \).
    With \( \vectorvar{x}_{0} \) being arbitrary,
    \( S \subseteq \abs{\mathcal{V}} \) follows readily from \aref{rule-of-inference:universal-generalization}{Universal Generalization}
    and the \aref{dfn:subset}{definition of subsets}.

    \vspace{0.3cm}
    \noindent
    Next, we let \( \smallsubscript{+}{S} \) and \( \smallsubscript{(\cdot)}{S} \) be the functions obtained by restricting \( + \) and \( (\cdot) \)
    to \( S^{2} \) and \( \prod \tuple{\abs{\mathcal{F}}, S} \) respectively.
    In other words,
    \begin{equation}
        \label{eq:vector-addition-is-restricted-to-S2-scalar-multiplication-is-restricted-to-F-times-S}
        \begin{gathered}
            \forall \vectorvar{x}\,{\in}\,S \, \forall \vectorvar{y}\,{\in}\,S \,
            \big[ \vectorvar{x} \,\, \smallsubscript{+}{S} \, \vectorvar{y} = \vectorvar{x} + \vectorvar{y} \big]\\
            \forall \vectorvar{x}\,{\in}\,S \, \forall a\,{\in}\,\abs{\mathcal{F}} \,
            \big[ a \, \smallsubscript{\cdot}{S} \, \vectorvar{x} = a \cdot \vectorvar{x} \big]
        \end{gathered}
    \end{equation}
    Naturally, the first order of business is to verify that \( \tuple{S, \smallsubscript{+}{S}} \) forms an \aref{dfn:abelian-group}{abelian group} since
    this is a requirement for \( \tuple{S, \mathcal{F}, \smallsubscript{+}{S}, \smallsubscript{\cdot}{S}} \) to be a vector space.
    For \( \tuple{S, \smallsubscript{+}{S}} \) to be a group, \( S \) has to be \aref{dfn:closed-operation}{closed under}
    the binary operation \( \smallsubscript{+}{S} \) first and foremost.
    To verify this, we let \( \vectorvar{x}_{0} \) and \( \vectorvar{y}_{0} \) represent arbitrary elements of \( S \)
    and then use them to \aref{rule-of-inference:universal-instantiation}{instantiate} the first formula in
    \eqref{eq:vector-addition-is-restricted-to-S2-scalar-multiplication-is-restricted-to-F-times-S}:
    \begin{equation}
        \label{eq:instantiated-vector-addition-restricted-to-S2}
        \vectorvar{x}_{0} \, \smallsubscript{+}{S} \, \vectorvar{y}_{0} = \vectorvar{x}_{0} + \vectorvar{y}_{0}
    \end{equation}
    Recall that \( X \) is a non-empty subset of \( \vectorsubspaces{\mathcal{V}} \).
    This means that \( X \) must contain at least one element, which we now denote as \( W \). 
    Since \( \vectorvar{x}_{0}, \vectorvar{y}_{0} \in S \) and \( W \in X \), it follows from \eqref{eq:S-is-the-intersection-of-X} that
    \( \vectorvar{x}_{0}, \vectorvar{y}_{0} \in W \). Furthermore, since \( X \subseteq \vectorsubspaces{V} \) and \( W \in X \), it follows that
    \( W \in \vectorsubspaces{\mathcal{V}} \). Thus, by \cref{dfn:set-of-all-subsets-that-form-vector-subspaces},
    the structure \( \tuple{W, \mathcal{F}, \smallsubscript{+}{W}, \smallsubscript{\cdot}{W}} \) is a
    \aref{dfn:vector-subspace}{vector subspace} of \( \mathcal{V} \) and, by extension, a vector space over \( \mathcal{F} \).
    This means that \( W \) is \aref{dfn:closed-operation}{closed under} \( \smallsubscript{+}{W} \) as required by
    the \aref{dfn:vector-space}{definition of vector spaces}.
    Thus,
    \begin{equation}
        \label{eq:W-is-closed-under-plus-W}
        (\vectorvar{x}_{0} \, \smallsubscript{+}{W} \, \vectorvar{y}_{0}) \in W
    \end{equation}
    Since the operation \( \smallsubscript{+}{W} \) is the operation \( + \) restricted to \( W^{2} \), it follows that,
    \begin{equation}
        \label{eq:vector-addition-is-restricted-to-W2}
        \vectorvar{x}_{0} \, \smallsubscript{+}{W} \, \vectorvar{y}_{0} = \vectorvar{x}_{0} + \vectorvar{y}_{0}
    \end{equation}
    It should be obvious from \eqref{eq:instantiated-vector-addition-restricted-to-S2} and \eqref{eq:vector-addition-is-restricted-to-W2} that
    \( \vectorvar{x}_{0} \, \smallsubscript{+}{W} \, \vectorvar{y}_{0} \) and \(  \vectorvar{x}_{0} \, \smallsubscript{+}{S} \, \vectorvar{y}_{0} \)
    are equal. This equality permits us to conclude from \eqref{eq:W-is-closed-under-plus-W} that \( \vectorvar{x}_{0} \, \smallsubscript{+}{S} \, \vectorvar{y}_{0} \in W \).
    Since \( W \) is an arbitrary element of \( X \), we have by \aref{rule-of-inference:universal-generalization}{Universal Generalization}
    \( \forall z\,{\in}\,X \, \mleft[ \mleft( \vectorvar{x}_{0} \, \smallsubscript{+}{S} \, \vectorvar{y}_{0} \mright) \in z \mright] \).
    And by \eqref{eq:S-is-the-intersection-of-X}, necessarily \( \mleft( \vectorvar{x}_{0} \, \smallsubscript{+}{S} \, \vectorvar{y}_{0} \mright) \in S \).
    Since \( \vectorvar{x}_{0} \) and \( \vectorvar{y}_{0} \) are both arbitrary elements of \( S \),
    we conclude by \aref{rule-of-inference:universal-generalization}{Universal Generalization} once again that
    \(  \forall \vectorvar{x}\,{\in}\,S \, \forall \vectorvar{y}\,{\in}\,S \,
        \mleft[ \mleft( \vectorvar{x} \, \smallsubscript{+}{S} \, \vectorvar{y} \mright) \in S \mright] \).
    In other words, \( S \) is closed under \( \smallsubscript{+}{S} \).

    \vspace{0.3cm}
    \noindent
    Next, we verify the associativity and commutativity of \( \smallsubscript{+}{S} \), which should be fairly obvious from
    \eqref{eq:vector-addition-is-restricted-to-S2-scalar-multiplication-is-restricted-to-F-times-S} and the fact that \( + \) itself is both associative and commutative.
    Thus, for any \( \vectorvar{x}, \vectorvar{y}, \vectorvar{z} \in S \):
    \begin{gather*}
        \vectorvar{x} \, \smallsubscript{+}{S} \, \mleft( \vectorvar{y} \, \smallsubscript{+}{S} \, \vectorvar{z} \mright)
        = \vectorvar{x} + \mleft( \vectorvar{y} + \vectorvar{z} \mright)
        = \mleft( \vectorvar{x} + \vectorvar{y} \mright) + \vectorvar{z}
        = \mleft( \vectorvar{x}_{0} \, \smallsubscript{+}{S} \, \vectorvar{y}_{0} \mright) \, \smallsubscript{+}{S} \, \vectorvar{z}\\
        \vectorvar{x} \, \smallsubscript{+}{S} \, \vectorvar{y}
        = \vectorvar{x} + \vectorvar{y}
        = \vectorvar{y} + \vectorvar{x}
        = \vectorvar{y} \, \smallsubscript{+}{S} \, \vectorvar{x}
    \end{gather*}
    Moreoever, since \( \zerovector \in S \), it follows that \( \zerovector \, \smallsubscript{+}{S} \, \vectorvar{x}
    = \zerovector + \vectorvar{x}
    = \vectorvar{x} + \zerovector
    = \vectorvar{x} \, \smallsubscript{+}{S} \, \zerovector
    = \zerovector + \vectorvar{x}
    = \vectorvar{x} \)
    for all \( \vectorvar{x} \in S \).
    In other words, \( \zerovector \) is an identity of \( S \) with respect to \( \smallsubscript{+}{S} \).
    To summarize, it has been established that \( \smallsubscript{+}{S} \) is a total, associative, and commutative binary operation on \( S \)
    with an identity element. Therefore, \( \tuple{S, \smallsubscript{+}{S}} \) is necessarily an abelian group.

    \vspace{0.3cm}
    \noindent
    We now turn our attention to the vector space properties involving scalar multiplication.
    We already know from \eqref{eq:vector-addition-is-restricted-to-S2-scalar-multiplication-is-restricted-to-F-times-S} that 
    \( \dom (\smallsubscript{\cdot}{S}) = \prod \tuple{\abs{\mathcal{F}}, S} \); however, it is still necessary to prove that
    \( \ran (\smallsubscript{\cdot}{S}) \subseteq S \) since this is one of the conditions stated in \cref{dfn:vector-space}.
    And we do this by showing that the following formula can be derived from all active assumptions:
    \begin{equation}
        \label{eq:S-is-closed-under-scalar-multiplication}
        \forall \vectorvar{x}\,{\in}\,S \, \forall a\,{\in}\,\abs{\mathcal{F}} \,
        \big[ a \, \smallsubscript{\cdot}{S} \, \vectorvar{x} \in S \big]
    \end{equation}
    Suppose that \( \vectorvar{x}_{0} \) is an arbitrary element of \( S \). Let \( W \) be an arbitrary element of \( X \) so that
    the structure \( \tuple{W, \mathcal{F}, \smallsubscript{+}{W}, \smallsubscript{\cdot}{W} } \) is a vector subspace of \( \mathcal{V} \)
    where \( \smallsubscript{(\cdot)}{W} \) is the restriction of \( (\cdot) \) to \( \prod \tuple{\abs{\mathcal{F}}, W} \).
    Since \( \vectorvar{x}_{0} \in S \) and \( W \in X \), it follows from \eqref{eq:S-is-the-intersection-of-X} that \( \vectorvar{x}_{0} \in W \).
    Let \( a_{0} \in \abs{\mathcal{F}} \) be arbitrary. Clearly, \( a_{0} \, \smallsubscript{\cdot}{W} \, \vectorvar{x}_{0} \) is defined and,
    \begin{equation}
        \label{eq:a0-cdotW-x0-equals-a0-cdot-x0-equals-a0-cdotS-x0}
        a_{0} \, \smallsubscript{\cdot}{W} \, \vectorvar{x}_{0} = a_{0} \cdot \vectorvar{x}_{0}
        = a_{0} \, \smallsubscript{\cdot}{S} \, \vectorvar{x}_{0}
    \end{equation}
    But because \( \tuple{W, \mathcal{F}, \smallsubscript{+}{W}, \smallsubscript{\cdot}{W}} \) is a \aref{dfn:vector-space}{vector space} over \( \mathcal{F} \),
    necessarily \( \ran (\smallsubscript{\cdot}{W}) \subseteq W \). In other words, \( a_{0} \, \smallsubscript{\cdot}{W} \, \vectorvar{x}_{0} \in W \), and by
    \eqref{eq:a0-cdotW-x0-equals-a0-cdot-x0-equals-a0-cdotS-x0}, it follows that \( a_{0} \, \smallsubscript{\cdot}{S} \, \vectorvar{x}_{0} \) also belongs to \( W \).
    Since \( W \) is an arbitrary element of \( X \), we have by \aref{rule-of-inference:universal-generalization}{Universal Generalization}
    \( \forall z\,{\in}\,X \, \mleft[ a_{0} \, \smallsubscript{\cdot}{S} \, \vectorvar{x}_{0} \in z \mright] \).
    By \eqref{eq:S-is-the-intersection-of-X}, this means that \( a_{0} \, \smallsubscript{\cdot}{S} \, \vectorvar{x}_{0} \in S \).
    Since \( a_{0} \in \abs{\mathcal{F}} \) and \( \vectorvar{x}_{0} \in S \) are also arbitrarily selected,
    \eqref{eq:S-is-closed-under-scalar-multiplication} can be derived immediately via \aref{rule-of-inference:universal-generalization}{Universal Generalization}.

    \vspace{0.3cm}
    \noindent
    Having established that \( \tuple{S, \smallsubscript{+}{S}} \) is an abelian group and that
    \( (\smallsubscript{\cdot}{S}) \) is a function on \( \prod \tuple{\abs{\mathcal{F}},S} \) into \( S \),
    it remains to verify that the rest of the properties of \aref{dfn:vector-space}{vector spaces} are satisfied,
    which is rather obvious since the difference between \( \smallsubscript{+}{S} \), \( \mleft( \smallsubscript{\cdot}{S} \mright) \)
    and \( + \), \( (\cdot) \) is purely symbolic---at least when it concerns the elements of \( S \).
    E.g., for any \( \vectorvar{x}, \vectorvar{y} \in S \) and \( a, b \in \abs{\mathcal{F}} \):
    \begin{gather*}
        a \, \smallsubscript{\cdot}{S} \, \mleft( b \, \smallsubscript{\cdot}{S} \, \vectorvar{x} \mright)
            = a \cdot \mleft( b \cdot \vectorvar{x} \mright)
            = \mleft( a \cdot b \mright) \cdot \vectorvar{x}
            = \mleft( a \, \smallsubscript{\cdot}{S} \, b \mright) \, \smallsubscript{\cdot}{S} \, \vectorvar{x}\\
        1 \, \smallsubscript{\cdot}{S} \, \vectorvar{x} = 1 \cdot \vectorvar{x} = \vectorvar{x}\\
        a \, \smallsubscript{\cdot}{S} \, \mleft( \vectorvar{x} \, \smallsubscript{+}{S} \, \vectorvar{y} \mright)
            = a \cdot \mleft( \vectorvar{x} + \vectorvar{y} \mright)
            = \mleft( a \cdot \vectorvar{x} \mright) + \mleft( a \cdot \vectorvar{y} \mright)
            = \mleft( a \, \smallsubscript{\cdot}{S} \, \vectorvar{x} \mright) \, \smallsubscript{+}{S} \, \mleft( a \, \smallsubscript{\cdot}{S} \, \vectorvar{y} \mright)
    \end{gather*}
    Therefore, \( \tuple{S, \mathcal{F}, \smallsubscript{+}{S}, \smallsubscript{\cdot}{S}} \) is a vector space over \( \mathcal{F} \).

    \vspace{0.3cm}
    \noindent
    Having proven claims \ref{enum:S-is-non-empty} and \ref{enum:S-is-a-vector-subspace-of-V},
    we can now conclude that \( S \in \vectorsubspaces{\mathcal{V}} \) and thus \( \Intersect X \in \vectorsubspaces{\mathcal{V}} \).
    By the \aref{rule-of-inference:deduction-theorem}{Deduction Theorem},
    \( X \subseteq \vectorsubspaces{\mathcal{V}} \,\,{\land}\,\, X \neq \emptyset {\implies} \Intersect X \in \vectorsubspaces{\mathcal{V}} \).
    Since \( X \) is arbitrary, the theorem follows from \aref{rule-of-inference:universal-generalization}{Universal Generalization}.
\end{proof}

\begin{dfn}[Span]
    \phantomsection \label{dfn:vector-span}
    Let \( \mathcal{V} \) be a vector space over a field \( \mathcal{F} \). Given any subset \( S \) of \( \abs{\mathcal{V}} \),
    the \keyword{span} of \( S \) is the intersection of all the elements of \( \vectorsubspaces{\mathcal{V}} \) that \( S \) is a
    subset of.
    \begin{equation*}
        \vectorspan{S} = \Intersect \setbuilderBig{X \subseteq \abs{\mathcal{V}}}{X \neq \emptyset \,\,{\land}\,\,
        \tuple{X, \mathcal{F}, \smallsubscript{+}{X}, \smallsubscript{\cdot}{X}} \leq \mathcal{V} \,\,{\land}\,\, S \subseteq X }
    \end{equation*}
\end{dfn}

\begin{thm}
    \phantomsection \label{thm:span-of-empty-set-is-the-singleton-containing-the-zero-vector}
    For any vector space \( \mathcal{V} \), \( \vectorspan{\emptyset} = \mleft\{ \zerovector \mright\} \).
\end{thm}

\begin{proof}
    Let \( \mathcal{V} \) be a vector space over a field \( \mathcal{F} \). It follows immediately from \cref{dfn:vector-span} that,
    \begin{equation*}
        \vectorspan{\emptyset} = \Intersect \setbuilderBig{X \subseteq \abs{\mathcal{V}}}{X \neq \emptyset \,\,{\land}\,\,
        \tuple{X, \mathcal{F}, \smallsubscript{+}{X}, \smallsubscript{\cdot}{X}} \leq \mathcal{V} \,\,{\land}\,\, \emptyset \subseteq X}
    \end{equation*}
    By \cref{thm:empty-set-universal-subset}, we know that \( \emptyset \subseteq X \) is trivially true for any set \( X \). Therefore,
    \begin{equation*}
        \vectorspan{\emptyset} = \Intersect \setbuilderBig{X \subseteq \abs{\mathcal{V}}}{X \neq \emptyset \,\,{\land}\,\,
        \tuple{X, \mathcal{F}, \smallsubscript{+}{X}, \smallsubscript{\cdot}{X}} \leq \mathcal{V}}
        = \Intersect \vectorsubspaces{\mathcal{V}}
    \end{equation*}
    But \( \Intersect \vectorsubspaces{\mathcal{V}} = \mleft\{ \zerovector \mright\} \) by
    \cref{thm:the-insersection-of-the-set-of-all-vector-subspaces-is-the-singleton-with-only-empty-set-as-its-element},
    so necessarily \( \vectorspan{\emptyset} = \mleft\{ \zerovector \mright\} \).
\end{proof}

\begin{thm}
    \phantomsection \label{thm:span-of-finite-set-of-vectors-is-equal-to-the-set-of-all-possible-linear-combinations-of-the-elements-of-the-set}
    If \( \mathcal{V} \) is a vector space over \( \mathcal{F} \) and \( S \) is a finite, non-empty subset of \( \abs{\mathcal{V}} \), then
    \( \vectorspan{S} \) is the set of all possible \aref{dfn:vector-space-linear-combination}{linear combinations} of the elements of \( S \).
    More specifically, if \( S = \mleft\{ \vectorvar{s}_{0}, \dots, \vectorvar{s}_{k-1} \mright\} \) for some \( k \in \naturalset \) and \( k \neq 0 \)
    where \( \vectorvar{s}_{i} \in \abs{\mathcal{V}} \) for all \( i \in k \), then,
    \begin{equation*}
        \vectorspan{S} = \setbuilderBig{\vectorvar{x}}{
            \exists a_{0}\,{\in}\,\abs{\mathcal{F}} \, \dots \exists a_{k-1}\,{\in}\,\abs{\mathcal{F}} \,
            \big( \vectorvar{x} = a_{0}\vectorvar{s}_{0} + \cdots + a_{k}\vectorvar{s}_{k-1} \big)
        }
    \end{equation*}
\end{thm}

\begin{proof}
    Let \( \mathcal{V} \) be a vector space over a field \( \mathcal{F} \) and \( S = \mleft\{ \vectorvar{s}_{0}, \dots, \vectorvar{s}_{k-1} \mright\} \)
    be a finite, non-empty subset of \( \abs{\mathcal{V}} \) with \( k \) elements where \( k \in \naturalset \) and \( k \neq 0 \). Let \( L \) be the set
    of all \aref{dfn:vector-space-linear-combination}{linear combinations} of \( \vectorvar{s}_{0}, \dots, \vectorvar{s}_{k-1} \). That is,
    \begin{equation}
        \label{eq:L-is-the-set-of-all-linear-combinations-of-all-the-elements-of-S}
        L = \setbuilderBig{\vectorvar{x}}{
            \exists a_{0}\,{\in}\,\abs{\mathcal{F}} \, \dots \exists a_{k-1}\,{\in}\,\abs{\mathcal{F}} \,
            \big( \vectorvar{x} = a_{0}\vectorvar{s}_{0} + \cdots + a_{k}\vectorvar{s}_{k-1} \big)}
    \end{equation}
    The goal, then, is to prove that \( \vectorspan{S} = L \), which can be achieved by showing that \( \vectorspan{S} \subseteq L \,\,{\land}\,\, L \subseteq \vectorspan{S} \).

    \vspace{0.5cm}
    \noindent
    {\scshape Claim:} \( \vectorspan{S} \subseteq L \)

    \vspace{0.1cm}
    \noindent
    It is obvious from \eqref{eq:L-is-the-set-of-all-linear-combinations-of-all-the-elements-of-S} that \( L \subseteq \abs{\mathcal{V}} \). Moreover,
    if \( 0 \) denotes the identity of \( \abs{\mathcal{F}} \) with respect to scalar addition of \( \mathcal{F} \) and
    \( \zerovector \) denotes the identity of \( \abs{\mathcal{V}} \) with respect to the vector addition of \( \mathcal{V} \), then,
    \begin{gather*}
        0 \cdot \vectorvar{s}_{0} + \cdots + 0 \cdot \vectorvar{s}_{k-1} = \zerovector
    \end{gather*}
    Since \( 0 \in \abs{\mathcal{F}} \), then it follows immediately from \eqref{eq:L-is-the-set-of-all-linear-combinations-of-all-the-elements-of-S}
    that \( \zerovector \in L \). Not only that, but each vector in \( S \) belongs to \( L \) as well since
    \( \vectorvar{s}_{i} = 0\cdot\vectorvar{s}_{0} + \cdots + 1\cdot\vectorvar{s}_{i} + \cdots + 0\cdot\vectorvar{s}_{k-1} \) for any \( i \in k \).
    Therefore, \( L \neq \emptyset \) and \( S \subseteq L \).

    \vspace{0.3cm}
    \noindent
    Let us now consider the structure \( \tuple{L, \mathcal{F}, \smallsubscript{+}{L}, \smallsubscript{\cdot}{L}} \) where \( \smallsubscript{+}{L} \)
    and \( \smallsubscript{(\cdot)}{L} \) are the vector addition and scalar multiplication of \( \mathcal{V} \)
    \aref{dfn:restricted-functions}{restricted} to \( L^{2} \) and \( \prod \tuple{\abs{\mathcal{F}}, L} \) respectively.
    Suppose that \( \vectorvar{x} \) and \( \vectorvar{y} \) are arbitrary elements of \( L \).
    Thus, we have by \eqref{eq:L-is-the-set-of-all-linear-combinations-of-all-the-elements-of-S},
    \begin{gather*}
            \vectorvar{x} = a_{0}\vectorvar{s}_{0} + \cdots + a_{k-1}\vectorvar{s}_{k-1}\\
            \vectorvar{y} = b_{0}\vectorvar{s}_{0} + \cdots + a_{k-1}\vectorvar{s}_{k-1}
    \end{gather*}
    for some \( a_{0}, \dots, a_{k-1} \in \abs{\mathcal{F}} \) and \( b_{0}, \dots, b_{k-1} \in \abs{\mathcal{F}} \).
    If we then compute \( \vectorvar{x} \, \smallsubscript{+}{L} \, \vectorvar{y} \), we obtain the following:
    \begin{align}
        \phantomsection \label{eq:x-plusL-y}
        \vectorvar{x} \, \smallsubscript{+}{L} \, \vectorvar{y}
        &= \big( a_{0}\vectorvar{s}_{0} + \cdots + a_{k-1}\vectorvar{s}_{k-1} \big) \, \smallsubscript{+}{L} \,
          \big( b_{0}\vectorvar{s}_{0} + \cdots + a_{k-1}\vectorvar{s}_{k-1} \big) \notag \\
        &= \big( a_{0}\vectorvar{s}_{0} + \cdots + a_{k-1}\vectorvar{s}_{k-1} \big) +
          \big( b_{0}\vectorvar{s}_{0} + \cdots + a_{k-1}\vectorvar{s}_{k-1} \big)
    \end{align}
    Since \( \tuple{\abs{\mathcal{V}}, +} \) is a group, the generalized associative and commutative laws are applicable
    and thus the terms in \eqref{eq:x-plusL-y} can be rearranged to:
    \begin{align}
        \label{eq:x-plusL-y-rearranged}
        \vectorvar{x} \, \smallsubscript{+}{L} \, \vectorvar{y}
         = \big( a_{0}\vectorvar{s}_{0} + b_{0}\vectorvar{s}_{0} \big) + \cdots + \big( a_{k-1}\vectorvar{s}_{k-1} + b_{k-1}\vectorvar{s}_{k-1} \big)
    \end{align}
    But \( \mathcal{V} \) is a \aref{dfn:vector-space}{vector space} over \( \mathcal{F} \), which means that scalar multiplication is distributive
    over vector addition. This allows \eqref{eq:x-plusL-y-rearranged} to be simplified even further:
    \begin{align}
        \label{eq:x-plusL-y-final}
        \vectorvar{x} \, \smallsubscript{+}{L} \, \vectorvar{y}
          = \mleft( a_{0} + b_{0} \mright)\vectorvar{s}_{0} + \cdots + \mleft( a_{k-1} + b_{k-1} \mright)\vectorvar{s}_{k-1}
    \end{align}
    Since \( \mathcal{F} \) is a field, \( \abs{\mathcal{F}} \) is closed under scalar addition by the \aref{dfn:field}{field axioms}.
    In other words, \( (a_{i} + b_{i}) \in \abs{\mathcal{F}} \) for all \( i \in k \). It should now be clear from \eqref{eq:x-plusL-y-final}
    that \( (\vectorvar{x} \, \smallsubscript{+}{L} \, \vectorvar{y}) \in L \) for any \( \vectorvar{x}, \vectorvar{y} \in L \). Simply put,
    \( L \) is \aref{dfn:closed-operation}{closed under} \( \smallsubscript{+}{L} \).
    Additionally, we have for every \( \vectorvar{x}, \vectorvar{y}, \vectorvar{z} \in L \):
    \begin{gather*}
        \vectorvar{x} \, \smallsubscript{+}{L} \, \mleft( \vectorvar{y} \, \smallsubscript{+}{L} \, \vectorvar{z} \mright)
        = \vectorvar{x} + \mleft( \vectorvar{y} + \vectorvar{z} \mright)
        = \mleft( \vectorvar{x} + \vectorvar{y} \mright) + \vectorvar{z}
        = \mleft( \vectorvar{x} \, \smallsubscript{+}{L} \, \vectorvar{y} \mright) \, \smallsubscript{+}{L} \, \vectorvar{z}\\
        \vectorvar{x} \, \smallsubscript{+}{L} \, \vectorvar{y}
        = \vectorvar{x} + \vectorvar{y}
        = \vectorvar{y} + \vectorvar{x}
        = \vectorvar{y} \, \smallsubscript{+}{L} \, \vectorvar{x}
    \end{gather*}
    To summarize, it has been shown that:
    \begin{enumerate*}[label=(\roman*)]
        \item \( \zerovector \in L \),
        \item \( \zerovector \) is the identity of \( L \) with respect to \( \smallsubscript{+}{L} \),
        \item \( L \) is closed under \( \smallsubscript{+}{L} \),
        \item \( \smallsubscript{+}{L} \) is associative,
        \item \( \smallsubscript{+}{L} \) is commutative.
    \end{enumerate*}
    In other words, the structure \( \tuple{L, \smallsubscript{+}{L}} \) is an abelian group.

    \vspace{0.3cm}
    \noindent
    Next, we verify that \( \ran (\smallsubscript{\cdot}{L}) \subseteq L \).
    Let \( \alpha \in \abs{\mathcal{F}} \) and \( \vectorvar{x} \in L \) be arbitrary.
    Thus,
    \begin{align}
        \phantomsection \label{eq:alpha-cdotL-x-undistributed}
        \alpha \, \smallsubscript{\cdot}{L} \, \vectorvar{x}
        = \alpha \cdot \vectorvar{x}
        = \alpha \cdot \big( a_{0}\vectorvar{s}_{0} + \cdots + a_{k-1}\vectorvar{s}_{k-1} \big)
    \end{align}
    By the generalized distributive law, \eqref{eq:alpha-cdotL-x-undistributed} can then be reduced to,
    \begin{equation}
        \alpha \, \smallsubscript{\cdot}{L} \, \vectorvar{x} =
        \mleft( \alpha a_{0} \mright)\vectorvar{s}_{0} + \cdots + \mleft( \alpha a_{k-1} \mright)\vectorvar{s}_{k-1}
    \end{equation}
    Since \( \mathcal{F} \) is a field, it follows that \( \alpha_{i}a_{i} \in \abs{\mathcal{F}} \) for all \( i \in k \). This means that
    \( \alpha \, \smallsubscript{\cdot}{L} \, \vectorvar{x} \) is also a \aref{dfn:vector-space-linear-combination}{linear combination}
    of \( \vectorvar{s}_{0}, \dots, \vectorvar{s}_{k-1} \).
    By \eqref{eq:L-is-the-set-of-all-linear-combinations-of-all-the-elements-of-S}, it follows that \( \alpha \, \smallsubscript{\cdot}{L} \, \vectorvar{x} \in L \).
    Thus, \( \ran (\smallsubscript{\cdot}{L}) \subseteq L \).
    With that, we can now verify that the structure \( \tuple{L, \mathcal{F}, \smallsubscript{+}{L}, \smallsubscript{\cdot}{L}} \) satisfies all the other properties of
    \aref{dfn:vector-space}{vector spaces}.
    Consider, for instance, the operation \( a \, \smallsubscript{\cdot}{L} \, \mleft( b \, \smallsubscript{\cdot}{L} \, \vectorvar{x} \mright) \)
    where \( a, b \in \abs{\mathcal{F}} \) and \( \vectorvar{x} \in L \).
    Since \( \ran (\smallsubscript{\cdot}{L}) \subseteq L \), it follows that \( b \, \smallsubscript{\cdot}{L} \, \vectorvar{x} \in L \).
    Since \( \dom (\smallsubscript{\cdot}{L}) = \prod \tuple{\abs{\mathcal{F}}, L} \) and \( \smallsubscript{(\cdot)}{L} \) is a restriction of \( (\cdot) \),
    it follows that \( a \, \smallsubscript{\cdot}{L} \, \mleft( b \, \smallsubscript{\cdot}{L} \, \vectorvar{x} \mright) \) is defined and that,
    \begin{equation*}
        a \, \smallsubscript{\cdot}{L} \, \mleft( b \, \smallsubscript{\cdot}{L} \, \vectorvar{x} \mright)
        = a \cdot \mleft( b \cdot \vectorvar{x} \mright)
        = \mleft( ab \mright) \cdot \vectorvar{x}
        = \mleft( ab \mright) \, \smallsubscript{\cdot}{L} \, \vectorvar{x}
    \end{equation*}
    Similarly, \( 1 \, \smallsubscript{\cdot}{L} \, \vectorvar{x} =  1 \cdot \vectorvar{x} = \vectorvar{x} \). The distributivity
    of \( \smallsubscript{(\cdot)}{L} \) over vector and scalar additions of \( \mathcal{V} \) also holds:
    \begin{gather*}
        a \, \smallsubscript{\cdot}{L} \, \mleft( \vectorvar{x} \, \smallsubscript{+}{L} \, \vectorvar{y} \mright)
        = a \cdot \mleft( \vectorvar{x} + \vectorvar{y} \mright)
        = a \cdot \vectorvar{x} + a \cdot \vectorvar{y}
        = a \, \smallsubscript{\cdot}{L} \vectorvar{x} \, \smallsubscript{+}{L} \, a \, \smallsubscript{\cdot}{L} \, \vectorvar{y}\\
        \mleft( a + b \mright) \, \smallsubscript{\cdot}{L} \, \vectorvar{x}
        = \mleft( a + b \mright) \cdot \vectorvar{x}
        = a \cdot \vectorvar{x} + b \cdot \vectorvar{x}
        = a \, \smallsubscript{\cdot}{L} \, \vectorvar{x} \, \smallsubscript{+}{L} \, b \, \smallsubscript{\cdot}{L} \, \vectorvar{x}
    \end{gather*}
    Thus, \( \tuple{L, \mathcal{F}, \smallsubscript{+}{L}, \smallsubscript{\cdot}{L}} \) is a \aref{dfn:vector-space}{vector space} over \( \mathcal{F} \)
    and a \aref{dfn:vector-subspace}{vector subspace} of \( \mathcal{V} \). In other words, \( L \in \vectorsubspaces{\mathcal{V}} \).
    Since \( \vectorspan{S} \) is the intersection of all the subsets of \( S \) that are in \( \vectorsubspaces{\mathcal{V}} \), it follows that
    \( \vectorspan{S} \subseteq L \).

    \vspace{0.3cm}
    \noindent
    {\scshape Claim:} \( L \subseteq \vectorspan{S} \)

    \vspace{0.1cm}
    \noindent
    Let \( \vectorvar{x} = a_{0}\vectorvar{s}_{0} + \cdots + a_{k-1}\vectorvar{s}_{k-1} \) be an arbitrary element of \( L \) where
    \( a_{0},\dots,a_{k-1} \in \abs{\mathcal{F}} \).
    Suppose that \( W \) is an arbitrary element of \( \vectorsubspaces{\mathcal{V}} \) and that \( S \subseteq W \).
    By \cref{dfn:set-of-all-subsets-that-form-vector-subspaces}, it follows that \( \tuple{W, \mathcal{F}, \smallsubscript{+}{W}, \smallsubscript{\cdot}{W}} \)
    is a \aref{dfn:vector-subspace}{vector subspace} of \( \mathcal{V} \).
    Since \( S \subseteq W \), it follows that \( \vectorvar{s}_{0}, \dots, \vectorvar{s}_{k-1} \in W \).
    Since \( W \) is closed under \( \smallsubscript{+}{W} \) and \( \ran (\smallsubscript{\cdot}{W}) \subseteq W \), it follows that
    \( \vectorvar{x} = a_{0}\vectorvar{s}_{0} + \cdots + a_{k-1}\vectorvar{s}_{k-1} \in W \).
    Thus, by the \aref{rule-of-inference:deduction-theorem}{Deduction Theorem},
    \begin{equation*}
        W \in \vectorsubspaces{\mathcal{V}} \,\,{\land}\,\, S \subseteq W \implies \vectorvar{x} \in W
    \end{equation*}
    Since \( W \) is arbitrary, we have by \aref{rule-of-inference:universal-generalization}{Universal Generalization},
    \begin{equation*}
        \forall Z \, \big[ Z \in \vectorsubspaces{\mathcal{V}} \,\,{\land}\,\, S \subseteq Z \implies \vectorvar{x} \in Z \big]
    \end{equation*}
    By \cref{dfn:vector-span}, we have \( \vectorvar{x} \in \vectorspan{S} \). And finally, by the \aref{rule-of-inference:deduction-theorem}{Deduction Theorem}
    and \aref{rule-of-inference:universal-generalization}{Universal Generalization},
    \begin{equation*}
        \forall \vectorvar{x} \, \big[ \vectorvar{x} \in L \implies \vectorvar{x} \in \vectorspan{S} \big]\\
    \end{equation*}
    In other words, \( L \subseteq \vectorspan{S} \).

    \vspace{0.3cm}
    \noindent
    Having proven both \( \vectorspan{S} \subseteq L \) and \( L \subseteq \vectorspan{S} \), the equality \( \vectorspan{S} = L \) follows immediately from
    \cref{thm:set-equality-in-terms-of-subsets}.
\end{proof}

\begin{dfn}[Basis]
    \phantomsection \label{dfn:vector-space-basis}
    Let \( \mathcal{V} \) be a vector space. A set \( B \) is a \keyword{basis} of \( \mathcal{V} \) if and only if
    the following conditions hold:
    \begin{enumerate}[topsep=2pt,noitemsep,label=\textup{\texttt{(\alph*)}}]
        \item \( B \subseteq \abs{\mathcal{V}} \)
        \item \aref{notation:linear-dependence-and-independence}{\( {\nparallel}B \)}
        \item \aref{dfn:vector-span}{\( \vectorspan{B} = \abs{\mathcal{V}} \)}
    \end{enumerate}
    \vspace{0.1cm}
    \noindent
    If the basis of \( \mathcal{V} \) is a finite set, then \( \mathcal{V} \) is called a \keyword{finite-dimensional} vector space; otherwise,
    \( \mathcal{V} \) is an \keyword{infinite dimensional} vector space.
\end{dfn}

\noindent
It can be shown that every vector space has a basis although doing so requires the use of Zorn's Lemma, which we have yet to cover. The interested reader
should consult a text on linear algebra for a more in-depth treatment of vector spaces.

\subsubsection{Field of Sets} \phantomsection \label{sec:field-of-sets}
\begin{dfn}
    \phantomsection \label{dfn:field-of-sets}
    If \( X \) is a non-empty set, then the structure \( \tuple{X, F} \) is a \keyword{field of sets} if and only if the following conditions hold:
    \begin{enumerate}[topsep=3pt,noitemsep,label=\textup{\texttt{(\alph*)}}]
        \item \( F \subseteq \powerset{X} \) \phantomsection \label{enum:field-of-sets-F-is-a-system-of-subsets-of-X}
        \item \( \emptyset \in F \) \phantomsection \label{enum:field-of-sets-contains-empty-set}
        \item \( \forall x\,{\in}\,F \, \forall y\,{\in}\,F \, \mleft[ x - y \in F \mright] \) \phantomsection \label{enum:field-of-sets-closure-under-complementation}
        \item \( \forall x\,{\in}\,F \, \forall y\,{\in}\,F \, \mleft[ x \union y \in F \mright] \) \phantomsection \label{enum:field-of-sets-closure-under-binary-union}
    \end{enumerate}
    The set \( F \) is sometimes called an \keyword{algebra} over \( X \). The elements of \( X \) are called \keyword{points} whereas the elements
    of \( F \) are called \keyword{complexes}.
\end{dfn}

\begin{rem}
    The term `field' in \cref{dfn:field-of-sets} is a misnomer since has nothing to do with fields in the sense of \aref{dfn:field}{commutative division rings}.
    In fact, if \( \tuple{X, F} \) is a field of sets, then \( \tuple{F, \symdiff, \intersect} \) cannot possibly be a \aref{dfn:ring-division}{division ring}
    because \( F \) does not possess an identity with respect to \( \intersect \); i.e., \( \forall x\,{\in}\,F \, \mleft( x \intersect x = x \mright) \).
    Rather than a field, \( \tuple{F, \symdiff, \intersect} \) is actually a \aref{dfn:ring-rings-of-sets}{ring of sets} as evident from
    \cref{dfn:ring-rings-of-sets}; more precisely, the structure \( \tuple{F, \symdiff, \intersect} \) is a \aref{dfn:ring-boolean}{Boolean ring}.
\end{rem}

\begin{dfn}[Sigma Field]
    \phantomsection \label{dfn:sigma-field}
    The structure \( \tuple{X, \Sigma} \) is a \keyword{sigma field} or \keyword{measurable space} if and only if the following conditions hold:
    \begin{enumerate}[topsep=3pt,noitemsep,label=\textup{\texttt{(\alph*)}}]
        \item \( \tuple{X, \Sigma} \) is a field of sets.
        \item \( \Sigma \) is closed under countable\footnote{The countability of a set refers to the possibility of defining an injective
            function between the set and \( \naturalset \). Intuitively, it means that the ``number of elements'' in the set does not exceed
            that of the natural numbers.}\,unions. In other words,
            \begin{equation*}
                \forall X\,{\in}\,\powerset{\Sigma} \, \Big[ \mleft( \abs{X} \leq \abs{\naturalset} \mright) \implies \Union X \in \Sigma \Big]
            \end{equation*}
    \end{enumerate}
    The set \( \Sigma \) is called a \keyword{sigma algebra} over \( X \).
\end{dfn}

\begin{rem}
    It is worth emphasizing that the second condition of \cref{dfn:sigma-field} applies also to subsets of \( \Sigma \) that are
    \emph{countably infinite}---having the same cardinality as \( \naturalset \).
    In fact, the \( < \) part of the condition is actually redundant since
    the closure of \( \Sigma \) under finite unions is a direct
    consequence\footnote{The proof for this requires the principle of mathematical induction, which we will cover in the next chapter.}
    of the fact that \( \tuple{X, \Sigma} \) is a field of sets.
\end{rem}

\subsection{Metric Structures}

Metric structures are structures that endow a set with the notion of `geometric distance' between its elements.

\begin{dfn}[Metric]
    \phantomsection \label{dfn:metric-space}
    If \( M \) is a non-empty set and \( \function{d}{M^{2}}{\realset} \),
    then the structure \( \tuple{M, d} \) is a \keyword{metric space} if and only if
    the following conditions hold:
    \begin{enumerate}[topsep=3pt,itemsep=2pt,label=\textup{\texttt{(\alph*)}}]
        \item \( \forall x\,{\in}\,M \, \mleft[ d(x,x) = 0 \mright] \)
        \item \( \forall x\,{\in}\,M \, \forall y\,{\in}\,M \, \mleft[ x \neq y {\implies} d(x,y) > 0 \mright] \) \textup{(non-negativity)}
        \item \( \forall x\,{\in}\,M \, \forall y\,{\in}\,M \, \mleft[ d(x,y) = d(y,x) \mright] \) \textup{(symmetry)}
        \item \( \forall x\,{\in}\,M \, \forall y\,{\in}\,M \, \forall z\,{\in}\,M \, \mleft[ d(x,z) \leq d(x,y) + d(y,z) \mright] \)
            \textup{(triangle inequality)}
    \end{enumerate}

    \vspace{0.1cm}
    \noindent
    The elements of \( M \) are often referred to as the \keyword{points} of \( M \) whereas the function \( d \) is commonly known as 
    a \keyword{metric} on \( M \).
\end{dfn}

\noindent
The simplest example of a metric structure is the structure \( \tuple{\realset, d} \) where the function \( d \) is defined as the absolute difference between any
two elements of \( \realset \):
\begin{equation*}
    d = \setbuilderbigg{\orderedpair{\orderedpair{x}{y}}{z} \in \realset^{2} \times \realset }{
        z = \sqrt{\mleft( x - y \mright)^{2}}
    }
\end{equation*}

\begin{dfn}[Subspace]
    \phantomsection \label{dfn:metric-subspace}
    Given any metric space \( \tuple{M, d} \), a \keyword{metric subspace} is a \aref{dfn:substructure}{substructure}
    of \( \tuple{M, d} \) that is itself a metric space.
    In other words, a structure \( \tuple{M^{\prime}, d^{\prime}} \) is a metric subspace of \( \tuple{M, d} \) if and only if:
    \begin{enumerate}[topsep=2pt,label=\textup{\texttt{(\alph*)}}]
        \item \( M^{\prime} \subseteq M \)
        \item \( d^{\prime} = \restrictedfunction{d}{(M^{\prime})^{2}} \)
    \end{enumerate}
\end{dfn}

\subsection{Topological Structures}

\begin{dfn}[Topology]
    \phantomsection \label{dfn:topology}
    A set \( \mathcal{T} \) is a \keyword{topology} on a non-empty set \( X \) if and only if the following conditions hold:
    \begin{enumerate}[topsep=2pt,noitemsep,label=\textup{\texttt{(\alph*)}}]
        \item \( \mathcal{T} \subseteq \powerset{X} \)
        \item \( \emptyset, X \in \mathcal{T} \)
        \item \( \forall S \subseteq \mathcal{T} \, \mleft( \Union S \in \mathcal{T} \mright) \)
        \item \( \forall S \subseteq \mathcal{T} \, \mleft( \abs{S} < \abs{\naturalset} {\implies} \Intersect S \in \mathcal{T} \mright) \)
    \end{enumerate}
    \vspace{0.1cm}
    \noindent
    The structure \( \tuple{X, \mathcal{T}} \) is known as a \keyword{topological space}.
\end{dfn}

\begin{cor}
    For any non-empty set \( X \), \( \powerset{X} \) and \( \mleft\{ X, \emptyset \mright\} \) are both
    topologies on \( X \).
\end{cor}

\begin{rem}
    \( \powerset{X} \) is known as the \keyword{discrete topology} on \( X \) whereas \( \mleft\{ X, \emptyset \mright\} \) is known as the \keyword{indiscrete} or
    \keyword{trivial topology} on \( X \).
\end{rem}

\begin{dfn}[Open Set]
    \phantomsection \label{dfn:topological-space-open-set}
    If \( \tuple{X, \mathcal{T}} \) is a topological space, then an element of \( \mathcal{T} \) is called an \keyword{open set} or \keyword{open subset} of \( X \).
\end{dfn}

\begin{dfn}[Closed Set]
    \phantomsection \label{dfn:topological-space-closed-set}
    If \( \tuple{X, \mathcal{T}} \) is a topological space, then a subset \( A \) of \( X \) is called a \keyword{closed set} of \( X \) if and only if
    \( X - A \) is an open set; that is, \( X - A \in \mathcal{T} \).
\end{dfn}

\begin{thm}
    \phantomsection \label{thm:topological-space-properties-of-closed-sets}
    For any topological space \( \tuple{X, \mathcal{T}} \), the following properties hold:
    \begin{enumerate}[topsep=2pt,label=\textup{\texttt{(\alph*)}}]
        \item \( \emptyset \) and \( X \) are closed sets. \phantomsection \label{enum:topological-space-properties-of-closed-sets-empty-and-X-are-closed}
        \item Arbitrary intersection of closed sets results in closed set. \phantomsection \label{enum:topological-space-properties-of-closed-sets-intersection-closed}
        \item Finite union of closed sets results in a closed set. \phantomsection \label{enum:topological-space-properties-of-closed-sets-finite-union-closed}
    \end{enumerate}
\end{thm}

\begin{proof}
    Let \( \tuple{X, \mathcal{T}} \) be a topological space. We now prove each property separately.
    
    \vspace{0.3cm}
    \noindent
    {\scshape Property} \ref{enum:topological-space-properties-of-closed-sets-empty-and-X-are-closed}: \( \emptyset \) and \( X \) are closed

    \vspace{0.1cm}
    \noindent
    Since \( \mathcal{T} \) is a topology on \( X \), it follows from \cref{dfn:topology} that \( \emptyset \in \mathcal{T} \) and \( X \in \mathcal{T} \).
    However, \( \emptyset = X - X \) and \( X = X - \emptyset \). By \cref{dfn:topological-space-closed-set}, \( \emptyset \) and \( X \) are closed sets.

    \vspace{0.3cm}
    \noindent
    {\scshape Property} \ref{enum:topological-space-properties-of-closed-sets-intersection-closed}: Arbitrary intersection of closed sets yields a closed set

    \vspace{0.1cm}
    \noindent
    Let \( S \subseteq \powerset{X} \) be an arbitrary system of subsets of \( X \) such that \( \forall x\,{\in}\,S \, \mleft[ (X - x) \in \mathcal{T} \mright] \).
    We have to show that \( \Intersect S \) is a closed set; that is, \( \mleft( X - \Intersect S \mright) \in \mathcal{T} \).
    Thus, we have by the \aref{set-identities:generalized-de-morgan-laws}{Generalized De Morgan's Law}:
    \begin{equation*}
        X - \Intersect S = \Union \setbuilderBig{x}{\exists z \, \big( z \in S \,\,{\land}\,\, x = X - z \big)}
    \end{equation*}
    Let \( M = \setbuilderbig{x}{\exists z \, \mleft( z \in S \,\,{\land}\,\, x = X - z \mright)} \). Let \( x_{0} \) be an arbitrary element of \( M \).
    Thus \(x_{0} = X - z_{0}\) for some \( z_{0} \in S \). Since \( \forall x\,{\in}\,S \, \mleft[ \mleft( X - x \mright) \in \mathcal{T} \mright] \) and
    \( z_{0} \in S \), it follows by \aref{rule-of-inference:universal-instantiation}{Universal Instantiation} that \( x_{0} = X - z_{0} \in \mathcal{T} \).
    In other words, \( x_{0} \in M {\implies} x_{0} \in \mathcal{T} \) for any arbitrary \( x_{0} \).
    By \aref{rule-of-inference:universal-generalization}{Universal Generalization}, it follows that \( M \subseteq \mathcal{T} \).
    Since \( \mathcal{T} \) is a topology on \( X \), we have by \cref{dfn:topology} \( X - \Intersect S = \Union M \in \mathcal{T} \).

    \vspace{0.3cm}
    \noindent
    {\scshape Property} \ref{enum:topological-space-properties-of-closed-sets-finite-union-closed}: Finite union of closed sets yields a closed set

    \vspace{0.1cm}
    \noindent
    Let \( S \subseteq \powerset{X} \) be an arbitrary system of subsets of \( X \)  such that
    \( \forall x\,{\in}\,S \, \mleft[ \mleft( X - x \mright) \in \mathcal{T} \mright] \) and \( \abs{S} < \abs{\naturalset} \).
    The aim is to show that \( (X - \Union S) \in \mathcal{T} \).
    Once again, we have by the \aref{set-identities:generalized-de-morgan-laws}{Generalized De Morgan's Laws}:
    \begin{equation*}
        X - \Union S = \Intersect \setbuilderBig{x}{\exists z \, \big( z \in S \,\,{\land}\,\, x = X - z \big)}
    \end{equation*}
    Let \( M = \setbuilderbig{x}{\exists z \, \mleft[ z \in S \,\,{\land}\,\, x = X - z \mright]} \) and \( x_{0} \)
    be an arbitrary element of \( M \). Thus, \( x_{0} = X - z_{0} \) for some \( z_{0} \in S \).
    Since \( \forall x\,{\in}\,S \, \mleft[ \mleft( X - x \mright) \in \mathcal{T} \mright] \) and \( z_{0} \in S \),
    we have by \aref{rule-of-inference:universal-instantiation}{Universal Instantiation} \( X - z_{0} \in \mathcal{T} \).
    Thus, \( x_{0} \in \mathcal{T} \). By the \aref{rule-of-inference:deduction-theorem}{Deduction Theorem},
    \( x_{0} \in M {\implies} x_{0} \in \mathcal{T} \) for any arbitrary \( x_{0} \).
    In other words, \( M \subseteq \mathcal{T} \). Since \( S \) is a finite set, \( M \) must also be finite. Thus,
    \( X - \Union S = \Intersect M \in \mathcal{T} \) by \cref{dfn:topology}.
\end{proof}

\begin{rem}
    One could very well use closed sets and the three properties stated in \cref{thm:topological-space-properties-of-closed-sets} to define what
    a topology is; however, due to historical reasons, most mathematicians prefer the open set definition.
\end{rem}

\begin{dfn}[Interior]
    \phantomsection \label{dfn:topological-space-interior-of-a-subset}
    Let \( \tuple{X, \mathcal{T}} \) be a topological space and \( A \) be a subset of \( X \). The \keyword{interior} of \( A \) is defined as
    the union of all open sets contained in \( A \). In other words,
    \begin{equation*}
        \forall x \, \Big[ x \in \interior_{\tuple{X, \mathcal{T}}}{A} \iff \big( x \in X \,\,{\land}\,\,
        \exists y\,{\in}\,\mathcal{T} \, \mleft[ y \subseteq A \mright] \big) \Big]
    \end{equation*}
    Or in \aref{dfn:set-builder-notation}{set-builder notation}, \( \interior_{\tuple{X, \mathcal{T}}}{A} = \Union \setbuilderbig{x \in \mathcal{T}}{x \subseteq A} \).
\end{dfn}

\begin{cor}
    \phantomsection \label{cor:topological-space-interior-is-an-openset}
    Since \( \setbuilderbig{x \in \mathcal{T}}{x \subseteq A} \subseteq \mathcal{T} \), it follows from
    \cref{dfn:topology} that its union belongs to \( \mathcal{T} \).
    In other words, for any \( A \subseteq X \), \( \interior_{\tuple{X, \mathcal{T}}}{A} \) is a open set.
\end{cor}

\begin{dfn}[Closure]
    \phantomsection \label{dfn:topological-space-closure}
    Let \( \tuple{X, \mathcal{T}} \) be a topological space and \( A \) be a subset of \( X \). The \keyword{closure}
    of \( A \) is defined as the intersection of all closed sets containing \( A \). That is,
    \begin{equation*}
        \forall x \, \Big[ x \in \closure_{\tuple{X, \mathcal{T}}}{A}
            \iff
            \forall y\,{\in}\,\powerset{X} \, \big[ \mleft(  X - y \in \mathcal{T} \mright) \,\,{\land}\,\, \mleft( A \subseteq y \mright) \implies \mleft( x \in y \mright) \big]
        \Big]
    \end{equation*}
    In \aref{dfn:set-builder-notation}{set-builder notation},
    \( \closure_{\tuple{X, \mathcal{T}}}{A} = \Intersect \setbuilderbig{x \in \powerset{X}}{X - x \in \mathcal{T} \,\,{\land}\,\, A \subseteq x} \).
\end{dfn}

\begin{notation}
    \phantomsection \label{notation:topological-space-closure-and-interior}
    When \( \mathcal{T} \) is understood to be a topology on \( X \),
    the notations for the closure and interior of \( A \) can be simplified to
    \( \closure_{X}{A} \) and \( \interior_{X}{A} \).
\end{notation}

\begin{thm}
    \phantomsection \label{thm:topological-space-closure-of-a-subset-is-a-closed-set}
    If \( \tuple{X, \mathcal{T}} \) is a topological space and \( A \) is a subset of \( X \), then
    \( \closure_{X}{A} \) is a closed set.
\end{thm}

\begin{proof}
    By \cref{dfn:topological-space-closure}, we have,
    \begin{equation*}
        \closure_{X}{A} = \Intersect \setbuilderbig{x \in \powerset{X}}{X - x \in \mathcal{T} \,\,{\land}\,\, A \subseteq x}
    \end{equation*}
    Now let \( S =  \setbuilderbig{x \in \powerset{X}}{X - x \in \mathcal{T} \,\,{\land}\,\, A \subseteq x} \)
    so that \( \closure_{X}{A} = \Intersect S \).
    Furthermore, let \( R = \setbuilderbig{x}{\exists z \, \mleft( z \in S \,\,{\land}\,\, x = X - z \mright)} \).
    Thus, by the \aref{set-identities:generalized-de-morgan-laws}{Generalized De Morgan's Law},
    \begin{equation*}
        X - \Intersect S = \Union R 
    \end{equation*}
    Let \( x_{0} \in R \) be arbitrary. Thus, \( x_{0} = X - z_{0} \) for some \( z_{0} \in \powerset{X} \) where \( X - z_{0} \in \mathcal{T} \) and \( A \subseteq z_{0}\).
    Clearly, \( x_{0} \in \mathcal{T} \), which means that \( \forall x\,{\in}\,R \, \mleft( x \in \mathcal{T} \mright) \), or simply \( R \subseteq \mathcal{T} \).
    By \cref{dfn:topology}, it thus follows that \( \Union R = X - \Intersect S = X - \closure_{X}{A} \) is an open set.
    Finally, by \cref{dfn:topological-space-closure}, it follows that \( \closure_{X}{A} \) is closed.
\end{proof}

\begin{thm}
    \phantomsection \label{thm:topological-space-inclusion-relations-between-interior-and-closure}
    If \( \tuple{X, \mathcal{T}} \) is a topological space and \( A \) is a subset of \( X \), then,
    \begin{equation*}
        \interior{A} \subseteq A \subseteq \closure{A}
    \end{equation*}
\end{thm}

\begin{proof}
    It is obvious from \cref{dfn:topological-space-interior-of-a-subset} that \( \interior{A} \subseteq A \) since
    \( \interior{A} \) is the union of a system of subsets of \( X \) that \( A \) is a subset of (every element of
    \( \interior{A} \) must therefore also belong to \( A \)).
    Now let \( a \in A \) be arbitrary. Additionally, let \( Z \) be a subset of \( X \) such that \( X - Z \in \mathcal{T} \) and \( A \subseteq Z \).
    Since \( A \subseteq Z \) and \( a \in A \), clearly \( a \in Z \).
    Since \( Z \) is arbitrary, we have by \aref{rule-of-inference:universal-generalization}{Universal Generalization},
    \begin{equation*}
        \forall Z\,{\in}\,\powerset{X} \, \Big[ (X - Z \in \mathcal{T}) \,\,{\land}\,\, \mleft( A \subseteq Z \mright) \implies a \in Z \Big]
    \end{equation*}
    Thus, \( a \in \closure{A} \) by \cref{dfn:topological-space-closure}. And finally, by the \aref{rule-of-inference:deduction-theorem}{Deduction Theorem}, we have,
    \begin{equation*}
        a \in A \implies a \in \closure{A}
    \end{equation*}
    Since \( a \in A \) is arbitrary, \aref{rule-of-inference:universal-generalization}{Universal Generalization} applies and we have
    \( \forall x\,{\in}\,A \, \big[ x \in A {\implies} x \in \closure{A} \big] \). In other words, \( A \subseteq \closure{A} \)
\end{proof}

\noindent
Many notions from calculus and analysis can also be generalized using topological spaces.

\begin{dfn}[Connectedness]
    \phantomsection \label{dfn:topological-space-connectedness}
    A topological space \( \tuple{X, \mathcal{T}} \) is \keyword{disconnected} if and only if \( X \) is the union of two disjoint open sets.
    In other words,
    \begin{equation*}
        \exists A\,{\in}\,\mathcal{T} \, \exists B\,{\in}\,\mathcal{T} \, \big[ (A \intersect B = \emptyset) \,\,{\land}\,\, (X = A \union B) \big]
    \end{equation*}
    A topological space is \keyword{connected} if and only if it is not disconnected.
\end{dfn}

\begin{dfn}[Continuity]
    \phantomsection \label{dfn:topological-space-continuity}
    If \( \tuple{X, \mathcal{T}_{X}} \) and \( \tuple{Y, \mathcal{T}_{Y}} \) are two topological spaces, then a function
    \( \function{f}{X}{Y} \) is \keyword{continuous} if and only if for all \( S \in \mathcal{T}_{Y} \),
    \( \imageSet{f^{-1}}{S} \in \mathcal{T}_{X} \). Formally,
    \begin{equation*}
        \forall S\,{\in}\,\mathcal{T}_{Y} \, \big[ \imageSet{f^{-1}}{S} \in \mathcal{T}_{X} \big]
    \end{equation*}
\end{dfn}

\subsection{Measures}
\begin{dfn}[Set Function]
    \phantomsection \label{dfn:set-function}
    A function \( f \) is a \keyword{set function} if and only if \( \dom f \subseteq \powerset{X} \) for some given set \( X \).
\end{dfn}

\begin{dfn}
    A set function \( f \) is \keyword{algebraic} if and only if \( \ran f \subseteq \abs{\mathcal{F}} \) for some \aref{dfn:field}{field} \( \mathcal{F} \).
\end{dfn}

\begin{dfn}[Additivity]
    \phantomsection \label{dfn:additive-set-function}
    An algebraic set function \( f \) is \keyword{additive} if and only if for any \( A, B \in \dom f \) such that \( A \) and \( B \)
    are disjoint, \( A \union B \in \dom f \) and \( f(A \union B) = f(A) + f(B) \).
\end{dfn}

\begin{dfn}[Finite Additivity]
    \phantomsection \label{dfn:finite-additive-set-function}
    An algebraic set function \( f \) is \keyword{finitely additive} if and only if for any \( S \subseteq \dom f \)
    such that \( \exists k\,{\in}\,\naturalset \, \big( \abs{S} = k \big) \) and all elements of \( S \)
    are \aref{dfn:pairwise-disjoint-sets}{pairwise disjoint}:
    \begin{enumerate}[topsep=2pt,label=\textup{\texttt{(\alph*)}}]
        \item \( \Union S \in \dom f \)
        \item \( f \mleft( \Union S \mright) = \sum_{X \,\in\, S} f(X) \)
    \end{enumerate}
\end{dfn}

\begin{rem}
    It can be shown by mathematical induction that finite additivity and (binary) additivity are in fact equivalent.
\end{rem}

\begin{thm}
    \phantomsection \label{thm:additive-set-function-maps-empty-set-to-0}
    Given any field \( \mathcal{F} \) and any additive algebraic set function \( f \) with \( \ran f \subseteq \abs{\mathcal{F}} \),
    \( \emptyset \in \dom f {\implies} f(\emptyset) = 0 \) where \( 0 \) denotes the additive identity of \( \mathcal{F} \).
\end{thm}

\begin{proof}
    Let \( \mathcal{F} \) be a field with additive identity \( 0 \). Let \( f \) be an algebraic set function with \( \ran f \subseteq \abs{\mathcal{F}} \).
    Suppose that \( \emptyset \in \dom f \). Now let \( A \in \dom f \) be arbitrary. Clearly, \( A \union \emptyset = \emptyset \union A = A \).
    Therefore, \( f(A \union \emptyset) = f(\emptyset \union A) = f(A) \). But \( f \) is also additive, so we have by \cref{dfn:additive-set-function}
    \( f(A \union \emptyset) = f(A) + f(\emptyset) \). In other words, \( f(A) = f(A) + f(\emptyset) \). Since \( \ran f \subseteq \abs{\mathcal{F}} \)
    and \( \mathcal{F} \) is a field, it is obvious that \( f(\emptyset) \) is the additive identity of \( \mathcal{F} \). Thus, \( f(\emptyset) = 0 \).
\end{proof}

\begin{dfn}[Sigma Additivity]
    \phantomsection \label{dfn:sigma-additive-set-function}
    An algebraic set function \( f \) is \keyword{sigma additive} if and only if for any
    countable subset \( S \) of \( \dom f \) consisting of only \aref{dfn:pairwise-disjoint-sets}{pairwise disjoint} elements:
    \begin{enumerate}[topsep=2pt,noitemsep,label=\textup{\texttt{(\alph*)}}]
        \item \( \Union S \in \dom f \)
        \item \( f \mleft( \Union S \mright) = \sum_{x \,\in\, S} f(x) \)
    \end{enumerate}
\end{dfn}

\begin{rem}
    Obviously, sigma additivity implies additivity.
\end{rem}

\begin{dfn}[Measure Space]
    \phantomsection \label{dfn:measure-space}
    For any set \( X \), the structure \( \tuple{X, \Sigma, \mu} \) is a \keyword{measure space} if and only if
    the following conditions hold:
    \begin{enumerate}[topsep=2pt,noitemsep,label=\textup{\texttt{(\alph*)}}]
        \item \( \Sigma \) is a \aref{dfn:sigma-field}{sigma algebra} over \( X \)
        \item \( \function{\mu}{\Sigma}{\realset} \) is an algebraic set function with the following properties:
            \begin{itemize}[topsep=2pt]
                \item \( \mu(x) \geq 0 \) for all \( x \in \Sigma \)
                \item \( \mu \) is sigma additive
            \end{itemize}
    \end{enumerate}
    The set function \( \mu \) is called a \keyword{measure} on \( \Sigma \).
\end{dfn}

\begin{thm}[Monotonicity]
    \phantomsection \label{thm:measure-space-monotonicity}
    If \( \tuple{X, \Sigma, \mu} \) is a measure space, then for all \( A, B \in \Sigma \), \( A \subseteq B {\implies} \mu(A) \leq \mu(B) \).
\end{thm}

\begin{proof}
    Let \( \tuple{X, \Sigma, \mu} \) be a measure space. Suppose that \( A, B \in \Sigma \) and that \( A \subseteq B \).
    Since \( \Sigma \) is a \aref{dfn:sigma-field}{sigma algebra} over \( X \), it follows that \( \tuple{X, \Sigma} \) is
    a field of sets, which means that the following statements must hold:
    \begin{enumerate}[topsep=2pt,label=\textup{\texttt{(\arabic*)}}]
        \item \( B - A \in \Sigma \)
        \item \( A \union B \in \Sigma \)
    \end{enumerate}
    By \cref{set-identities:associativity-of-intersection-difference},
    we have \( A \intersect \mleft( B - A \mright) = \mleft( A - A \mright) \intersect \mleft( B - A \mright) = \emptyset \).
    In other words, \( A \) and \( B - A \) are disjoint sets.
    Furthermore, \( A \union \mleft( B - A \mright) = \mleft( A \union B \mright) - \mleft( A - A \mright) = \mleft( A \union B \mright) - \emptyset = A \union B \) by
    \cref{set-identities:union-difference}.
    In other words,
    \begin{equation*}
        \mu(A \union \mleft( B - A \mright)) = \mu(A \union B)
    \end{equation*}
    But \( \function{\mu}{\Sigma}{\realset} \) is a \aref{dfn:measure-space}{measure} and \( A \) and \( B - A \) are disjoint sets,
    so we have,
    \begin{equation*}
        \mu(A \union B) = \mu(A) + \mu(B - A)
    \end{equation*}
    Also, \( A \union B = B \) since \( A \subseteq B \). Therefore, \( \mu(A \union B) = \mu(B) = \mu(A) + \mu(B - A) \).
    Since \( \mu \) is a measure, it follows from \cref{dfn:measure-space} that \( \mu(A) \geq 0 \) and \( \mu(B - A) \geq 0 \).
    Using the familiar property of real numbers, we thus infer that \( \mu(A) + \mu(B - A) \geq 0 + \mu(A)\).
    In other words,
    \begin{equation*}
        \mu(B) = \mu(A) + \mu(B - A) \geq \mu(A)
    \end{equation*}
    Thus, \( \mu(A) \leq \mu(B) \).

    \vspace{0.3cm}
    \noindent
    By the \aref{rule-of-inference:deduction-theorem}{Deduction Theorem},
    \( A \subseteq B {\implies} \mu(A) \leq \mu(B) \).
    Since \( A \) and \( B \) are arbitrary, the theorem follows from \aref{rule-of-inference:universal-generalization}{Universal Generalization}.
\end{proof}

\noindent
The concept of measure spaces is a generalization of geometric measurements (length, area, and volume) as well other measurable quantities such
as mass, energy, probability, etc. Measure spaces are studied extensively in measure theory.

\subsection{Graphs}
Many real-world systems can be elegantly described by means of a network of interconnected points (also called nodes).
A typical example of such systems is a local area network of say, an office building.
The mathematical abstraction of such systems gives rise to the class of mathematical structures known as graphs.

\begin{dfn}[Graph]
    \phantomsection \label{dfn:graph}
    A structure \( \mathcal{G} = \tuple{V, E, \psi} \) is a \keyword{graph} if and only if the following conditions hold:
    \begin{enumerate}[topsep=2pt,label=\textup{\texttt{(\alph*)}}]
        \item \( V \) and \( E \) are sets such that \( V \intersect E = \emptyset \)
        \item \( \function{\psi}{E}{\setbuilderbig{ \mleft\{ x,y \mright\} }{x \in V \,\,{\land}\,\, y \in V}} \)
    \end{enumerate}
    The elements of \( V \) are called the \keyword{vertices} (or \keyword{nodes}) of \( \mathcal{G} \).
    The elements of \( E \) are called the \keyword{edges} of \( \mathcal{G} \).
    The function \( \psi \) is called the \keyword{incidence function} of \( \mathcal{G} \).
\end{dfn}

\begin{notation}
    If the need arises to mention two distinct graphs within the same context,
    it may help to distinguish the corresponding components for each graph using superscripts.
    For example, \( \mathcal{G} = \tuple{V^{\mathcal{G}}, E^{\mathcal{G}}, \psi^{\mathcal{G}}} \) and
    \( \mathcal{H} = \tuple{V^{\mathcal{H}}, E^{\mathcal{H}}, \psi^{\mathcal{H}}} \) represent two graphs.
\end{notation}

\begin{dfn}
    If \( \mathcal{G} = \tuple{V, E, \psi} \) is a graph, then the following terminology applies:
    \begin{enumerate}[topsep=2pt,label=\textup{\texttt{(\alph*)}}]
        \item \( e \) is said to \keyword{join} \( u \) and \( v \) if and only if,
            \begin{itemize}[topsep=2pt]
                \item \( e \in E \)
                \item \( u, v \in V \)
                \item \( \psi(e) = \mleft\{ u, v \mright\} \)
            \end{itemize}
        \item \( u \) and \( v \) are the \keyword{ends} of \( e \) if and only if \( e \) joins \( u \) and \( v \)
        \item \( u \) and \( v \) are \keyword{incident with} \( e \) if and only if \( e \) joins \( u \) and \( v \)
        \item \( e \) is \keyword{incident with} \( u \) and \( v \) if and only if \( e \) joins \( u \) and \( v \)
        \item Two vertices that are incident with a common edge are \keyword{adjacent}. \phantomsection \label{enum:graph-adjacent-vertices}
        \item Two edges that are incident with a common vertex are \keyword{adjacent}.
        \item Two or more distinct edges that have the same ends are \keyword{parallel}. \phantomsection \label{enum:graph-parallel-edges}
        \item Two distinct adjacent vertices are \keyword{neighbours}.
        \item \( e \) is a \keyword{loop} if and only if \( e \in E \) and \( \exists v\,{\in}\,V \, \mleft[ \psi(e) = \mleft\{ v \mright\} \mright] \)
            \phantomsection \label{enum:graph-loop}
        \item \( e \) is a \keyword{link} if and only if \( e \in E \) and
            \( \exists u\,{\in}\,V \, \exists v\,{\in}\,V \, \mleft[ u \neq v \,\,{\land}\,\, \psi(e) = \mleft\{ u, v \mright\} \mright] \)
            \phantomsection \label{enum:graph-link}
        \item \( \abs{V} \) is called the \keyword{order} of \( \mathcal{G} \)
        \item \( \abs{E} \) is called the \keyword{size} of \( \mathcal{G} \)
    \end{enumerate}
\end{dfn}

\begin{dfn}[Classification of Graphs]
    \phantomsection \label{dfn:classification-of-graphs}
    A graph \( \mathcal{G} = \tuple{V, E, \psi} \) is:
    \begin{enumerate}[topsep=2pt,label=\textup{\texttt{(\alph*)}}]
        \item a \keyword{null graph} if and only if \( V = \emptyset \) and \( E = \emptyset \)
        \item \keyword{trivial} if and only if \( \abs{V} = 1 \)
        \item \keyword{simple} if and only if it has no \aref{enum:graph-loop}{loops} and \aref{enum:graph-parallel-edges}{parallel edges}
        \item \keyword{complete} if and only if it is \keyword{simple} and any two vertices are \aref{enum:graph-adjacent-vertices}{adjacent}
        \item \keyword{empty} if and only if no two vertices are \aref{enum:graph-adjacent-vertices}{adjacent}; that is, \( E = \emptyset \)
    \end{enumerate}
\end{dfn}

\noindent
For normal graphs, the orientation of the edges doesn't matter; that is why in \cref{dfn:graph} the incidence function maps each edge to an \emph{unordered} pair.
However, there are a great many situations in which it is valuable to incorporate directional information into a graph's structure.
We call graphs with directional information \emph{directed graphs}.

\begin{dfn}[Digraphs]
    \phantomsection \label{dfn:directed-graph}
    A structure \( \mathcal{D} = \tuple{V, E, \psi} \) is a \keyword{directed graph} (or \keyword{digraph} for short)
    if and only if:
    \begin{enumerate}[topsep=2pt,label=\textup{\texttt{(\alph*)}}]
        \item \( V \intersect E = \emptyset \) 
        \item \( \function{\psi}{E}{\setbuilderbig{\orderedpair{x}{y}}{x \in V \,\,{\land}\,\, y \in V}} \)
    \end{enumerate}
\end{dfn}

\begin{rem}
    Note how the incidence function for \( \mathcal{D} \) maps to \aref{dfn:kuratowski-ordered-pairs}{ordered pairs} in order to capture directional information.
\end{rem}
